\newcommand{\DoNotLoadEpstopdf}{}
\documentclass[11pt,letterpaper]{article}
\usepackage[T1]{fontenc}
\usepackage{lmodern,microtype}
\usepackage[margin=0.88in,headheight=15pt,headsep=18pt,footskip=27pt]{geometry}
\usepackage{amsmath,amssymb,amsthm,mathtools,bm}
\usepackage{booktabs,longtable,array,tabularx,graphicx,float,caption}
\usepackage{enumitem,fancyhdr,xcolor,needspace,placeins,xurl,aliascnt}
\usepackage[hidelinks]{hyperref}
\usepackage[nameinlink,noabbrev]{cleveref}
\allowdisplaybreaks
\setlength{\emergencystretch}{3em}
\setlist[enumerate]{topsep=4pt,itemsep=2pt}
\setlist[itemize]{topsep=4pt,itemsep=2pt}
\setlength{\parskip}{3pt}
\setlength{\parindent}{1.3em}
\setcounter{tocdepth}{2}
\makeatletter
\renewcommand*\l@section{\@dottedtocline{1}{0em}{4.3em}}
\renewcommand*\l@subsection{\@dottedtocline{2}{1.5em}{6.0em}}
\makeatother
\hypersetup{pdftitle={Projector-Resolved Seed Calculus VII: Final Traditional Edition},pdfauthor={Komorebi P.},pdfsubject={Complete corrected finite core, strengthened quantum packet development, and mathematical proof completions}}
\pagestyle{fancy}\fancyhf{}
\fancyhead[L]{\small\textsc{Projector-Resolved Seed Calculus VII}}
\fancyhead[R]{\small Komorebi P.}
\fancyfoot[C]{\thepage}
\newtheorem{theorem}{Theorem}[section]
\newaliascnt{proposition}{theorem}\newtheorem{proposition}[proposition]{Proposition}\aliascntresetthe{proposition}
\newaliascnt{lemma}{theorem}\newtheorem{lemma}[lemma]{Lemma}\aliascntresetthe{lemma}
\newaliascnt{corollary}{theorem}\newtheorem{corollary}[corollary]{Corollary}\aliascntresetthe{corollary}
\theoremstyle{definition}
\newaliascnt{definition}{theorem}\newtheorem{definition}[definition]{Definition}\aliascntresetthe{definition}
\newaliascnt{example}{theorem}\newtheorem{example}[example]{Example}\aliascntresetthe{example}
\newaliascnt{construction}{theorem}\newtheorem{construction}[construction]{Construction}\aliascntresetthe{construction}
\newaliascnt{conjecture}{theorem}\newtheorem{conjecture}[conjecture]{Conjecture}\aliascntresetthe{conjecture}
\newaliascnt{standardfact}{theorem}\newtheorem{standardfact}[standardfact]{Standard Fact}\aliascntresetthe{standardfact}
\newtheorem*{statement}{Supplementary statement}
\theoremstyle{remark}
\newaliascnt{remark}{theorem}\newtheorem{remark}[remark]{Remark}\aliascntresetthe{remark}
\newaliascnt{boundary}{theorem}\newtheorem{boundary}[boundary]{Boundary}\aliascntresetthe{boundary}
\newaliascnt{certificate}{theorem}\newtheorem{certificate}[certificate]{Recorded Numerical Check}\aliascntresetthe{certificate}
\newaliascnt{classical}{theorem}\newtheorem{classical}[classical]{Classical Input}\aliascntresetthe{classical}
\crefname{classical}{classical input}{classical inputs}
\crefname{standardfact}{standard fact}{standard facts}
\newcommand{\cA}{\mathcal A}\newcommand{\cB}{\mathcal B}\newcommand{\cC}{\mathcal C}
\newcommand{\cE}{\mathcal E}\newcommand{\cF}{\mathcal F}\newcommand{\cG}{\mathcal G}
\newcommand{\cH}{\mathcal H}\newcommand{\cK}{\mathcal K}\newcommand{\cL}{\mathcal L}
\newcommand{\cO}{\mathcal O}\newcommand{\cD}{\mathcal D}\newcommand{\cP}{\mathcal P}
\newcommand{\cR}{\mathcal R}\newcommand{\cS}{\mathcal S}\newcommand{\cT}{\mathcal T}
\newcommand{\bbC}{\mathbb C}\newcommand{\bbR}{\mathbb R}\newcommand{\bbZ}{\mathbb Z}\newcommand{\bbN}{\mathbb N}
\newcommand{\CC}{\mathbb C}\newcommand{\RR}{\mathbb R}\newcommand{\ZZ}{\mathbb Z}\newcommand{\NN}{\mathbb N}
\newcommand{\C}{\mathbb C}\newcommand{\R}{\mathbb R}\newcommand{\N}{\mathbb N}
\newcommand{\Hcal}{\mathcal H}\newcommand{\Acal}{\mathcal A}\newcommand{\Bcal}{\mathcal B}
\newcommand{\Ccal}{\mathcal C}\newcommand{\Dcal}{\mathcal D}\newcommand{\Ecal}{\mathcal E}
\newcommand{\Fcal}{\mathcal F}\newcommand{\Gcal}{\mathcal G}\newcommand{\Lcal}{\mathcal L}
\newcommand{\Mcal}{\mathcal M}\newcommand{\Ocal}{\mathcal O}\newcommand{\Pcal}{\mathcal P}
\newcommand{\Qcal}{\mathcal Q}\newcommand{\Rcal}{\mathcal R}\newcommand{\Scal}{\mathcal S}
\newcommand{\Tcal}{\mathcal T}\newcommand{\Ucal}{\mathcal U}\newcommand{\Vcal}{\mathcal V}
\newcommand{\Id}{\operatorname{Id}}\newcommand{\id}{\mathrm{id}}
\newcommand{\Tr}{\operatorname{Tr}}\newcommand{\tr}{\operatorname{tr}}
\newcommand{\rank}{\operatorname{rank}}\newcommand{\Ran}{\operatorname{Ran}}\newcommand{\Ker}{\operatorname{Ker}}
\newcommand{\Spec}{\operatorname{Spec}}\newcommand{\diag}{\operatorname{diag}}
\newcommand{\End}{\operatorname{End}}\newcommand{\Hom}{\operatorname{Hom}}
\newcommand{\dist}{\operatorname{dist}}\newcommand{\ad}{\operatorname{ad}}\newcommand{\Ad}{\operatorname{Ad}}
\newcommand{\Ob}{\operatorname{Ob}}\newcommand{\ob}{\mathfrak o}
\newcommand{\HS}{\mathrm{HS}}\newcommand{\phys}{\mathrm{phys}}
\newcommand{\e}{\mathrm e}\newcommand{\ii}{\mathrm i}\newcommand{\dd}{\,\mathrm d}\newcommand{\one}{\bm 1}
\newcommand{\norm}[1]{\left\lVert#1\right\rVert}\newcommand{\abs}[1]{\left\lvert#1\right\rvert}
\newcommand{\ket}[1]{\left|#1\right\rangle}\newcommand{\bra}[1]{\left\langle#1\right|}
\newcommand{\ip}[2]{\left\langle#1,#2\right\rangle}
\newcommand{\comm}[2]{[#1,#2]}\newcommand{\deph}{\mathcal D_P}
\newcommand{\spanof}{\operatorname{span}}\newcommand{\srank}{\operatorname{srank}}
\newcommand{\code}[1]{{\small\nolinkurl{#1}}}\newcommand{\lean}[1]{{\small\nolinkurl{#1}}}
\newcolumntype{L}[1]{>{\raggedright\arraybackslash}p{#1}}
\newcolumntype{Y}{>{\raggedright\arraybackslash}X}
\newcommand{\startreadingpart}[1]{\clearpage\FloatBarrier\part{#1}\setcounter{section}{0}\renewcommand{\thesection}{\Roman{part}.\arabic{section}}\renewcommand{\theHsection}{\arabic{part}.\arabic{section}}}

% Operation Komorebi: shared publication profile.
\usepackage{komorebi-publication}
\renewcommand{\KomorebiTitle}{Quantum Seed and Packet Calculus}
\renewcommand{\KomorebiSubtitle}{Symmetry, circuits, geometric transport, channels, effective dynamics, and continuum towers}
\renewcommand{\KomorebiSeries}{Core series VII}
\renewcommand{\KomorebiKind}{Traditional paper}
\renewcommand{\KomorebiEdition}{Standardized publication edition / 8 October 2026}
\renewcommand{\KomorebiShortTitle}{Paper VII}
\renewcommand{\KomorebiTitleSize}{38}
\renewcommand{\KomorebiTitleLeading}{43}
\renewcommand{\KomorebiProfile}{classical}
\renewcommand{\KomorebiOrientation}{portrait}
\KomorebiSetup
\begin{document}
\KomorebiCover
\KomorebiOpeningNote{The complete corrected finite manuscript, the complete strengthened mathematical development, and the substantive proof-completion supplements are collected in one conventional reading edition. The visual guide and Lean/evidence guide are separate companions.}

\hypersetup{pageanchor=false}

\hypersetup{pageanchor=true}
\pagenumbering{roman}
\section*{Editorial reading note}
\addcontentsline{toc}{section}{Editorial reading note}
This edition preserves the original author's attribution and mathematical scope while bringing the recovered Paper VII sources into one reading order. Part I uses the corrected integrated finite manuscript and the latest finite completion. Part II retains the strengthened manuscript in full, including all 79 theorem, proposition, lemma, and corollary proofs. Part III retains every substantive mathematical section of the six formal proof supplements, followed by valid earlier derivations that were otherwise absent from the integrated paper.

The original finite item numbers are preserved inside Part I with the prefix I. The strengthened development has independent numbering with prefix II. The supplementary constructions have prefix III. Repeated results are retained because the later versions supply different scope, proof detail, or connections; the combined count is not a count of distinct new discoveries.

The selected saved finite formal checkpoint records 19 covered manuscript items, 323 theorem declarations, and 25 active Lean modules. These figures describe the archived formal development. This editorial build does not constitute a new Lean compilation, and the strengthened Part II is not added to that formal count. Compiler records, exact hypotheses, supporting declarations, numerical evidence, and reproduction instructions are presented in the separate Lean/evidence guide.

The earlier Feshbach statement is read with both eliminated-block and full-resolvent invertibility. Exact common-kernel compression is treated as automatic and is distinguished from operator closure. Normalized channels and unnormalized error families are treated with their correct recovery scales. The projected connection uses the actual range bundle and the Lie-bracket term. Part II also makes explicit the nonzero-space, common-space, self-adjointness, finite-address, and trace-class hypotheses needed in its limit statements.

The five physical continuum targets remain open in this corpus. The bounded direct and norm limits in Part II do not provide an interacting continuum field theory, unbounded gravitational constraint closure, Lorentzian reconstruction, an Einstein infrared action, or a physical massless spin-two sector.
\clearpage
\begingroup\small\tableofcontents\endgroup
\clearpage
\pagenumbering{arabic}
\startreadingpart{Corrected finite quantum seed calculus}
\graphicspath{{assets/core/}}
\section{The quantum seed realization}\label{core:section:001}

\paragraph{Conventions.}
Unless a zero-dimensional boundary case is stated explicitly, quantum state spaces are nonzero finite-dimensional complex Hilbert spaces. A complete packet family is a finite collection satisfying
\[
P_\alpha^*=P_\alpha,\qquad
P_\alpha P_\beta=\delta_{\alpha\beta}P_\alpha,\qquad
\sum_\alpha P_\alpha=I.
\]
All operator norms in the analytic estimates are induced by the Hilbert-space norm, not the entrywise maximum norm or the Frobenius norm. We write $[A,B]=AB-BA$ and use units with $\hbar=1$. Parameter derivatives are real derivatives. A matrix $X$ is supported on a code projector $P$ when $PXP=X$.

The common operator sequence of the series becomes
\[
\boxed{
\begin{gathered}
\text{quantum seed}
\longrightarrow
\text{gate or observable address}
\longrightarrow
\text{symmetry projection}\\
\longrightarrow
\text{packet transfer}
\longrightarrow
\text{constraint}
\longrightarrow
\text{refinement}
\longrightarrow
\text{reconstruction}
\end{gathered}
}
\]

\begin{figure}[htbp]
\centering
\includegraphics[width=0.98\textwidth]{01_quantum_seed_architecture.pdf}
\caption{The quantum realization of the common seed-calculus operator chain. The same finite packet language is used for states, gates, Hamiltonians, channels, geometric projectors, and effective dynamics.}
\label{core:fig:architecture}
\end{figure}

\begin{definition}[Finite quantum seed system]\label{core:definition:001}
A finite quantum seed system is a tuple
\[
\mathfrak Q=
\left(
\cH,
\psi_0,
\cO,
\cR,
\{P_{\ell,\alpha}\},
H,
\tau
\right),
\]
where $\cH$ is a finite-dimensional Hilbert space, $\psi_0\in\cH$ is a nonzero seed, $\cO$ is the declared operator grammar, $\cR$ is a collection of algebraic or physical constraints, $\{P_{\ell,\alpha}\}$ is an exact packet system at each refinement level, $H=H^*$ is a Hamiltonian when dynamics are present, and $\tau$ specifies the topology used in a later completion.
\end{definition}

The generated state space is
\[
\cH(\psi_0;\cO)
=
\operatorname{span}
\left\{
O_r\cdots O_1\psi_0:
O_j\in\cO,
\ r\ge0
\right\}.
\]
The quantum seed is cyclic for the grammar when
\[
\cH(\psi_0;\cO)=\cH.
\]
This is a statement about the pair $(\psi_0,\cO)$, not about the vector alone. A state may be cyclic for the full matrix algebra and noncyclic for a restricted physical gate set.

\section{Quantum seed rank and cyclicity}\label{core:section:002}

\subsection{Finite observable algebras}

Let a finite-dimensional unital $*$-subalgebra of operators act in its standard multiplicity form
\[
\cA
\cong
\bigoplus_{i=1}^r
\left(M_{d_i}(\bbC)\otimes I_{m_i}\right)
\]
on
\[
\cH
\cong
\bigoplus_{i=1}^r
\left(\bbC^{d_i}\otimes\bbC^{m_i}\right).
\]
Here every $d_i$ is a positive integer and every multiplicity $m_i$ is a nonnegative integer. An empty block family represents the zero space and has seed rank zero; an empty maximum below is assigned the value zero. A vector $\psi_i\in\bbC^{d_i}\otimes\bbC^{m_i}$ is identified with a $d_i\times m_i$ coefficient matrix $C_i$.

\begin{theorem}[Orbit dimension of one quantum seed]\label{core:theorem:001}\label{core:thm:seed-orbit}
For integers $d\ge1$ and $m\ge0$, consider the action of $M_d(\bbC)\otimes I_m$ on $\bbC^d\otimes\bbC^m$. Let $C$ be the coefficient matrix of a vector $\psi$. Then
\[
\boxed{
\dim\operatorname{span}\{(A\otimes I_m)\psi:A\in M_d(\bbC)\}
=d\,\rank C.
}
\]
The seed is cyclic if and only if $\rank C=m$, which is possible only when $m\le d$.
\end{theorem}

\begin{proof}\phantomsection\label{core:proof:001}
Under the coefficient-matrix identification, $(A\otimes I_m)\psi$ corresponds to $AC$. Every row of $AC$ lies in the row space of $C$. Conversely, any $d\times m$ matrix whose rows lie in the row space of $C$ equals $AC$ for a suitable $A$: choose each row of $A$ to produce the desired row combination. The orbit therefore consists of all matrices with $d$ independently chosen rows in the $\rank C$-dimensional row space of $C$, and has dimension $d\rank C$. It equals the full $dm$-dimensional space exactly when $\rank C=m$.
\phantomsection\label{core:proof:001:end}\end{proof}

\begin{theorem}[Quantum seed-rank formula]\label{core:theorem:002}\label{core:thm:seed-rank}
For
\[
\cA=\bigoplus_i(M_{d_i}(\bbC)\otimes I_{m_i})
\quad\text{on}\quad
\cH=\bigoplus_i(\bbC^{d_i}\otimes\bbC^{m_i}),
\]
the minimum number of seeds required to generate $\cH$ is
\[
\boxed{
\operatorname{srank}_{\cA}(\cH)
=
\max_i\left\lceil\frac{m_i}{d_i}\right\rceil.
}
\]
\end{theorem}

\begin{proof}\phantomsection\label{core:proof:002}
Each seed contributes a coefficient matrix of rank at most $d_i$ in the $i$th block. Thus $k$ seeds can cover at most a $kd_i$-dimensional subspace of the multiplicity factor $\bbC^{m_i}$, so $k\ge\lceil m_i/d_i\rceil$ for every $i$.

For the converse, take
\[
k=\max_i\left\lceil\frac{m_i}{d_i}\right\rceil.
\]
In each block, partition a basis of $\bbC^{m_i}$ into at most $k$ groups of size at most $d_i$. Choose the coefficient matrix of seed $a$ so that its nonzero rows span the $a$th group. The combined row spaces of the $k$ seeds then span $\bbC^{m_i}$ in every block. The preceding theorem shows that their $\cA$-orbits generate the whole direct sum.
\phantomsection\label{core:proof:002:end}\end{proof}

\begin{corollary}[Local orbit and Schmidt rank]\label{core:corollary:001}\label{core:thm:schmidt}
Let $\psi\in\cH_A\otimes\cH_B$, with $\dim\cH_A=d_A\ge1$, and let $r$ be its Schmidt rank. The reduced positive operator on $B$ is $\rho_B=\Tr_A(\ket\psi\bra\psi)$; normalization of $\psi$ is immaterial to its support. Then
\[
\boxed{
\dim\operatorname{span}\{(A\otimes I_B)\psi:A\in M_{d_A}(\bbC)\}
=d_A r.
}
\]
In particular, local operators on $A$ generate the complete bipartite Hilbert space exactly when the reduced state on $B$ has full support.
\end{corollary}

\begin{proof}\phantomsection\label{core:proof:003}
A Schmidt decomposition of length $r$ gives a coefficient matrix of rank $r$. Applying Theorem~\ref{core:thm:seed-orbit} gives the orbit dimension $d_A r$. The range of $\rho_B$ is the span of the $r$ Schmidt vectors on $B$, so full support is equivalent to $r=\dim\cH_B$. Because $d_A>0$, this is equivalent to the orbit having the full bipartite dimension.
\phantomsection\label{core:proof:003:end}\end{proof}

This is an algebraic span statement for arbitrary local operators. It is not a claim that a normalized unitary gate orbit contains every vector, nor a deterministic protocol for preparing arbitrary remote states. Entanglement controls the multiplicity support appearing in that algebraic span.

\subsection{A computational-basis seed}

\begin{theorem}[One seed for the computational basis]\label{core:theorem:003}\label{core:thm:computational}
For an integer $n\ge0$, let
\[
\psi_0=\ket{0}^{\otimes n}
\]
and let $X_j$ denote the Pauli $X$ operator on qubit $j$. Then
\[
\left\{
X_1^{a_1}\cdots X_n^{a_n}\psi_0:
(a_1,\ldots,a_n)\in\{0,1\}^n
\right\}
\]
is the complete computational basis. Hence the seed is cyclic under the algebra generated by the $X_j$.
\end{theorem}

\begin{proof}\phantomsection\label{core:proof:004}
For a binary address $a=(a_1,\ldots,a_n)$,
\[
X_1^{a_1}\cdots X_n^{a_n}\ket{0\cdots0}
=
\ket{a_1\cdots a_n}.
\]
Distinct binary addresses produce distinct orthonormal basis vectors. There are $2^n$ of them, equal to the dimension of the Hilbert space.
\phantomsection\label{core:proof:004:end}\end{proof}

\begin{certificate}[Three-qubit seed orbit]\label{core:certificate:001}
The implementation generated all eight computational basis states from $\ket{000}$ using the three local $X$ address operators. The measured orbit rank was $8/8$.
\end{certificate}

\section{Cyclic symmetry packets}\label{core:section:003}

\subsection{Character projectors}

Let $N\ge1$ be an integer, and let $S$ be a unitary symmetry whose order divides $N$:
\[
S^N=I.
\]
Set $\omega=e^{2\pi i/N}$.

\begin{theorem}[Cyclic quantum projectors]\label{core:theorem:004}\label{core:thm:cyclic}
The operators
\[
\boxed{
P_j
=
\frac1N\sum_{r=0}^{N-1}\omega^{-jr}S^r,
\qquad
0\le j<N,
}
\]
are self-adjoint, pairwise orthogonal projectors satisfying
\[
P_jP_k=\delta_{jk}P_j,
\qquad
\sum_{j=0}^{N-1}P_j=I,
\qquad
SP_j=\omega^jP_j.
\]
\end{theorem}

\begin{proof}\phantomsection\label{core:proof:005}
Using $S^{r+s}=S^{r+s\bmod N}$ and the finite character orthogonality relation,
\[
\frac1N\sum_{r=0}^{N-1}\omega^{(k-j)r}=\delta_{jk},
\]
one obtains $P_jP_k=\delta_{jk}P_j$. Summing over $j$ cancels every term indexed by $r\ne0$ and retains the $r=0$ term, proving completeness. Unitarity and the change of residue $r\mapsto-r$ give
\[
P_j^*=\frac1N\sum_{r=0}^{N-1}\omega^{jr}S^{-r}=P_j.
\]
Finally, multiplication by $S$ shifts the summation index and contributes the factor $\omega^j$.
\phantomsection\label{core:proof:005:end}\end{proof}

For parity,
\[
S=Z^{\otimes n},
\qquad
P_{\mathrm{even}}=\frac{I+S}{2},
\qquad
P_{\mathrm{odd}}=\frac{I-S}{2}.
\]

\subsection{Tensor refinement}

Fix $N\ge1$. Suppose every factor $a$ has a complete orthogonal family $P^{(a)}_0,\ldots,P^{(a)}_{N-1}$. The fine tensor packets are
\[
P_{\mathbf j}^{\mathrm{fine}}
=
P^{(1)}_{j_1}\otimes\cdots\otimes P^{(m)}_{j_m}.
\]

\begin{theorem}[Tensor/polyphase refinement]\label{core:theorem:005}\label{core:thm:tensor}
For each residue $r\in\bbZ/N\bbZ$, define
\[
\boxed{
P_r^{\mathrm{coarse}}
=
\sum_{j_1+\cdots+j_m\equiv r\pmod N}
P^{(1)}_{j_1}\otimes\cdots\otimes P^{(m)}_{j_m}.
}
\]
Then the $P_r^{\mathrm{coarse}}$ are complete orthogonal projectors, and every coarse packet is the exact sum of its fine descendants.
\end{theorem}

\begin{proof}\phantomsection\label{core:proof:006}
The fine tensor projectors are mutually orthogonal and sum to the identity. The stated coarse families partition their address set by total residue. Sums over disjoint sets of orthogonal projectors remain orthogonal projectors, and the union of the residue classes is the full fine address set.
\phantomsection\label{core:proof:006:end}\end{proof}

\begin{figure}[htbp]
\centering
\includegraphics[width=0.86\textwidth]{10_tensor_parity_refinement.pdf}
\caption{The eight computational packets of three qubits refine the two total-parity packets exactly.}
\label{core:fig:tensor-refinement}
\end{figure}

\begin{certificate}[Exact parity refinement]\label{core:certificate:002}
For three qubits, the local computational projectors were grouped by Hamming-weight parity. The resulting coarse projectors agreed with $(I\pm Z^{\otimes3})/2$ with residual $0$.
\end{certificate}

\section{Quantum circuits as packet paths}\label{core:section:004}

Let
\[
U=U_mU_{m-1}\cdots U_1
\]
be a circuit with $m\ge1$, and let $\{P_\alpha\}$ be a complete packet system. For the empty circuit $m=0$, set $U=I$; its boundary block is $P_\beta P_\alpha=\delta_{\beta\alpha}P_\alpha$.

\begin{theorem}[Exact circuit path expansion]\label{core:theorem:006}\label{core:thm:paths}
For source packet $\alpha_0$ and target packet $\alpha_m$,
\[
\boxed{
P_{\alpha_m}UP_{\alpha_0}
=
\sum_{\alpha_1,\ldots,\alpha_{m-1}}
P_{\alpha_m}U_mP_{\alpha_{m-1}}
\cdots
P_{\alpha_1}U_1P_{\alpha_0}.
}
\]
Summing over the source and target packets reconstructs the complete circuit operator.
\end{theorem}

\begin{proof}\phantomsection\label{core:proof:007}
Insert the identity resolution $I=\sum_\alpha P_\alpha$ between every adjacent pair of gates. Distributing the finite sums gives the path expansion. Summing the left and right boundary projectors reconstructs $U$.
\phantomsection\label{core:proof:007:end}\end{proof}

\begin{theorem}[Packetwise unitarity conservation]\label{core:theorem:007}\label{core:thm:unitarity}
If $U$ is unitary and $U_{\beta\alpha}=P_\beta UP_\alpha$, then
\[
\boxed{
\sum_\beta U_{\beta\alpha}^*U_{\beta\gamma}
=
\delta_{\alpha\gamma}P_\alpha.
}
\]
Similarly,
\[
\sum_\alpha U_{\beta\alpha}U_{\delta\alpha}^*
=
\delta_{\beta\delta}P_\beta.
\]
\end{theorem}

\begin{proof}\phantomsection\label{core:proof:008}
Using completeness and unitarity,
\[
\sum_\beta U_{\beta\alpha}^*U_{\beta\gamma}
=
P_\alpha U^*\left(\sum_\beta P_\beta\right)UP_\gamma
=
P_\alpha P_\gamma.
\]
The second identity is analogous.
\phantomsection\label{core:proof:008:end}\end{proof}

\begin{example}[Bell seed circuit]\label{core:example:001}
In this example only, $H$ denotes the Hadamard gate, not a Hamiltonian. Starting from $\ket{00}$,
\[
(H\otimes I)\ket{00}
=
\frac{\ket{00}+\ket{10}}{\sqrt2},
\]
and then
\[
\operatorname{CNOT}(H\otimes I)\ket{00}
=
\frac{\ket{00}+\ket{11}}{\sqrt2}.
\]
The parity packets show that the final state lies entirely in the even sector, while the intermediate circuit path passes through both even and odd addresses.
\end{example}

\begin{figure}[htbp]
\centering
\begin{minipage}[t]{0.48\textwidth}
\centering
\includegraphics[width=\textwidth]{02_bell_seed_probabilities.pdf}
\captionof{figure}{The Bell state generated from the seed $\ket{00}$.}
\label{core:fig:bell}
\end{minipage}\hfill
\begin{minipage}[t]{0.48\textwidth}
\centering
\includegraphics[width=\textwidth]{04_circuit_packet_paths.pdf}
\captionof{figure}{Circuit reconstruction as a sum over packet paths.}
\label{core:fig:circuit-paths}
\end{minipage}
\end{figure}

\begin{certificate}[Bell and path reconstruction]\label{core:certificate:003}
The Bell-state residual and the complete parity-path reconstruction residual were both $0$ in the numerical implementation.
\end{certificate}

\section{Symmetry-resolved Hamiltonians and leakage}\label{core:section:005}

Let $H=H^*$ and let $P$ be a packet projector. If $[H,P]=0$, then $P\cH$ is invariant under $e^{-itH}$.

\begin{theorem}[Exact symmetry packets]\label{core:theorem:008}\label{core:thm:symmetry}
If $[H,P_\alpha]=0$ for every packet, then
\[
H=\sum_\alpha P_\alpha HP_\alpha
\]
and
\[
e^{-itH}
=
\sum_\alpha P_\alpha e^{-itH}P_\alpha.
\]
No amplitude transfers between distinct packets.
\end{theorem}

\begin{proof}\phantomsection\label{core:proof:009}
The commutation relation gives $P_\beta HP_\alpha=0$ for $\beta\ne\alpha$. Every polynomial in $H$ has the same block form, and the finite-dimensional exponential is the norm limit of its polynomial series.
\phantomsection\label{core:proof:009:end}\end{proof}

For two qubits and real parameters $J,h$, consider
\[
H
=
JX\otimes X
+h(Z\otimes I+I\otimes Z).
\]
In the ordered even basis $\{\ket{00},\ket{11}\}$ and odd basis $\{\ket{01},\ket{10}\}$,
\[
H_{\mathrm{even}}
=
\begin{pmatrix}
2h&J\\
J&-2h
\end{pmatrix},
\qquad
H_{\mathrm{odd}}
=
\begin{pmatrix}
0&J\\
J&0
\end{pmatrix}.
\]
Therefore
\[
\Spec H_{\mathrm{even}}
=\{\pm\sqrt{4h^2+J^2}\},
\qquad
\Spec H_{\mathrm{odd}}=\{\pm J\}.
\]

\begin{figure}[htbp]
\centering
\includegraphics[width=0.66\textwidth]{03_parity_block_hamiltonian.pdf}
\caption{The parity-resolved Hamiltonian is exactly block diagonal after the even and odd basis vectors are grouped.}
\label{core:fig:parity-hamiltonian}
\end{figure}

\subsection{Leakage under broken symmetry}

Set $Q=I-P$.

\begin{theorem}[Commutator leakage bound]\label{core:theorem:009}\label{core:thm:leakage}
For every self-adjoint $H$, orthogonal projector $P$, and real time $t$,
\[
\boxed{
\norm{Qe^{-itH}P}
\le
|t|\,\norm{[H,P]}.
}
\]
In particular, exact commutation gives zero leakage for all time.
\end{theorem}

\begin{proof}\phantomsection\label{core:proof:010}
Set $F(s)=e^{isH}Pe^{-isH}$. Differentiation of the actual exponential gives
\[
F'(s)=i\,e^{isH}[H,P]e^{-isH},\qquad F(0)=P.
\]
Unitary conjugation preserves the operator norm, so $\|F'(s)\|=\|[H,P]\|$. The mean-value estimate on the interval between $0$ and $t$ yields
\[
\|F(t)-P\|\le |t|\,\|[H,P]\|.
\]
Finally, $e^{itH}Qe^{-itH}P=(P-F(t))P$ and $\|P\|\le1$. Left multiplication by $e^{itH}$ preserves the norm and gives the desired inequality. This argument includes positive, negative and zero time.
\phantomsection\label{core:proof:010:end}\end{proof}

\begin{figure}[htbp]
\centering
\includegraphics[width=0.78\textwidth]{05_leakage_bound.pdf}
\caption{A parity-breaking perturbation creates nonzero transfer between the packets, while the measured leakage remains below the commutator bound.}
\label{core:fig:leakage}
\end{figure}

\begin{certificate}[Symmetry and negative control]\label{core:certificate:004}
For the parity-preserving Hamiltonian, the off-packet block norm and maximum measured leakage were both $0$. After adding a parity-breaking $X\otimes I$ perturbation, the maximum leakage was approximately $0.79196$, and the maximum violation of the proved bound was $0$.
\end{certificate}

\section{Moving packets, connections, and curvature}\label{core:section:006}

\subsection{Moving quantum coordinates}

\begin{theorem}[Moving-frame Hamiltonian]\label{core:theorem:010}\label{core:thm:moving}
Let $U(t)$ and $\psi(t)$ be differentiable on an open real time interval, let $U(t)$ be unitary, and let $H(t)=H(t)^*$. Suppose
\[
i\dot\psi(t)=H(t)\psi(t).
\]
Writing
\[
\psi(t)=U(t)y(t)
\]
gives
\[
\boxed{
i\dot y(t)
=
\left(U(t)^*H(t)U(t)-iU(t)^*\dot U(t)\right)y(t).
}
\]
Both the connection term $-iU^*\dot U$ and the effective Hamiltonian $U^*HU-iU^*\dot U$ are Hermitian.
\end{theorem}

\begin{proof}\phantomsection\label{core:proof:011}
Differentiate $\psi=Uy$:
\[
i\dot U y+iU\dot y=HUy.
\]
Multiplying by $U^*$ gives the formula. Differentiating $U^*U=I$ shows $U^*\dot U$ is skew-Hermitian, so $-iU^*\dot U$ is Hermitian.
\phantomsection\label{core:proof:011:end}\end{proof}

\begin{figure}[htbp]
\centering
\includegraphics[width=0.70\textwidth]{08_moving_frame_connection.pdf}
\caption{The moving-frame equation closes to numerical precision only when the connection term is retained.}
\label{core:fig:moving-frame}
\end{figure}

\begin{certificate}[Moving-frame correction]\label{core:certificate:005}
The complete moving-frame equation had residual $1.0645\times10^{-10}$, while omitting the connection produced residual $0.7000$.
\end{certificate}

\subsection{A canonical projector connection}

Let $P_1(t),\ldots,P_m(t)$ be differentiable on an open real time interval, with $P_a(t)^*=P_a(t)$ and
\[
P_aP_b=\delta_{ab}P_a,
\qquad
\sum_aP_a=I.
\]

\begin{theorem}[Canonical packet connection]\label{core:theorem:011}\label{core:thm:canonical}
Define
\[
\boxed{
C(t)
=
\frac12\sum_a[\dot P_a(t),P_a(t)]
=
\sum_a\dot P_a(t)P_a(t).
}
\]
Then
\[
C(t)^*=-C(t)
\]
and
\[
\boxed{
[C(t),P_b(t)]=\dot P_b(t)
}
\]
for every packet $b$.
\end{theorem}

\begin{proof}\phantomsection\label{core:proof:012}
Differentiating $P_a^2=P_a$ gives
\[
\dot P_aP_a+P_a\dot P_a=\dot P_a.
\]
Summing over $a$ and using $\sum_a\dot P_a=0$ yields
\[
\sum_aP_a\dot P_a=-\sum_a\dot P_aP_a,
\]
which proves both displayed formulas for $C$ and its skew-Hermiticity.

For fixed $b$,
\[
CP_b=\dot P_bP_b.
\]
When $a\ne b$, differentiating $P_bP_a=0$ gives
\[
P_b\dot P_a=-\dot P_bP_a.
\]
Hence
\[
P_bC
=\sum_aP_b\dot P_aP_a
=-\dot P_b(I-P_b),
\]
and subtraction gives $[C,P_b]=\dot P_b$.
\phantomsection\label{core:proof:012:end}\end{proof}

The connection is the infinitesimal transport operator of the moving packet frame. It is determined by the packet family itself, without choosing a basis inside any packet.

\subsection{Projector curvature}

\begin{classical}[Projected connection]\label{core:classical:001}\label{core:thm:projected}
Let $M$ be a smooth finite-dimensional real manifold without boundary, let $V$ be a finite-dimensional real normed vector space, and let $P:M\to\operatorname{End}_{\bbR}(V)$ be smooth with $P(x)^2=P(x)$. In the quantum application, $V$ is the underlying real space of a finite complex Hilbert space and $P$ is an orthogonal complex-linear projector.

The rank of $P$ is locally constant. On each open rank region, its ranges form a smooth vector bundle. For a smooth range section $s$, so that $Ps=s$, define
\[
\nabla_Xs=P\,ds(X).
\]
This is a connection on smooth range sections. For smooth vector fields $X,Y$, its curvature satisfies
\[
\boxed{
(\nabla_X\nabla_Y-\nabla_Y\nabla_X-\nabla_{[X,Y]})s
=P[dP(X),dP(Y)]P\,s.
}
\]
Equivalently, its endomorphism-valued curvature is
\[
\boxed{F=P(dP\wedge dP)P.}
\]
No global frame, globally constant rank or commuting-vector-field assumption is imposed. The geometric formalization concerns real smooth structures, not complex-analytic manifolds or boundary extensions.
\end{classical}

\begin{proof}\phantomsection\label{core:proof:013}
\textit{Range coordinates.} Fix $a\in M$ and put
\[
T_a(x)=P(x)P(a)+(I-P(x))(I-P(a)).
\]
Then $T_a(a)=I$ and $P(x)T_a(x)=T_a(x)P(a)$. Invertibility is open, so on a neighborhood of $a$, $T_a(x)$ restricts to an isomorphism from $\Ran P(a)$ to $\Ran P(x)$. This proves local constancy of rank. The local coordinates and their transitions are
\[
(x,v)\longmapsto (x,T_a(x)^{-1}v),\qquad
u\longmapsto T_b(x)^{-1}T_a(x)u.
\]
Their smoothness follows from smoothness of $P$ and inversion. The transition maps satisfy the cocycle identity by cancellation. These charts construct the smooth range bundles on the open rank regions, including when different connected components have different ranks.

\textit{Connection rules.} Smoothness of $P\,ds(X)$ and $P^2=P$ show that $\nabla_Xs$ is a smooth range section. Differentiation and projection give
\[
\begin{aligned}
\nabla_{X+Y}s&=\nabla_Xs+\nabla_Ys,&
\nabla_{fX}s&=f\nabla_Xs,\\
\nabla_X(s+t)&=\nabla_Xs+\nabla_Xt,&
\nabla_X(fs)&=f\nabla_Xs+(Xf)s
\end{aligned}
\]
for smooth real $f$. The range coordinates above identify ambient smooth range sections with smooth sections of the constructed bundle.

\textit{Curvature.} Differentiating $P^2=P$ gives $P(dP(X))P=0$. Write $D_Xs=ds(X)$. For ambient smooth sections,
\[
D_XD_Ys-D_YD_Xs-D_{[X,Y]}s=0.
\]
Expansion of the projected derivatives therefore leaves
\[
R(X,Y)s=P(dP(X))D_Ys-P(dP(Y))D_Xs.
\]
Since $Ps=s$, differentiation gives $D_Ys=(dP(Y))s+PD_Ys$. The terms with a middle $P$ vanish. Thus
\[
R(X,Y)s=P\bigl(dP(X)dP(Y)-dP(Y)dP(X)\bigr)P\,s.
\]
This expression is alternating, bilinear in the tangent vectors and range-preserving. The ordinary chain rule and the Lie-bracket transformation law identify the local calculations on chart overlaps, yielding the global formula.
\phantomsection\label{core:proof:013:end}\end{proof}

For an orthogonal rank-one complex projector, use the scalar connection convention $A=-i\langle u,du\rangle$ in a local unit frame. The associated scalar curvature density is
\[
\cF_{\mu\nu}
=-i\Tr\left(P[\partial_\mu P,\partial_\nu P]\right).
\]
On the Bloch sphere,
\[
P(\theta,\phi)
=
\frac12\left(I+\mathbf n(\theta,\phi)\cdot\boldsymbol\sigma\right),
\]
where
\[
\mathbf n=(\sin\theta\cos\phi,\sin\theta\sin\phi,\cos\theta).
\]
The Pauli multiplication law gives
\[
[\partial_\theta P,\partial_\phi P]
=\frac{i}{2}(\partial_\theta\mathbf n\mathbin{\times}\partial_\phi\mathbf n)\cdot\boldsymbol\sigma.
\]
Using $\Tr(P\,\mathbf a\cdot\boldsymbol\sigma)=\mathbf n\cdot\mathbf a$ therefore yields
\[
\boxed{\cF_{\theta\phi}
=\frac12\mathbf n\cdot(\partial_\theta\mathbf n\mathbin{\times}\partial_\phi\mathbf n)
=\frac12\sin\theta.}
\]
The scalar sign changes if the opposite convention for $A$ is chosen. This explicit worked calculation is not a separate entry in the 19-item Lean register.

\begin{figure}[htbp]
\centering
\includegraphics[width=0.76\textwidth]{09_bloch_curvature.pdf}
\caption{Finite-difference evaluation of the Bloch projector curvature against the exact density $\frac12\sin\theta$.}
\label{core:fig:bloch}
\end{figure}

\begin{certificate}[Geometric packet identities]\label{core:certificate:006}
The canonical connection reconstructed every packet derivative with residual $0$. The maximum Bloch-curvature residual was $2.0803\times10^{-11}$.
\end{certificate}

\section{Measurements and quantum channels}\label{core:section:007}

\subsection{Packet measurements}

Let $\rho$ be a density matrix and let $\{E_k\}$ be a POVM:
\[
E_k\ge0,
\qquad
\sum_kE_k=I.
\]
The outcome probabilities are
\[
p_k=\Tr(E_k\rho).
\]

\begin{theorem}[Block-diagonal measurement invariance]\label{core:theorem:012}\label{core:thm:measurement}
Let $\{P_\alpha\}$ be a complete packet family, and suppose every measurement effect is packet diagonal:
\[
E_k=\sum_\alpha P_\alpha E_kP_\alpha.
\]
Define the packet-dephased state
\[
\cD_P(\rho)=\sum_\alpha P_\alpha\rho P_\alpha.
\]
Then
\[
\boxed{
\Tr(E_k\rho)
=
\Tr(E_k\cD_P(\rho))
}
\]
for every outcome $k$.
\end{theorem}

\begin{proof}\phantomsection\label{core:proof:014}
Insert the packet decomposition of $E_k$ and use cyclicity of the trace:
\[
\Tr(E_k\rho)
=
\sum_\alpha\Tr(P_\alpha E_kP_\alpha\rho)
=
\sum_\alpha\Tr(E_kP_\alpha\rho P_\alpha).
\]
\phantomsection\label{core:proof:014:end}\end{proof}

This theorem says exactly what coherence a declared measurement can and cannot see. Coherence between different packets is invisible to every packet-diagonal observable, while coherence inside a packet remains available.

\subsection{Kraus transfer}

\begin{classical}[Finite Kraus channel]\label{core:classical:002}\label{core:thm:kraus}
Let $K_1,\ldots,K_r$ satisfy
\[
\sum_aK_a^*K_a=I.
\]
Then
\[
\boxed{
\Phi(\rho)=\sum_aK_a\rho K_a^*
}
\]
is completely positive and trace preserving.
\end{classical}

\begin{proof}\phantomsection\label{core:proof:015}
Trace preservation follows from
\[
\Tr\Phi(\rho)
=\Tr\left(\rho\sum_aK_a^*K_a\right)
=\Tr\rho.
\]
For any auxiliary dimension and positive matrix $X$, each term
\[
(I\otimes K_a)X(I\otimes K_a)^*
\]
is positive, so their sum is positive. This is complete positivity.
\phantomsection\label{core:proof:015:end}\end{proof}

Packet superoperator blocks are defined by
\[
\Phi_{\alpha\beta;\gamma\delta}(X)
=
P_\alpha
\Phi(P_\gamma XP_\delta)
P_\beta.
\]
Summing all four packet indices reconstructs $\Phi(X)$ exactly.

For amplitude damping with a real parameter $0\le\gamma\le1$,
\[
K_0=
\begin{pmatrix}
1&0\\0&\sqrt{1-\gamma}
\end{pmatrix},
\qquad
K_1=
\begin{pmatrix}
0&\sqrt\gamma\\0&0
\end{pmatrix}.
\]
The excited state becomes
\[
\Phi(\ket1\bra1)
=
\gamma\ket0\bra0+(1-\gamma)\ket1\bra1.
\]

\begin{figure}[htbp]
\centering
\includegraphics[width=0.56\textwidth]{16_amplitude_damping_channel.pdf}
\caption{Amplitude damping transfers population from the excited packet to the ground packet while preserving trace and positivity.}
\label{core:fig:damping}
\end{figure}

\begin{certificate}[Channel reconstruction]\label{core:certificate:007}
For $\gamma=0.23$, the Kraus completeness and output-trace residuals were $0$. The computed output from $\ket1\bra1$ was $\diag(0.23,0.77)$. Packet dephasing changed no computational-basis measurement probability in the Bell-state test.
\end{certificate}

\section{Quantum error correction as obstruction-packet annihilation}\label{core:section:008}

Let $P=P^*=P^2$ project onto a code space $\cC\subset\cH$, and let $E_1,\ldots,E_r$ be a finite family of complex error matrices. Neither a nonzero code nor channel normalization is assumed unless stated. The operators
\[
P E_a^*E_bP
\]
are the code-level obstruction packets: if they retain logical-state information beyond one scalar multiple of $P$, the environment can distinguish logical states and exact correction is impossible.

For precision, a complete recovery is a Kraus map
\[
\cR(X)=\sum_jK_jXK_j^*,\qquad \sum_jK_j^*K_j=I.
\]
The family is \emph{exactly correctable} when one such recovery works for every coherent error combination $A_w=\sum_aw_aE_a$: for every $w$ there is a real $t_w\ge0$, independent of the logical input, such that
\[
PXP=X\quad\Longrightarrow\quad
\cR(A_wXA_w^*)=t_wX.
\]
The coefficient is an unnormalized weight. It is not necessarily a probability for an arbitrary matrix $A_w$.

\begin{classical}[Finite Knill--Laflamme criterion]\label{core:classical:003}\label{core:thm:kl}
The finite error family is exactly correctable in the sense just defined if and only if there is a positive semidefinite complex matrix $C=(c_{ab})$ such that
\[
\boxed{PE_a^*E_bP=c_{ab}P\qquad\text{for every }a,b.}
\]
If $\sum_aE_a^*E_a=I$, the equivalent normalized-channel formulation is the existence of a complete recovery satisfying
\[
\cR\!\left(\sum_aE_aXE_a^*\right)=X
\quad\text{for every }X\text{ with }PXP=X.
\]
The zero-code and empty-error-family cases are included. An empty family is not a trace-preserving channel on a nonzero ambient space.
\end{classical}

\begin{proof}\phantomsection\label{core:proof:016}
\textit{Necessity.} Fix an error $E_a$ and let $L_j=K_jE_aP$. Perfect recovery on rank-one code operators gives
\[
\sum_j\ket{L_j\psi}\bra{L_j\psi}
=t_a\ket\psi\bra\psi\qquad(\psi\in\cC).
\]
For every vector orthogonal to $\psi$, the quadratic form on the left is a sum of nonnegative squared absolute values and equals zero. Each $L_j\psi$ therefore lies in the span of $\psi$. A linear operator that preserves every line in $\cC$ is scalar there: apply it to two independent vectors and their sum. The one-dimensional case is immediate, and the zero-code case is vacuous. Consequently
\[
K_jE_aP=d_{ja}P.
\]
Recovery completeness now gives
\[
\begin{aligned}
PE_a^*E_bP
&=\sum_j (K_jE_aP)^*(K_jE_bP)\\
&=\left(\sum_j\overline{d_{ja}}d_{jb}\right)P.
\end{aligned}
\]
The coefficients form a Gram matrix, hence are positive semidefinite. For a normalized channel, apply the same rank-one argument to all composite branches $K_jE_aP$ at once to obtain the necessary condition.

\textit{Sufficiency.} Diagonalize $C=U\operatorname{diag}(\lambda_r)U^*$ with $U$ unitary and $\lambda_r\ge0$. Define compressed transformed errors
\[
F_r=\sum_aU_{ar}E_aP.
\]
Then $F_rP=F_r$ and $F_r^*F_s=\lambda_r\delta_{rs}P$. If $\lambda_r=0$, $F_r^*F_r=0$ forces $F_r=0$. For positive eigenvalues only, set $V_r=F_r/\sqrt{\lambda_r}$. These satisfy
\[
V_rP=V_r,\qquad V_r^*V_s=\delta_{rs}P.
\]
Let $S=\sum_{\lambda_r>0}V_rV_r^*$. Orthogonality makes $S$ an orthogonal projector. Use the recovery branches
\[
K_r=V_r^*,\qquad K_{\mathrm{unused}}=I-S.
\]
They are complete on the whole ambient space because
\[
\sum_rK_r^*K_r+K_{\mathrm{unused}}^*K_{\mathrm{unused}}
=S+(I-S)^2=I.
\]
The inverse unitary change of coordinates reconstructs each $E_aP$ from the $F_r$. Every recovery branch of a coherent combination is therefore scalar on the code, while its unused-space branch vanishes. For all supported matrices, the resulting recovery is
\[
\cR(A_wXA_w^*)=(w^*Cw)X.
\]
This establishes the arbitrary-family statement without assuming normalization. Under $\sum_aE_a^*E_a=I$ and $P\ne0$, summing the diagonal KL equations gives $\Tr C=1$, so the normalized-channel output is $X$. When $P=0$, every supported $X$ is zero and the conclusion is automatic. The recovery is trace preserving by completeness and completely positive at every finite ancillary dimension by the Kraus calculation in Classical Input~\ref{core:thm:kraus}.
\phantomsection\label{core:proof:016:end}\end{proof}

For the three-qubit bit-flip code,
\[
\ket{0_L}=\ket{000},
\qquad
\ket{1_L}=\ket{111},
\]
and the error family is
\[
I,
\quad X\otimes I\otimes I,
\quad I\otimes X\otimes I,
\quad I\otimes I\otimes X.
\]
The four error images of the code are pairwise orthogonal two-dimensional syndrome packets whose projectors sum to the identity on the eight-dimensional Hilbert space.

\begin{figure}[htbp]
\centering
\includegraphics[width=0.88\textwidth]{19_error_correction_syndromes.pdf}
\caption{The four orthogonal syndrome packets of the three-qubit bit-flip code. Each packet identifies one correctable error without resolving the encoded logical state.}
\label{core:fig:error-correction}
\end{figure}

\begin{certificate}[Exact bit-flip correction]\label{core:certificate:008}
The Knill--Laflamme residual and syndrome-completeness residual were both $0$. A mixed single-bit-flip channel applied to a nontrivial logical superposition was reversed with density-matrix residual $1.0385\times10^{-16}$.
\end{certificate}

\section{Exact effective dynamics}\label{core:section:009}

\subsection{Schur--Feshbach reduction}

Let $P$ be a retained packet projector and $Q=I-P$. In packet coordinates,
\[
H=
\begin{pmatrix}
H_{PP}&H_{PQ}\\
H_{QP}&H_{QQ}
\end{pmatrix}.
\]

\begin{theorem}[Exact Feshbach resolvent]\label{core:theorem:013}\label{core:thm:feshbach}
Assume that both $zI_Q-H_{QQ}$ and $zI-H$ are invertible. Define
\[
\boxed{
H_{\mathrm{eff}}(z)
=
H_{PP}
+H_{PQ}(z-H_{QQ})^{-1}H_{QP}.
}
\]
Then
\[
\boxed{
P(z-H)^{-1}P
=
\left(z-H_{\mathrm{eff}}(z)\right)^{-1}
}
\]
where the left side is restricted to $P\cH$, the right side is an inverse on that retained space, and both are written in the same orthonormal basis. Here $z-H_{QQ}$ and $z-H_{\mathrm{eff}}$ abbreviate the appropriately sized scalar matrices minus the indicated blocks.
\end{theorem}

\begin{proof}\phantomsection\label{core:proof:017}
The block matrix $z-H$ has lower-right block $z-H_{QQ}$. Gaussian block elimination produces the Schur complement
\[
z-H_{PP}-H_{PQ}(z-H_{QQ})^{-1}H_{QP}.
\]
The assumed invertibility of the lower-right block and of the full matrix implies invertibility of this Schur complement. The upper-left block of the full inverse is its inverse, which is exactly the displayed operator.
\phantomsection\label{core:proof:017:end}\end{proof}

The second premise cannot be dropped: at $z=0$ and $H=\operatorname{diag}(0,-1)$, the eliminated block is $1$ but $zI-H$ is singular. This counterexample is also included in the Lean module.

The self-energy
\[
\Sigma_P(z)
=H_{PQ}(z-H_{QQ})^{-1}H_{QP}
\]
contains the complete influence of the eliminated packet. Bare truncation drops $\Sigma_P(z)$ and is therefore exact only when the cross-packet coupling vanishes or the self-energy cancels for another stated reason.

\begin{figure}[htbp]
\centering
\begin{minipage}[t]{0.48\textwidth}
\centering
\includegraphics[width=\textwidth]{06_feshbach_vs_truncation.pdf}
\captionof{figure}{Exact Feshbach reduction versus bare truncation.}
\label{core:fig:feshbach}
\end{minipage}\hfill
\begin{minipage}[t]{0.48\textwidth}
\centering
\includegraphics[width=\textwidth]{07_packet_green_blocks.pdf}
\captionof{figure}{Norms of packet Green-function blocks.}
\label{core:fig:green-blocks}
\end{minipage}
\end{figure}

\begin{certificate}[Effective resolvent]\label{core:certificate:009}
For a random six-dimensional Hermitian test, the exact Feshbach residual was $3.2291\times10^{-17}$. The corresponding bare-truncation residual was $1.3091\times10^{-2}$.
\end{certificate}

\subsection{The hidden packet as memory}

For a time-independent finite Hamiltonian $H$ and fixed projectors $P,Q$, write $\psi=x+y$, with $x=P\psi$ and $y=Q\psi$. Then
\[
\begin{aligned}
i\dot x&=H_{PP}x+H_{PQ}y,\\
i\dot y&=H_{QP}x+H_{QQ}y.
\end{aligned}
\]
Solving the second equation gives
\[
y(t)
=e^{-itH_{QQ}}y(0)
-i\int_0^t e^{-i(t-s)H_{QQ}}H_{QP}x(s)\,ds.
\]
Substitution yields the exact Volterra equation
\[
\boxed{
\begin{aligned}
i\dot x(t)
={}&H_{PP}x(t)
+H_{PQ}e^{-itH_{QQ}}y(0)\\
&-i\int_0^t
K(t-s)x(s)\,ds,
\end{aligned}
}
\]
with memory kernel
\[
\boxed{
K(t)=H_{PQ}e^{-itH_{QQ}}H_{QP}.
}
\]
Thus the eliminated degrees of freedom are retained as a frequency-dependent self-energy or an exact memory term. For forward evolution the integral is causal. A time-dependent Hamiltonian requires its two-time propagator instead of the stationary exponential used here. This Volterra derivation is not a separate entry in the finite Lean register.

\section{Finite quantum geometry}\label{core:section:010}

\subsection{A projector-resolved geometry}

At resolution $\ell$, define
\[
\boxed{
\mathfrak Q_\ell
=
\left(
\cA_\ell,
\cH_\ell,
D_\ell,
\{P_{\ell,\alpha}\},
E_\ell,
C_\ell,
\Pi_\ell^{\mathrm{phys}},
T_\ell
\right).
}
\]
The entries are an observable algebra, Hilbert space, geometric or Dirac-type operator, packet projectors, fine embedding, coarse reconstruction, physical constraint projector, and transfer operator. The maps have types $E_\ell:\cH_\ell\to\cH_{\ell+1}$ and $C_\ell:\cH_{\ell+1}\to\cH_\ell$. A left inverse $C_\ell E_\ell=I$ does not by itself assert $C_\ell=E_\ell^*$ or an isometric embedding.

A coherent refinement should satisfy
\[
C_\ell E_\ell=I
\]
and
\[
\boxed{
P_{\ell,\alpha}
=
C_\ell
\left(
\sum_{\beta\succ\alpha}P_{\ell+1,\beta}
\right)
E_\ell.
}
\]
For an operator $D$, define the refinement defect
\[
\Delta_\ell(D)
=C_\ell D_{\ell+1}E_\ell-D_\ell.
\]
Exact resolution independence is $\Delta_\ell(D)=0$. When exact equality is too strong, the correct continuum problem is to prove convergence of the resolvents and uniform control of the effective self-energies.

\begin{figure}[htbp]
\centering
\includegraphics[width=0.84\textwidth]{17_quantum_spacetime_tower.pdf}
\caption{A conjectural nested packet tower for quantum geometry. The diagram is an organizational model; it is not a numerical reconstruction of physical spacetime.}
\label{core:fig:spacetime-tower}
\end{figure}

\subsection{Euclidean transfer}

\begin{theorem}[Finite Euclidean transfer]\label{core:theorem:014}\label{core:thm:euclidean}
Let $H$ be Hermitian on a nonzero finite-dimensional Hilbert space. For all real $\tau,\sigma\ge0$,
\[
T(\tau)=e^{-\tau(H-E_0I)},
\qquad
E_0=\min\Spec H,
\]
is positive definite and satisfies
\[
\boxed{
T(\tau+\sigma)=T(\tau)T(\sigma).
}
\]
Its spectrum is contained in $(0,1]$.
\end{theorem}

\begin{proof}\phantomsection\label{core:proof:018}
Diagonalize $H=\sum_jE_jP_j$. Then
\[
T(\tau)=\sum_j e^{-\tau(E_j-E_0)}P_j.
\]
Every eigenvalue is strictly positive and at most one. Multiplication adds the exponents, proving the semigroup law and $T(0)=I$. Positivity here is operator positivity, not an assertion that every matrix entry is nonnegative.
\phantomsection\label{core:proof:018:end}\end{proof}

The paper tests a controlled volume-state model
\[
H_{nn}=2\kappa+V(n),
\qquad
H_{n,n+1}=H_{n+1,n}=-\kappa,
\]
with $\kappa>0$ and a real confining potential. This is a finite transfer-matrix test only. Operator positivity alone is not reflection positivity of a Euclidean field measure. It is not derived from a gravitational path integral and is not evidence for a unique microscopic spacetime dynamics.

\begin{figure}[htbp]
\centering
\begin{minipage}[t]{0.48\textwidth}
\centering
\includegraphics[width=\textwidth]{13_euclidean_transfer_spectrum.pdf}
\captionof{figure}{Positive spectrum of the Euclidean transfer operator.}
\label{core:fig:euclidean-spectrum}
\end{minipage}\hfill
\begin{minipage}[t]{0.48\textwidth}
\centering
\includegraphics[width=\textwidth]{14_causal_volume_ground_state.pdf}
\captionof{figure}{Ground-state volume packet in the controlled toy model.}
\label{core:fig:volume-ground}
\end{minipage}
\end{figure}

\begin{figure}[htbp]
\centering
\includegraphics[width=0.68\textwidth]{15_causal_volume_transfer.pdf}
\caption{The finite Euclidean volume-transfer kernel. Positivity and semigroup composition are exact finite statements; its gravitational interpretation remains provisional.}
\label{core:fig:volume-transfer}
\end{figure}

\begin{certificate}[Finite transfer model]\label{core:certificate:010}
The minimum transfer eigenvalue was $0.0767054$, the semigroup residual was $0$, and reconstruction from four volume-residue packets had residual $0$.
\end{certificate}

\section{Spectral dimension and continuum diagnostics}\label{core:section:011}

Let $L\ge0$ be a finite Laplacian on $N\ge1$ states, with eigenvalues $\lambda_j$ and at least one zero eigenvalue. For diffusion time $\tau>0$, define the normalized heat trace
\[
K(\tau)=\frac1N\sum_j e^{-\tau\lambda_j}
\]
and the scale-dependent spectral dimension
\[
\boxed{
d_s(\tau)
=-2\frac{d\log K(\tau)}{d\log\tau}.
}
\]
For a finite system, $d_s(\tau)$ is a diagnostic, not a topological invariant. Indeed, finite differentiation gives
\[
d_s(\tau)=2\tau\,
\frac{\sum_j\lambda_j e^{-\tau\lambda_j}}{\sum_j e^{-\tau\lambda_j}}.
\]
It tends to zero as $\tau\downarrow0$ and, under the stated zero-mode assumption, as $\tau\to\infty$. If a general positive operator has no zero mode, its large-time behavior is instead asymptotic to $2\tau\lambda_{\min}$; the zero-mode conclusion must not be applied to it. A geometric dimension can only be inferred from a stable intermediate window and a controlled refinement limit.

\subsection{The circle dispersion bound}

For a periodic grid with spacing $h>0$ and physical wave number $k\in\bbR$,
\[
\lambda_{k,h}
=
\frac4{h^2}\sin^2\left(\frac{kh}{2}\right).
\]

\begin{theorem}[Low-mode second-order bound]\label{core:theorem:015}\label{core:thm:lowmode}
For every real $k$ and every nonzero real $h$, define $\lambda_{k,h}=4h^{-2}\sin^2(kh/2)$. Then
\[
\boxed{
0\le k^2-\lambda_{k,h}
\le
\frac{k^4h^2}{12}.
}
\]
\end{theorem}

\begin{proof}\phantomsection\label{core:proof:019}
Put $x=kh/2$. The all-real inequalities $|\sin x|\le|x|$ and $|x-\sin x|\le |x|^3/6$ give
\[
0\le x^2-\sin^2x
=|x-\sin x|\,|x+\sin x|
\le\frac{|x|^3}{6}\,2|x|
=\frac{x^4}{3}.
\]
Multiplying by $4/h^2$ proves both bounds. The estimate holds for every frequency, but its second-order convergence interpretation keeps $k$ fixed while $h\to0$; it is not a uniform small relative error across all grid modes.
\phantomsection\label{core:proof:019:end}\end{proof}

\begin{figure}[htbp]
\centering
\includegraphics[width=0.76\textwidth]{11_circle_low_mode_convergence.pdf}
\caption{The first nonzero circle mode converges at second order and remains below the proved dispersion bound.}
\label{core:fig:circle-convergence}
\end{figure}

\begin{certificate}[Continuum dispersion]\label{core:certificate:011}
The observed convergence order was $1.99936$. The maximum error-minus-bound value was negative, so no tested mode violated the theorem.
\end{certificate}

\subsection{One- and two-dimensional plateaus}

\begin{figure}[htbp]
\centering
\includegraphics[width=0.80\textwidth]{12_spectral_dimension.pdf}
\caption{Intermediate spectral-dimension windows for a finite cycle and a finite two-dimensional torus. The ultraviolet and infrared departures are finite-size effects, not additional dimensions.}
\label{core:fig:spectral-dimension}
\end{figure}

\begin{certificate}[Spectral dimension]\label{core:certificate:012}
The median intermediate plateaus were
\[
d_s^{(1)}=1.00062,
\qquad
d_s^{(2)}=2.03781.
\]
These tests recover the expected dimensions of the chosen finite geometries; they do not select a physical quantum-spacetime theory.
\end{certificate}

\section{Constraint packets and the quantum-spacetime boundary}\label{core:section:012}

At one finite level, let $\widehat C_A$ be a finite family of matrix constraints, with fixed scalar structure coefficients, and let $\Pi^{\mathrm{phys}}$ be the orthogonal projector onto their exact common kernel. If the classical constraint algebra has structure coefficients $f_{AB}^{\ \ C}$, define the anomaly operator
\[
\boxed{
\mathfrak A_{AB}
=
[\widehat C_A,\widehat C_B]
-i\sum_C f_{AB}^{\ \ C}\widehat C_C.
}
\]
Its physical packet is
\[
\Pi^{\mathrm{phys}}\mathfrak A_{AB}\Pi^{\mathrm{phys}}.
\]
In this exact finite setting, $\widehat C_C\Pi^{\mathrm{phys}}=0$ for every $C$. Associativity gives
\[
\mathfrak A_{AB}\Pi^{\mathrm{phys}}=0,
\qquad
\Pi^{\mathrm{phys}}\mathfrak A_{AB}\Pi^{\mathrm{phys}}=0
\]
automatically. This compression therefore cannot certify operator closure away from the physical range. A nonzero measured compression contradicts the exact-kernel assumptions or indicates an approximation or implementation error. For example, $C_1=\operatorname{diag}(1,0,0)$ and $C_2=E_{12}+E_{21}$ annihilate the same third coordinate, while their commutator $E_{12}-E_{21}$ is nonzero. With zero proposed structure coefficients, the anomaly is nonzero but its physical compression vanishes.

The closure target must be stated at the operator level, not inferred from that automatic compression. Unbounded constraints additionally require a specified common invariant domain. Neither such a domain theory nor gravitational structure functions are supplied by the finite check.

\begin{conjecture}[Anomaly-free refinement]\label{core:conjecture:001}
For a specified nontrivial quantum-geometric constraint model, there exists a coherent packet tower with nonzero physical spaces, compatible physical projectors and exact operator closure
\[
\mathfrak A^{(\ell)}_{AB}=0
\]
at every finite level. The compatible physical compression then vanishes as a consequence, not as a substitute for closure.
\end{conjecture}
The operator condition is the substantive research target in this edition; the earlier compression-only wording was too weak. A choice of model, dynamics and refinement maps is still needed before this program becomes a fully specified physical existence problem.

\begin{conjecture}[Uniform Feshbach control]\label{core:conjecture:002}
For the same specified tower, the self-energies from exact elimination of fine packets admit bounds uniform in cutoff and volume on a common spectral domain. Together with consistent Hilbert-space identifications and a proved resolvent convergence criterion, they yield a nontrivial strong-resolvent continuum limit. Uniform boundedness alone is not asserted to imply convergence.
\end{conjecture}

\begin{conjecture}[Lorentzian reconstruction]\label{core:conjecture:003}
The Euclidean packet theory satisfies a positivity or reconstruction principle strong enough to produce a Lorentzian Hilbert space with causal propagation and positive physical norm.
\end{conjecture}

\begin{conjecture}[Einstein infrared limit]\label{core:conjecture:004}
There exist semiclassical packet states $\Psi_{g,\ell}$ such that the effective action converges to
\[
\frac1{16\pi G}
\int\sqrt{-g}\,(R-2\Lambda)\,d^4x
\]
with controlled higher-order corrections.
\end{conjecture}

\begin{conjecture}[Massless spin-two sector]\label{core:conjecture:005}
The continuum two-point function possesses a positive transverse-traceless massless spin-two pole, while all unphysical packet sectors are removed by the constraints.
\end{conjecture}

None of these physical conclusions is established by the exact finite identities. No formal logical-independence theorem is claimed. Exact packet reconstruction is a kinematic and computational advantage; it does not determine the physical action, prove continuum existence, or guarantee the general-relativistic limit.

\begin{figure}[htbp]
\centering
\includegraphics[width=0.86\textwidth]{18_proof_boundary.pdf}
\caption{The proof boundary. Finite packet, circuit, leakage, Feshbach, connection, and spectral identities are established here. Constraint closure and the physical continuum claims remain conjectural.}
\label{core:fig:proof-boundary}
\end{figure}

\section{Numerical certificates}\label{core:section:013}

The original manuscript reports CPU experiments and a 22-test suite; it makes no GPU claim. The values below reproduce the original historical report. They are numerical observations and are not counted among the 323 declarations in the selected finite formal checkpoint. The recovered numerical source and raw records accompany the source corpus; the present reading copy distinguishes those observations from mathematical proof. A displayed floating-point residual of zero denotes the reported computation, not an independent proof of exact arithmetic equality.

\begin{table}[htbp]
\centering
\caption{Reported finite numerical checks, retained from the original manuscript and not rerun for this edition.}
\label{core:tab:certificates}
\begin{tabular}{L{0.64\textwidth}r}
\toprule
Certificate & Value\\
\midrule
Bell-state residual & $0$\\
Three-qubit seed-orbit rank & $8/8$\\
Parity off-block norm & $0$\\
Circuit path reconstruction residual & $0$\\
Exact-symmetry maximum leakage & $0$\\
Maximum leakage-bound violation & $0$\\
Feshbach projected-resolvent residual & $3.2291\times10^{-17}$\\
Bare-truncation projected-resolvent residual & $1.3091\times10^{-2}$\\
Moving-frame residual with connection & $1.0645\times10^{-10}$\\
Moving-frame residual without connection & $7.0000\times10^{-1}$\\
Canonical packet-connection residual & $0$\\
Bloch-curvature residual & $2.0803\times10^{-11}$\\
Kraus completeness residual & $0$\\
Output trace residual & $0$\\
Knill--Laflamme residual & $0$\\
Bit-flip recovery residual & $1.0385\times10^{-16}$\\
Tensor parity-refinement residual & $0$\\
Circle low-mode convergence order & $1.99936$\\
One-dimensional spectral plateau & $1.00062$\\
Two-dimensional spectral plateau & $2.03781$\\
Minimum Euclidean-transfer eigenvalue & $0.0767054$\\
Transfer semigroup residual & $0$\\
Volume-packet reconstruction residual & $0$\\
\bottomrule
\end{tabular}
\end{table}

The original reported count is
\[
\boxed{22/22\text{ historical numerical tests, not rerun here}.}
\]
The 19 original figures are retained from the recovered numerical companion. The separate visual guide supplies a reading route through these constructions and the later strengthened results.

\setcounter{section}{14}
\section{Conclusion}\label{core:section:014}

The quantum realization of the seed calculus can now be stated precisely. A primitive state is addressed by a declared operator grammar; symmetry projectors split the state and every operator into exact packets; circuit products become sums over packet paths; Hamiltonian symmetry and its controlled breaking are visible as diagonal blocks and leakage; moving packets carry a canonical connection and curvature; channels, measurements, and error-correction conditions admit exact packet formulations; and eliminated sectors return through self-energy and memory rather than disappearing.

The finite theory therefore provides a coherent language for
\[
\boxed{
\begin{gathered}
\text{seed reachability}
+
\text{symmetry resolution}
+
\text{circuit paths}\\
+
\text{effective dynamics}
+
\text{geometric transport}
+
\text{constraint diagnostics}
\end{gathered}
}
\]
What it does not yet provide is the physical continuum theorem. The next decisive results would be anomaly-free refinement, uniform effective-operator control, Lorentzian reconstruction, the Einstein infrared limit, and a positive massless spin-two sector. Those requirements are explicit enough to be tested and too substantial to be replaced by finite numerical resemblance.



\startreadingpart{Strengthened quantum packet calculus and continuum towers}
\graphicspath{{assets/strengthened/}}
\section{Introduction}\label{strong:section:001}\label{strong:sec:introduction}

Quantum mechanics repeatedly uses the same operation. A Hilbert space is split into sectors, an operator is resolved between those sectors, and the pieces are propagated or recombined. Symmetry sectors, eigenspaces, syndrome spaces, adiabatic subspaces, and retained model spaces are familiar examples. This article treats them as instances of one finite packet calculus.

The elementary coordinate object is a complete family of orthogonal projectors
\[
P_aP_b=\delta_{ab}P_a,
\qquad
\sum_{a\in I}P_a=I.
\]
Every vector has packet coordinates \(P_a\psi\). Every operator has transfer blocks
\[
A_{ba}=P_bAP_a.
\]
These identities are simple. Their value appears when the same coordinates are used across representation theory, circuits, channels, error correction, effective dynamics, and geometric limits.

Finite dimensional quantum information and operator theory supply the standard background \cite{NielsenChuang2010,Watrous2018,Bhatia1997}. Finite group projectors come from character theory \cite{Serre1977,FultonHarris1991}. Moving spectral subspaces and geometric phase originate in the work of Kato, Berry, Wilczek, and Zee \cite{Kato1950,Kato1995,Berry1984,WilczekZee1984}. Completely positive maps and quantum dynamical semigroups are treated in \cite{Stinespring1955,Choi1975,Kraus1983,Lindblad1976}. Exact error correction is governed by the Knill Laflamme condition \cite{KnillLaflamme1997,SchumacherNielsen1996,BenyKempfKribs2007}. Projection methods and hidden state memory are linked by the Feshbach, Nakajima, and Zwanzig constructions \cite{Feshbach1958,Nakajima1958,Zwanzig1960,BravyiDiVincenzoLoss2011}.

The first part begins with quantum seeds. A state is cyclic only relative to a declared observable algebra. For an algebra in multiplicity form, the local orbit dimension is controlled by the rank of a coefficient matrix. This recovers the relation between Schmidt rank and the dimension reached by local operations. The finite group section then constructs central isotypic projectors. Cyclic symmetry gives a sharper charge selection rule. An operator with charge \(q\) can move a state only from packet \(j\) to packet \(j+q\).

A circuit is resolved by inserting the identity between successive gates. This produces a sum over packet paths. The target probability contains diagonal path terms and interference between distinct paths. Symmetry can remove entire path families before any matrix multiplication is performed. Unitarity gives exact row and column relations among the transfer blocks.

The moving packet section treats time dependent projectors. The canonical generator
\[
K(t)=\sum_a\dot P_a(t)P_a(t)
\]
is skew adjoint and transports every packet. In the moving frame, the physical Hamiltonian acquires the geometric term \(-\ii W(t)^*\dot W(t)\). A leakage estimate shows that the relevant off packet coupling is not \(QHP\) alone. It is
\[
Q(t)\bigl(H(t)-\ii K(t)\bigr)P(t).
\]
This is the exact finite distinction between following a moving subspace and remaining in a fixed one.

Density matrices require two packet labels. The superoperators
\[
\Pcal_{ab}(X)=P_aXP_b
\]
form a complete orthogonal packet system on Hilbert Schmidt operator space. Dephasing is the diagonal projection. Every channel has input and output superblocks, with the same composition law as ordinary matrix blocks. A covariant channel preserves the isotypic packets of the conjugation action. In the cyclic case these are coherence orders. The Choi matrix carries the equivalent symmetry under \(U\otimes\overline U\).

Error correction is expressed through compression to a code. Let \(P\) be the code projector. The class of \(PAP\) modulo scalar multiples of \(P\) is the obstruction to making \(A\) invisible inside the code. The Knill Laflamme criterion says that this obstruction vanishes for every product \(E_a^*E_b\) from the error family. Diagonalizing the resulting positive coefficient matrix produces orthogonal syndrome packets and an explicit recovery on the error span.

The effective dynamics section returns to ordinary Hilbert space packets. Schur complementation gives the exact retained resolvent and a frequency dependent self energy. The same hidden block gives a causal memory kernel in time. For a self adjoint hidden Hamiltonian, the self energy has a fixed sign in the upper half plane, and the time kernel is positive definite. Uniform spectral separation gives direct bounds that can be carried through a refinement tower.

The geometric part has two finite components. A positive transfer matrix gives a reflection positive Hankel kernel and reconstructs a nonnegative Hamiltonian through its logarithm. A finite Laplacian gives a heat trace and a scale dependent spectral dimension. The article proves fixed time convergence from periodic lattices to the continuum torus. Poisson summation then gives the continuum small time dimension.

Constraint closure needs a different obstruction. If \(C_a\) are self adjoint constraints and \(P_{\phys}\) projects onto their common kernel, then
\[
P_{\phys}[C_a,C_b]P_{\phys}=0
\]
holds automatically. This compression cannot test closure. The correct finite obstruction is the class of \([C_a,C_b]\) in the quotient of \(\End(\Hcal)\) by the span of the constraints. Vanishing means that the constraint span is closed under commutators.

The last part gives two bounded continuum tower theorems. Exact compatible finite Hamiltonians define a bounded self adjoint operator on the Hilbert direct limit. A norm summable sequence in a common Hilbert space converges to a bounded self adjoint limit, and its unitary and Euclidean evolutions converge uniformly on bounded time intervals. Reflection positivity and fixed algebraic closure pass to limits under the stated convergence assumptions.

These results do not construct an interacting continuum quantum field theory. They do not prove anomaly freedom on an unbounded common domain, Osterwalder Schrader reconstruction for a continuum measure, an Einstein infrared action, or a positive massless spin two sector. Each of those conclusions needs hypotheses absent from the finite calculations.

\begin{figure}[t]
\centering
\includegraphics[width=\textwidth]{01_quantum_architecture.pdf}
\caption{The finite mathematical layers developed in the article and the separate limit criteria needed at the continuum boundary.}
\label{strong:fig:architecture}
\end{figure}

\begin{table}[t]
\centering
\small
\begin{tabular}{>{\raggedright\arraybackslash}p{.20\textwidth}>{\raggedright\arraybackslash}p{.31\textwidth}>{\raggedright\arraybackslash}p{.39\textwidth}}
\toprule
Object & Packet coordinate & Exact structural condition \\
\midrule
State & \(P_a\psi\) & completeness of the projector family \\
Operator & \(P_bAP_a\) & block multiplication \\
Circuit & intermediate packet path & insertion of complete resolutions \\
Density matrix & \(P_a\rho P_b\) & two sided operator packets \\
Channel & superoperator block & complete positivity and covariance when declared \\
Code & syndrome packet & scalar compression of error products \\
Effective model & Feshbach map & invertibility of the hidden resolvent \\
Euclidean model & positive transfer matrix & positivity and self adjointness \\
Constraint system & quotient obstruction & commutator closure modulo the constraint span \\
Continuum tower & direct or norm limit & compatibility plus explicit convergence \\
\bottomrule
\end{tabular}
\caption{The packet coordinates used in the article.}
\label{strong:tab:dictionary}
\end{table}

\section{Quantum seeds and observable algebras}\label{strong:section:002}\label{strong:sec:seeds}

\begin{definition}[Finite quantum seed system]\label{strong:definition:001}\label{strong:def:quantum_seed}
A finite quantum seed system is a triple
\[
(\Hcal,\Acal,S),
\]
where \(\Hcal\) is a finite dimensional Hilbert space, \(\Acal\subseteq\End(\Hcal)\) is a unital star subalgebra, and \(S\subseteq\Hcal\) is a finite set. Its generated hull is
\[
\Acal S
=
\operatorname{span}\{A\psi:A\in\Acal,\ \psi\in S\}.
\]
The seed set is cyclic when \(\Acal S=\Hcal\). The quantum seed rank is the least cardinality of a cyclic seed set.
\end{definition}

The algebra is part of the data. A vector may be cyclic for \(\End(\Hcal)\) and fail to be cyclic for a restricted observable or gate algebra.

Let
\[
\Acal
\cong
\bigoplus_{r=1}^{s}
\bigl(M_{d_r}(\CC)\otimes I_{m_r}\bigr)
\]
act on
\[
\Hcal
\cong
\bigoplus_{r=1}^{s}
\bigl(\CC^{d_r}\otimes\CC^{m_r}\bigr).
\]
Assume \(d_r,m_r\ge1\) for each displayed block. For an empty block family, the seed rank and the displayed finite supremum are defined as zero. A vector in one block is identified with a \(d_r\) by \(m_r\) coefficient matrix.

\begin{theorem}[Orbit dimension of one seed]\label{strong:theorem:001}\label{strong:thm:orbit_dimension}
For integers \(d\ge1\), \(m\ge0\), and the action of \(M_d(\CC)\otimes I_m\) on \(\CC^d\otimes\CC^m\), let \(C_\psi\) be the coefficient matrix of \(\psi\). Then
\[
\dim\operatorname{span}\{(A\otimes I_m)\psi:A\in M_d(\CC)\}
=
 d\rank C_\psi.
\]
The vector is cyclic exactly when \(\rank C_\psi=m\).
\end{theorem}

\begin{proof}\phantomsection\label{strong:proof:001}
Under the coefficient matrix identification, \((A\otimes I_m)\psi\) corresponds to \(AC_\psi\). Every row of \(AC_\psi\) lies in the row space of \(C_\psi\). Conversely, a matrix whose rows lie in that row space equals \(AC_\psi\) for a suitable \(A\). There are \(d\) independent row choices in a space of dimension \(\rank C_\psi\). The orbit dimension is therefore \(d\rank C_\psi\). It equals \(dm\) exactly when \(C_\psi\) has full column rank.
\phantomsection\label{strong:proof:001:end}\end{proof}

\begin{theorem}[Quantum seed rank]\label{strong:theorem:002}\label{strong:thm:quantum_seed_rank}
For the displayed multiplicity decomposition,
\[
\operatorname{srank}_{\Acal}(\Hcal)
=
\max_{1\le r\le s}
\left\lceil\frac{m_r}{d_r}\right\rceil.
\]
\end{theorem}

\begin{proof}\phantomsection\label{strong:proof:002}
A single seed contributes a row space of dimension at most \(d_r\) to the \(r\)th multiplicity factor. Hence \(k\) seeds cover at most \(kd_r\) multiplicity directions, which gives the lower bound.

Set
\[
k=\max_r\left\lceil\frac{m_r}{d_r}\right\rceil.
\]
For each \(r\), partition a basis of \(\CC^{m_r}\) into at most \(k\) groups of size at most \(d_r\). Choose the coefficient matrix of the \(j\)th seed so that its row space spans the \(j\)th group. The combined row spaces span every multiplicity factor. The preceding theorem then shows that the \(k\) seeds generate every summand.
\phantomsection\label{strong:proof:002:end}\end{proof}

\begin{corollary}[Local orbit and Schmidt rank]\label{strong:corollary:001}\label{strong:cor:schmidt_orbit}
Let \(\psi\in\Hcal_A\otimes\Hcal_B\), let \(d_A=\dim\Hcal_A\ge1\), and let \(r\) be the Schmidt rank of \(\psi\). Then
\[
\dim\operatorname{span}\{(A\otimes I_B)\psi:A\in\End(\Hcal_A)\}
=
 d_A r.
\]
Local operators on \(A\) generate the whole bipartite Hilbert space exactly when the reduced state on \(B\) has full support.
\end{corollary}

\begin{proof}\phantomsection\label{strong:proof:003}
The Schmidt rank is the rank of the coefficient matrix in any product basis. Apply \cref{strong:thm:orbit_dimension}.
\phantomsection\label{strong:proof:003:end}\end{proof}

\begin{figure}[t]
\centering
\includegraphics[width=.68\textwidth]{02_seed_rank.pdf}
\caption{The dimension reached by all local operators is linear in Schmidt rank.}
\label{strong:fig:seed_rank}
\end{figure}

\begin{proposition}[Computational address orbit]\label{strong:proposition:001}\label{strong:prop:computational_orbit}
For \(n\) qubits, let
\[
\psi_0=\ket{0}^{\otimes n}
\]
and let \(X_j\) be the Pauli \(X\) operator on the \(j\)th qubit. Then
\[
\left\{
X_1^{a_1}\cdots X_n^{a_n}\psi_0:
(a_1,\ldots,a_n)\in\{0,1\}^n
\right\}
\]
is the computational basis.
\end{proposition}

\begin{proof}\phantomsection\label{strong:proof:004}
The addressed vector is \(\ket{a_1\cdots a_n}\). Distinct binary addresses produce the \(2^n\) distinct basis vectors.
\phantomsection\label{strong:proof:004:end}\end{proof}

\begin{remark}\label{strong:remark:001}\label{strong:rem:seed_scope}
The seed rank records generation by a declared algebra. It does not measure circuit depth, locality cost, control complexity, or physical reachability under a restricted time dependent Hamiltonian.
\end{remark}

\section{Symmetry packets and selection rules}\label{strong:section:003}\label{strong:sec:symmetry}

Let \(G\) be a finite group and let
\[
U:G\longrightarrow\Ucal(\Hcal)
\]
be a unitary representation. Write \(\widehat G\) for the set of irreducible unitary representations. For \(\lambda\in\widehat G\), let \(d_\lambda\) and \(\chi_\lambda\) denote its dimension and character.

\begin{theorem}[Finite group isotypic projectors]\label{strong:theorem:003}\label{strong:thm:isotypic_projectors}
The operators
\[
P_\lambda
=
\frac{d_\lambda}{\abs{G}}
\sum_{g\in G}
\overline{\chi_\lambda(g)}U(g)
\]
are orthogonal projectors. They satisfy
\[
P_\lambda P_\mu=\delta_{\lambda\mu}P_\lambda,
\qquad
\sum_{\lambda\in\widehat G}P_\lambda=I_{\Hcal}.
\]
The range of \(P_\lambda\) is the \(\lambda\) isotypic component of \(\Hcal\).
\end{theorem}

\begin{proof}\phantomsection\label{strong:proof:005}
Complete reducibility gives a unitary decomposition
\[
\Hcal
\cong
\bigoplus_{\mu\in\widehat G}
V_\mu\otimes M_\mu,
\qquad
U(g)
\cong
\bigoplus_\mu
\pi_\mu(g)\otimes I_{M_\mu}.
\]
On the \(\mu\) summand, the displayed average is central in the irreducible representation \(\pi_\mu\), so it is a scalar by Schur's lemma. Taking its trace and applying character orthogonality shows that the scalar is one when \(\mu=\lambda\) and zero otherwise. The asserted identities follow on the direct sum.
\phantomsection\label{strong:proof:005:end}\end{proof}

The central projectors separate irreducible types but do not choose coordinates inside the representation factor. Matrix units provide that finer resolution.

\begin{theorem}[Representation matrix units]\label{strong:theorem:004}\label{strong:thm:representation_matrix_units}
Fix an irreducible unitary representation \(\pi_\lambda:G\to U(d_\lambda)\) and an orthonormal basis of its representation space. Define
\[
E_{ij}^{\lambda}
=
\frac{d_\lambda}{\abs{G}}
\sum_{g\in G}
\bigl[\pi_\lambda(g^{-1})\bigr]_{ji}U(g).
\]
Then
\[
E_{ij}^{\lambda}E_{k\ell}^{\mu}
=
\delta_{\lambda\mu}\delta_{jk}E_{i\ell}^{\lambda},
\qquad
\bigl(E_{ij}^{\lambda}\bigr)^*=E_{ji}^{\lambda},
\]
and
\[
P_\lambda=\sum_{i=1}^{d_\lambda}E_{ii}^{\lambda}.
\]
\end{theorem}

\begin{proof}\phantomsection\label{strong:proof:006}
Apply the matrix coefficient form of Schur orthogonality to the two group averages. It gives the matrix unit multiplication law. Unitarity gives the adjoint relation after replacing \(g\) by \(g^{-1}\). Summing the diagonal coefficients gives the character average in \cref{strong:thm:isotypic_projectors}.
\phantomsection\label{strong:proof:006:end}\end{proof}

\subsection{Cyclic charge packets}\label{strong:subsec:cyclic_charge}

Let \(N\ge1\) be an integer, \(S\) unitary with \(S^N=I\), and \(\omega=\e^{2\pi\ii/N}\).

\begin{corollary}[Cyclic charge projectors]\label{strong:corollary:002}\label{strong:cor:cyclic_projectors}
The operators
\[
P_j
=
\frac1N
\sum_{r=0}^{N-1}
\omega^{-jr}S^r,
\qquad
j\in\ZZ/N\ZZ,
\]
form a complete orthogonal packet system and satisfy
\[
SP_j=\omega^jP_j.
\]
\end{corollary}

\begin{proof}\phantomsection\label{strong:proof:007}
Apply \cref{strong:thm:isotypic_projectors} to the cyclic group. Its irreducible representations are the one dimensional characters \(r\mapsto\omega^{jr}\).
\phantomsection\label{strong:proof:007:end}\end{proof}

\begin{definition}[Operator charge]\label{strong:definition:002}\label{strong:def:operator_charge}
An operator \(A\in\End(\Hcal)\) has cyclic charge \(q\in\ZZ/N\ZZ\) when
\[
SAS^*=\omega^qA.
\]
\end{definition}

\begin{theorem}[Charge selection rule]\label{strong:theorem:005}\label{strong:thm:charge_selection}
If \(A\) has charge \(q\), then
\[
P_kAP_j=0
\]
unless \(k=j+q\) in \(\ZZ/N\ZZ\). In fact,
\[
AP_j=P_{j+q}AP_j.
\]
\end{theorem}

\begin{proof}\phantomsection\label{strong:proof:008}
For \(x\in P_j\Hcal\),
\[
S(Ax)
=(SAS^*)Sx
=\omega^{q+j}Ax.
\]
Thus \(Ax\in P_{j+q}\Hcal\). The block statement follows.
\phantomsection\label{strong:proof:008:end}\end{proof}

\begin{corollary}[Charge addition]\label{strong:corollary:003}\label{strong:cor:charge_addition}
If \(A\) and \(B\) have charges \(q\) and \(r\), then \(AB\) has charge \(q+r\). The adjoint \(A^*\) has charge \(-q\).
\end{corollary}

\begin{proof}\phantomsection\label{strong:proof:009}
Conjugation by \(S\) is a star automorphism. Apply it to \(AB\) and \(A^*\).
\phantomsection\label{strong:proof:009:end}\end{proof}

\begin{figure}[t]
\centering
\includegraphics[width=.82\textwidth]{03_symmetry_selection.pdf}
\caption{A charge one operator moves each cyclic packet to the next packet and has no other transfer blocks.}
\label{strong:fig:symmetry_selection}
\end{figure}

\subsection{Tensor refinement}\label{strong:subsec:tensor_refinement}

Suppose the \(a\)th tensor factor carries a cyclic packet family
\[
\bigl(P_j^{(a)}\bigr)_{j\in\ZZ/N\ZZ}.
\]
The fine tensor packets are
\[
P_{\bm j}^{\rm fine}
=
P_{j_1}^{(1)}\otimes\cdots\otimes P_{j_m}^{(m)}.
\]

\begin{theorem}[Tensor charge refinement]\label{strong:theorem:006}\label{strong:thm:tensor_refinement}
For \(r\in\ZZ/N\ZZ\), define
\[
P_r^{\rm total}
=
\sum_{j_1+\cdots+j_m=r}
P_{j_1}^{(1)}\otimes\cdots\otimes P_{j_m}^{(m)}.
\]
Then the \(P_r^{\rm total}\) form a complete orthogonal packet system. Every total charge packet is the exact sum of its fine tensor descendants.
\end{theorem}

\begin{proof}\phantomsection\label{strong:proof:010}
The fine tensor projectors are mutually orthogonal and sum to the identity. The total residue classes partition the fine address set. Summing the fine projectors over each class gives the result.
\phantomsection\label{strong:proof:010:end}\end{proof}

\begin{corollary}[Tensor selection rule]\label{strong:corollary:004}\label{strong:cor:tensor_selection}
If \(A_a\) has local charge \(q_a\), then
\[
A_1\otimes\cdots\otimes A_m
\]
has total charge \(q_1+\cdots+q_m\).
\end{corollary}

\begin{proof}\phantomsection\label{strong:proof:011}
Conjugate the tensor product by the total cyclic symmetry and multiply the local phase factors.
\phantomsection\label{strong:proof:011:end}\end{proof}

\begin{remark}\label{strong:remark:002}\label{strong:rem:nonabelian_fusion}
For a nonabelian group, tensor products are refined by irreducible fusion multiplicities rather than residue addition. The central isotypic projectors remain exact, but a complete tensor refinement also requires a choice of Clebsch Gordan intertwiners.
\end{remark}

\section{Quantum circuits as packet paths}\label{strong:section:004}\label{strong:sec:circuits}

Let
\[
U=U_mU_{m-1}\cdots U_1
\]
be a finite circuit with \(m\ge1\), and let \((P_a)_{a\in I}\) be a complete orthogonal packet family. For the empty circuit, \(U=I\) and its transfer block is \(P_bP_a=\delta_{ba}P_a\), whose endpoint sum reconstructs \(I\).

\begin{theorem}[Exact circuit path expansion]\label{strong:theorem:007}\label{strong:thm:circuit_path}
For source packet \(a_0\) and target packet \(a_m\),
\[
P_{a_m}UP_{a_0}
=
\sum_{a_1,\ldots,a_{m-1}}
P_{a_m}U_mP_{a_{m-1}}
\cdots
P_{a_1}U_1P_{a_0}.
\]
Summing over \(a_0\) and \(a_m\) reconstructs \(U\).
\end{theorem}

\begin{proof}\phantomsection\label{strong:proof:012}
Insert \(I=\sum_aP_a\) between every adjacent pair of gates and distribute the finite sums. Completeness at the two endpoints gives the full operator.
\phantomsection\label{strong:proof:012:end}\end{proof}

For a packet path
\[
\gamma=(a_0,a_1,\ldots,a_m),
\]
write
\[
U_\gamma
=
P_{a_m}U_mP_{a_{m-1}}
\cdots
P_{a_1}U_1P_{a_0}.
\]

\begin{theorem}[Packet unitarity identities]\label{strong:theorem:008}\label{strong:thm:packet_unitarity}
If \(U\) is unitary and \(U_{ba}=P_bUP_a\), then
\[
\sum_bU_{ba}^*U_{bc}=\delta_{ac}P_a
\]
and
\[
\sum_aU_{ba}U_{da}^*=\delta_{bd}P_b.
\]
\end{theorem}

\begin{proof}\phantomsection\label{strong:proof:013}
Use completeness and unitarity:
\[
\sum_bU_{ba}^*U_{bc}
=
P_aU^*\left(\sum_bP_b\right)UP_c
=P_aP_c.
\]
The second identity follows in the same way.
\phantomsection\label{strong:proof:013:end}\end{proof}

\begin{corollary}[Probability conservation from one packet]\label{strong:corollary:005}\label{strong:cor:packet_probability}
If \(\psi\in P_a\Hcal\), then
\[
\sum_b\norm{P_bU\psi}^2=\norm{\psi}^2.
\]
\end{corollary}

\begin{proof}\phantomsection\label{strong:proof:014}
Take the inner product of the first identity in \cref{strong:thm:packet_unitarity} with \(\psi\) and set \(c=a\).
\phantomsection\label{strong:proof:014:end}\end{proof}

\begin{theorem}[Path interference formula]\label{strong:theorem:009}\label{strong:thm:path_interference}
Fix a source vector \(\psi\in P_{a_0}\Hcal\) and a target packet \(b\). Let \(\Gamma_b\) be the set of packet paths from \(a_0\) to \(b\), and set
\[
\eta_\gamma=U_\gamma\psi.
\]
Then
\[
\norm{P_bU\psi}^2
=
\sum_{\gamma\in\Gamma_b}\norm{\eta_\gamma}^2
+
2\operatorname{Re}
\sum_{\substack{\gamma,\gamma'\in\Gamma_b\\ \gamma<\gamma'}}
\ip{\eta_\gamma}{\eta_{\gamma'}}.
\]
Consequently,
\[
\norm{P_bU\psi}
\leq
\sum_{\gamma\in\Gamma_b}\norm{\eta_\gamma}.
\]
If the path outputs are pairwise orthogonal, the target probability is the sum of the path probabilities.
\end{theorem}

\begin{proof}\phantomsection\label{strong:proof:015}
The path expansion gives
\[
P_bU\psi=\sum_{\gamma\in\Gamma_b}\eta_\gamma.
\]
Expand the squared norm and group conjugate cross terms. The norm bound is the triangle inequality. Orthogonality removes every cross term.
\phantomsection\label{strong:proof:015:end}\end{proof}

\begin{proposition}[Charge pruning of circuit paths]\label{strong:proposition:002}\label{strong:prop:charge_pruning}
Suppose every gate has a finite charge decomposition
\[
U_r=\sum_qU_r^{(q)},
\qquad
SU_r^{(q)}S^*=\omega^qU_r^{(q)}.
\]
A term with gate charges \(q_1,\ldots,q_m\) can contribute from source charge \(a_0\) only to target charge
\[
a_m=a_0+q_1+\cdots+q_m.
\]
Every packet path violating one of the intermediate charge updates vanishes.
\end{proposition}

\begin{proof}\phantomsection\label{strong:proof:016}
Apply \cref{strong:thm:charge_selection} after each gate component. The packet address is forced at every step.
\phantomsection\label{strong:proof:016:end}\end{proof}

\begin{figure}[t]
\centering
\includegraphics[width=.88\textwidth]{04_circuit_paths.pdf}
\caption{A circuit block is the coherent sum of its allowed packet paths. Selection rules remove forbidden paths before amplitudes are evaluated.}
\label{strong:fig:circuit_paths}
\end{figure}

\begin{example}[Bell circuit]\label{strong:example:001}\label{strong:ex:bell_circuit}
Starting from \(\ket{00}\), apply a Hadamard gate to the first qubit and then a controlled not gate. The final state is
\[
\frac{\ket{00}+\ket{11}}{\sqrt2}.
\]
It lies in the even parity packet. The intermediate state has one even and one odd computational contribution, so the final even amplitude is reconstructed by more than one packet path.
\end{example}

\section{Symmetry resolved dynamics and leakage}\label{strong:section:005}\label{strong:sec:dynamics}

Let \(H=H^*\) and let \((P_a)_{a\in I}\) be a complete orthogonal packet system.

\begin{theorem}[Exact invariant packets]\label{strong:theorem:010}\label{strong:thm:invariant_packets}
The following conditions are equivalent:
\begin{enumerate}
\item \([H,P_a]=0\) for every \(a\)
\item \(P_bHP_a=0\) whenever \(a\neq b\)
\item every packet \(P_a\Hcal\) is invariant under \(\e^{-\ii tH}\) for every real \(t\)
\end{enumerate}
When they hold,
\[
H=\sum_aP_aHP_a
\]
and
\[
\e^{-\ii tH}=\sum_aP_a\e^{-\ii tH}P_a.
\]
\end{theorem}

\begin{proof}\phantomsection\label{strong:proof:017}
Commutation is equivalent to vanishing of every off diagonal block. The same block decomposition passes to polynomials in \(H\), then to the exponential by norm convergence. Invariance of every packet implies invariance under both \(\e^{-\ii tH}\) and its adjoint. Differentiation at \(t=0\) gives commutation with \(H\).
\phantomsection\label{strong:proof:017:end}\end{proof}

Fix one packet projector \(P\) and set \(Q=I-P\).

\begin{lemma}[Off diagonal norm]\label{strong:lemma:001}\label{strong:lem:off_diagonal_norm}
For self adjoint \(H\),
\[
\norm{[H,P]}=\norm{QHP}=\norm{PHQ}.
\]
\end{lemma}

\begin{proof}\phantomsection\label{strong:proof:018}
Relative to \(P\Hcal\oplus Q\Hcal\),
\[
[H,P]
=
\begin{pmatrix}
0&-PHQ\\
QHP&0
\end{pmatrix}.
\]
Its square has diagonal blocks \(-PHQHP\) and \(-QHPHQ\). The norm is therefore the largest singular value of \(QHP\).
\phantomsection\label{strong:proof:018:end}\end{proof}

\begin{theorem}[Fixed packet leakage]\label{strong:theorem:011}\label{strong:thm:fixed_leakage}
For every real \(t\),
\[
\norm{Q\e^{-\ii tH}P}
\leq
\abs{t}\norm{QHP}
=
\abs{t}\norm{[H,P]}.
\]
\end{theorem}

\begin{proof}\phantomsection\label{strong:proof:019}
Let
\[
H_{\rm d}=PHP+QHQ,
\qquad
H_{\rm o}=PHQ+QHP.
\]
The propagator generated by \(H_{\rm d}\) preserves \(P\Hcal\) and \(Q\Hcal\). Duhamel's formula relative to \(H_{\rm d}\) gives
\[
Q\e^{-\ii tH}P
=
-\ii
\int_0^t
Q\e^{-\ii(t-s)H}H_{\rm o}\e^{-\ii sH_{\rm d}}P
\dd s.
\]
Since \(H_{\rm o}P=QHP\) and both propagators are unitary, the norm is at most \(\abs{t}\norm{QHP}\). Use \cref{strong:lem:off_diagonal_norm} for the last equality.
\phantomsection\label{strong:proof:019:end}\end{proof}

\begin{figure}[t]
\centering
\includegraphics[width=.72\textwidth]{05_leakage_bound.pdf}
\caption{Exact transition amplitude in a two packet model and the linear commutator bound.}
\label{strong:fig:leakage}
\end{figure}

\begin{corollary}[Leakage probability]\label{strong:corollary:006}\label{strong:cor:leakage_probability}
If \(\psi\in P\Hcal\) is a unit vector, then
\[
\norm{Q\e^{-\ii tH}\psi}^2
\leq
 t^2\norm{QHP}^2.
\]
\end{corollary}

\begin{proof}\phantomsection\label{strong:proof:020}
Apply \cref{strong:thm:fixed_leakage} and square the result.
\phantomsection\label{strong:proof:020:end}\end{proof}

\section{Moving packets and geometric transport}\label{strong:section:006}\label{strong:sec:moving_packets}

Let \((P_a(t))_{a\in I}\) be a twice continuously differentiable family of complete orthogonal packet systems on a finite dimensional Hilbert space.

\begin{theorem}[Canonical packet connection]\label{strong:theorem:012}\label{strong:thm:kato_connection}
Define
\[
K(t)
=
\sum_a\dot P_a(t)P_a(t).
\]
Then
\[
K(t)^*=-K(t)
\]
and
\[
[K(t),P_a(t)]=\dot P_a(t)
\]
for every \(a\).
\end{theorem}

\begin{proof}\phantomsection\label{strong:proof:021}
Differentiating \(P_a^2=P_a\) gives
\[
\dot P_aP_a+P_a\dot P_a=\dot P_a.
\]
Summing over \(a\) and using \(\sum_a\dot P_a=0\) yields
\[
\sum_aP_a\dot P_a=-\sum_a\dot P_aP_a.
\]
This proves skew adjointness. For fixed \(b\),
\[
KP_b=\dot P_bP_b.
\]
Differentiating \(P_bP_a=0\) for \(a\neq b\) gives
\[
P_b\dot P_a=-\dot P_bP_a.
\]
Hence
\[
P_bK=-\dot P_b(I-P_b).
\]
Subtracting gives \([K,P_b]=\dot P_b\).
\phantomsection\label{strong:proof:021:end}\end{proof}

\begin{theorem}[Exact Kato transport]\label{strong:theorem:013}\label{strong:thm:kato_transport}
Let \(W(t)\) solve
\[
\dot W(t)=K(t)W(t),
\qquad
W(0)=I.
\]
Then \(W(t)\) is unitary and
\[
P_a(t)W(t)=W(t)P_a(0)
\]
for every packet \(a\).
\end{theorem}

\begin{proof}\phantomsection\label{strong:proof:022}
Skew adjointness gives
\[
\frac{\dd}{\dd t}W(t)^*W(t)=0.
\]
Thus \(W(t)\) is unitary. Next,
\[
\frac{\dd}{\dd t}
\left(W(t)^*P_a(t)W(t)\right)
=
W(t)^*
\left(
\dot P_a-[K,P_a]
\right)
W(t)=0.
\]
The expression equals \(P_a(0)\) at \(t=0\), which proves the intertwining identity.
\phantomsection\label{strong:proof:022:end}\end{proof}

\begin{theorem}[Moving frame Hamiltonian]\label{strong:theorem:014}\label{strong:thm:moving_frame}
Let \(U(t)\) solve
\[
\ii\dot U(t)=H(t)U(t),
\qquad
U(0)=I,
\]
where \(H(t)=H(t)^*\) is continuous on the compact time interval under consideration. Set
\[
V(t)=W(t)^*U(t).
\]
Then
\[
\ii\dot V(t)=H_{\rm mov}(t)V(t),
\]
where
\[
H_{\rm mov}(t)
=
W(t)^*H(t)W(t)-\ii W(t)^*\dot W(t)
=
W(t)^*\bigl(H(t)-\ii K(t)\bigr)W(t).
\]
The moving frame Hamiltonian is self adjoint.
\end{theorem}

\begin{proof}\phantomsection\label{strong:proof:023}
Differentiate \(V=W^*U\). Since \(\dot W=KW\) and \(K^*=-K\),
\[
\dot W^*=-W^*K.
\]
Substitution gives the equation. The first term in \(H_{\rm mov}\) is self adjoint. The operator \(-\ii W^*\dot W\) is self adjoint because \(W^*\dot W\) is skew adjoint.
\phantomsection\label{strong:proof:023:end}\end{proof}

\begin{theorem}[Moving packet leakage]\label{strong:theorem:015}\label{strong:thm:moving_leakage}
Fix a packet \(P(t)=P_a(t)\) and put \(Q(t)=I-P(t)\). For \(0\leq t\leq T\),
\[
\norm{Q(t)U(t)P(0)}
\leq
\int_0^t
\norm{Q(s)\bigl(H(s)-\ii K(s)\bigr)P(s)}
\dd s.
\]
In particular, the physical evolution follows the moving packet exactly on \([0,T]\) when the displayed off packet block vanishes for every time in that interval.
\end{theorem}

\begin{proof}\phantomsection\label{strong:proof:024}
In the Kato frame, \(P(0)\) is fixed and
\[
Q(t)U(t)P(0)
=
W(t)Q(0)V(t)P(0).
\]
Apply the time dependent version of the Duhamel argument in \cref{strong:thm:fixed_leakage} to \(H_{\rm mov}(t)\). The off diagonal block is
\[
Q(0)H_{\rm mov}(s)P(0)
=
W(s)^*Q(s)\bigl(H(s)-\ii K(s)\bigr)P(s)W(s).
\]
Unitary conjugation preserves the norm.
\phantomsection\label{strong:proof:024:end}\end{proof}

The geometric term is therefore part of the exact coupling test. A Hamiltonian may be packet diagonal at every fixed time and still fail to transport the moving packet because the packet itself rotates.

\subsection{Discrete packet transport}\label{strong:subsec:discrete_transport}

Let \((P_a)\) and \((Q_a)\) be two complete packet systems with the same index set. Define
\[
S(\Qcal,\Pcal)=\sum_aQ_aP_a.
\]

\begin{theorem}[Direct packet rotation]\label{strong:theorem:016}\label{strong:thm:direct_rotation}
Assume
\[
\norm{S(\Qcal,\Pcal)-I}<1.
\]
Then \(S(\Qcal,\Pcal)\) is invertible. Its unitary polar factor
\[
\Ucal(\Qcal,\Pcal)
=
S(S^*S)^{-1/2}
\]
satisfies
\[
Q_a\Ucal(\Qcal,\Pcal)
=
\Ucal(\Qcal,\Pcal)P_a
\]
for every \(a\).
\end{theorem}

\begin{proof}\phantomsection\label{strong:proof:025}
The norm condition gives invertibility by the Neumann series. The identity
\[
Q_aS=SP_a=Q_aP_a
\]
shows that \(S^*S\) commutes with every \(P_a\). Its inverse square root also commutes with every \(P_a\). Hence
\[
Q_a\Ucal
=Q_aS(S^*S)^{-1/2}
=SP_a(S^*S)^{-1/2}
=\Ucal P_a.
\]
The polar factor is unitary because \(S\) is invertible.
\phantomsection\label{strong:proof:025:end}\end{proof}

\begin{theorem}[Convergence of short packet rotations]\label{strong:theorem:017}\label{strong:thm:short_rotation_convergence}
Let \((P_a(t))\) be twice continuously differentiable on \([0,T]\). For a partition
\[
0=t_0<t_1<\cdots<t_n=T,
\]
let
\[
W_\pi
=
\Ucal\bigl(\Pcal(t_n),\Pcal(t_{n-1})\bigr)
\cdots
\Ucal\bigl(\Pcal(t_1),\Pcal(t_0)\bigr).
\]
For sufficiently fine partitions the factors are defined, and
\[
\norm{W_\pi-W(T)}
\leq
C_T\abs{\pi}
\]
for a constant depending on the packet path but not on the partition.
\end{theorem}

\begin{proof}\phantomsection\label{strong:proof:026}
Taylor expansion gives
\[
S\bigl(\Pcal(t+h),\Pcal(t)\bigr)
=
I+hK(t)+O(h^2)
\]
uniformly on \([0,T]\). Since \(K(t)^*=-K(t)\),
\[
S^*S=I+O(h^2),
\]
so the polar factor satisfies
\[
\Ucal\bigl(\Pcal(t+h),\Pcal(t)\bigr)
=
I+hK(t)+O(h^2).
\]
This is a first order one step approximation to \(\dot W=KW\). A telescoping product estimate over the partition gives a total error bounded by a constant times \(\sum_j(t_{j+1}-t_j)^2\), which is at most \(T\abs{\pi}\).
\phantomsection\label{strong:proof:026:end}\end{proof}

\begin{figure}[t]
\centering
\includegraphics[width=.62\textwidth]{06_kato_transport.pdf}
\caption{A moving rank one packet on the Bloch sphere. The Kato evolution transports the range without choosing a phase convention for its eigenvector.}
\label{strong:fig:kato_transport}
\end{figure}

\subsection{Local frames and curvature}\label{strong:subsec:curvature}

Let \(P(x)\) be a smooth rank \(r\) orthogonal projector over a parameter manifold. Choose a local isometry
\[
V(x):\CC^r\longrightarrow\Hcal,
\qquad
V(x)V(x)^*=P(x).
\]

\begin{theorem}[Nonabelian packet connection]\label{strong:theorem:018}\label{strong:thm:nonabelian_connection}
The matrix valued one form
\[
\Acal=V^*\dd V
\]
is skew adjoint. Under a frame change \(V'=VG\), where \(G(x)\in U(r)\),
\[
\Acal'
=
G^*\Acal G+G^*\dd G.
\]
Its curvature
\[
\Fcal=\dd\Acal+\Acal\wedge\Acal
\]
satisfies
\[
\Fcal=V^*(\dd P\wedge\dd P)V.
\]
\end{theorem}

\begin{proof}\phantomsection\label{strong:proof:027}
Differentiate \(V^*V=I\) to obtain \(\Acal^*=-\Acal\). The gauge transformation law follows by differentiating \(VG\). Since \(P=VV^*\) and \((I-P)\dd V=\dd P\,V\), direct expansion gives
\[
\dd\Acal+\Acal\wedge\Acal
=
\dd V^*(I-P)\wedge\dd V
=
V^*(\dd P\wedge\dd P)V.
\]
\phantomsection\label{strong:proof:027:end}\end{proof}

\begin{corollary}[Bloch sphere curvature]\label{strong:corollary:007}\label{strong:cor:bloch_curvature}
Let
\[
P(\theta,\phi)
=
\frac12\left(I+\bm n(\theta,\phi)\cdot\bm\sigma\right),
\]
where
\[
\bm n=(\sin\theta\cos\phi,\sin\theta\sin\phi,\cos\theta).
\]
With the convention
\[
\Fcal_{\theta\phi}
=-\ii\Tr\bigl(P[\partial_\theta P,\partial_\phi P]\bigr),
\]
one has
\[
\Fcal_{\theta\phi}=\frac12\sin\theta
\]
and
\[
\frac1{2\pi}
\int_0^{2\pi}\int_0^\pi
\Fcal_{\theta\phi}\dd\theta\dd\phi
=1.
\]
\end{corollary}

\begin{proof}\phantomsection\label{strong:proof:028}
The Pauli identity
\[
(\bm a\cdot\bm\sigma)(\bm b\cdot\bm\sigma)
=(\bm a\cdot\bm b)I
+\ii(\bm a\times\bm b)\cdot\bm\sigma
\]
gives
\[
-\ii\Tr\bigl(P[\partial_\theta P,\partial_\phi P]\bigr)
=\frac12
\bm n\cdot
(\partial_\theta\bm n\times\partial_\phi\bm n)
=\frac12\sin\theta.
\]
Integration gives one.
\phantomsection\label{strong:proof:028:end}\end{proof}

\section{Operator packets, measurements, and quantum channels}\label{strong:section:007}\label{strong:sec:channels}

Density matrices and observables carry both an input packet and an output packet. Equip \(\End(\Hcal)\) with the Hilbert Schmidt inner product
\[
\ip{X}{Y}_{\HS}=\Tr(X^*Y).
\]

\begin{definition}[Two sided operator packet]\label{strong:definition:003}\label{strong:def:operator_packet}
For a complete orthogonal family \((P_a)\), define
\[
\Pcal_{ab}(X)=P_aXP_b.
\]
The range
\[
\Bcal_{ab}=P_a\End(\Hcal)P_b
\]
is the operator packet with output address \(a\) and input address \(b\).
\end{definition}

\begin{theorem}[Complete operator packet system]\label{strong:theorem:019}\label{strong:thm:operator_packets}
The superoperators \((\Pcal_{ab})_{a,b\in I}\) are mutually orthogonal projections on Hilbert Schmidt operator space. They satisfy
\[
X=\sum_{a,b}P_aXP_b
\]
and
\[
\norm{X}_{\HS}^2
=
\sum_{a,b}\norm{P_aXP_b}_{\HS}^2.
\]
\end{theorem}

\begin{proof}\phantomsection\label{strong:proof:029}
Idempotence and completeness follow from the projector identities. For Hilbert Schmidt orthogonality,
\[
\begin{aligned}
\ip{P_aXP_b}{P_cYP_d}_{\HS}
&=
\Tr(P_bX^*P_aP_cYP_d)\\
&=
\delta_{ac}\delta_{bd}
\Tr(P_bX^*P_aYP_b).
\end{aligned}
\]
The Pythagorean identity follows.
\phantomsection\label{strong:proof:029:end}\end{proof}

\begin{figure}[t]
\centering
\includegraphics[width=.58\textwidth]{07_operator_packets.pdf}
\caption{A density matrix or observable carries two packet addresses. Diagonal packets contain populations and internal observables. Off diagonal packets contain coherences.}
\label{strong:fig:operator_packets}
\end{figure}

\begin{definition}[Packet dephasing]\label{strong:definition:004}\label{strong:def:dephasing}
The packet dephasing map is
\[
\Dcal_{\Pcal}(X)=\sum_aP_aXP_a.
\]
\end{definition}

\begin{theorem}[Dephasing as conditional expectation]\label{strong:theorem:020}\label{strong:thm:dephasing_expectation}
The map \(\Dcal_{\Pcal}\) is completely positive, trace preserving, unital, and idempotent. It is the Hilbert Schmidt orthogonal projection onto the block diagonal algebra
\[
\Acal_{\Pcal}
=
\{X:[X,P_a]=0\text{ for every }a\}.
\]
In particular,
\[
\norm{\Dcal_{\Pcal}(X)}_{\HS}
\leq
\norm{X}_{\HS}.
\]
\end{theorem}

\begin{proof}\phantomsection\label{strong:proof:030}
The projectors \(P_a\) form a Kraus family because \(\sum_aP_a^*P_a=I\). Thus the map is completely positive and trace preserving. It is unital for the same reason. Orthogonality and idempotence follow from \cref{strong:thm:operator_packets}. Its range consists exactly of the diagonal operator packets, which is the commutant of the projector family.
\phantomsection\label{strong:proof:030:end}\end{proof}

\begin{corollary}[Block diagonal measurement invariance]\label{strong:corollary:008}\label{strong:cor:measurement_invariance}
Let \((E_k)\) be a positive operator valued measurement, and suppose
\[
E_k=\Dcal_{\Pcal}(E_k)
\]
for every \(k\). Then every density matrix \(\rho\) satisfies
\[
\Tr(E_k\rho)
=
\Tr\bigl(E_k\Dcal_{\Pcal}(\rho)\bigr).
\]
\end{corollary}

\begin{proof}\phantomsection\label{strong:proof:031}
Since \(\Dcal_{\Pcal}\) is self adjoint for the Hilbert Schmidt inner product,
\[
\Tr(E_k\rho)
=
\Tr\bigl(\Dcal_{\Pcal}(E_k)\rho\bigr)
=
\Tr\bigl(E_k\Dcal_{\Pcal}(\rho)\bigr).
\]
\phantomsection\label{strong:proof:031:end}\end{proof}

\subsection{Channel superblocks}\label{strong:subsec:channel_blocks}

Let \(\Phi:\End(\Hcal_X)\to\End(\Hcal_Y)\) be linear. Choose packet systems \((P_a)\) on \(\Hcal_X\) and \((Q_c)\) on \(\Hcal_Y\). Define
\[
\Phi_{cd,ab}
=
\Qcal_{cd}\Phi\Pcal_{ab}.
\]

\begin{theorem}[Exact channel block calculus]\label{strong:theorem:021}\label{strong:thm:channel_blocks}
Every linear map has the exact reconstruction
\[
\Phi=\sum_{a,b,c,d}\Phi_{cd,ab}.
\]
If \(\Psi:\End(\Hcal_Y)\to\End(\Hcal_Z)\), then
\[
(\Psi\Phi)_{ef,ab}
=
\sum_{c,d}
\Psi_{ef,cd}\Phi_{cd,ab}.
\]
\end{theorem}

\begin{proof}\phantomsection\label{strong:proof:032}
Insert the complete operator packet resolutions before and after \(\Phi\). For composition, insert the middle resolution between \(\Psi\) and \(\Phi\) and collect the output and input packets.
\phantomsection\label{strong:proof:032:end}\end{proof}

\begin{standardfact}[Finite Kraus and Choi criteria]\label{strong:standardfact:001}\label{strong:fact:kraus_choi}
A linear map \(\Phi:M_n(\CC)\to M_m(\CC)\) is completely positive exactly when it has a Kraus representation
\[
\Phi(X)=\sum_rK_rXK_r^*.
\]
It is trace preserving exactly when
\[
\sum_rK_r^*K_r=I.
\]
Equivalently, its Choi matrix
\[
J(\Phi)
=
\sum_{i,j}\Phi(\ket{i}\bra{j})\otimes\ket{i}\bra{j}
\]
is positive, and trace preservation is
\[
\Tr_{\rm out}J(\Phi)=I_{\rm in}.
\]
\end{standardfact}

These finite criteria are classical \cite{Stinespring1955,Choi1975,Kraus1983,Watrous2018}.

\subsection{Covariance and coherence order}\label{strong:subsec:channel_covariance}

Let \(S\) be unitary and \(N\ge1\) an integer with \(S^N=I\), and define the conjugation action
\[
\Ad_S(X)=SXS^*.
\]

\begin{theorem}[Coherence order projectors]\label{strong:theorem:022}\label{strong:thm:coherence_projectors}
For \(q\in\ZZ/N\ZZ\), define
\[
\Lcal_q(X)
=
\frac1N
\sum_{r=0}^{N-1}
\omega^{-qr}S^rXS^{-r}.
\]
Then the \(\Lcal_q\) form a complete orthogonal packet system on Hilbert Schmidt operator space. Their ranges are the operators satisfying
\[
SXS^*=\omega^qX.
\]
If \(P_j\) are the state charge projectors, then
\[
\Ran\Lcal_q
=
\bigoplus_jP_{j+q}\End(\Hcal)P_j.
\]
\end{theorem}

\begin{proof}\phantomsection\label{strong:proof:033}
Apply the cyclic character projector formula to the unitary representation \(\Ad_S\) on operator space. The last identity follows from \cref{strong:thm:charge_selection} applied to each operator packet.
\phantomsection\label{strong:proof:033:end}\end{proof}

\begin{theorem}[Covariant channels preserve coherence order]\label{strong:theorem:023}\label{strong:thm:covariant_channel}
Suppose
\[
\Phi(SXS^*)=S\Phi(X)S^*
\]
for every \(X\). Then
\[
\Phi\Lcal_q=\Lcal_q\Phi
\]
for every \(q\). Thus \(\Phi\) cannot mix distinct coherence orders.
\end{theorem}

\begin{proof}\phantomsection\label{strong:proof:034}
The covariance identity says that \(\Phi\) commutes with every power of \(\Ad_S\). It therefore commutes with each finite character average \(\Lcal_q\).
\phantomsection\label{strong:proof:034:end}\end{proof}

\begin{theorem}[Choi symmetry criterion]\label{strong:theorem:024}\label{strong:thm:choi_covariance}
Let \(U:G\to\Ucal(\Hcal)\) be a finite group representation. A channel \(\Phi\) is covariant,
\[
\Phi\bigl(U(g)XU(g)^*\bigr)
=
U(g)\Phi(X)U(g)^*,
\]
if and only if
\[
\bigl(U(g)\otimes\overline{U(g)}\bigr)
J(\Phi)
\bigl(U(g)^*\otimes U(g)^{\mathsf T}\bigr)
=
J(\Phi)
\]
for every \(g\in G\).
\end{theorem}

\begin{proof}\phantomsection\label{strong:proof:035}
Let
\[
\Omega=\sum_i\ket{i}\otimes\ket{i}.
\]
Then
\[
J(\Phi)=(\Phi\otimes\id)(\ket{\Omega}\bra{\Omega}).
\]
The vector \(\Omega\) is invariant under \(U\otimes\overline U\). Covariance therefore implies the displayed Choi invariance. Conversely, the displayed transformation is the Choi matrix of
\[
\Ad_{U(g)}\circ\Phi\circ\Ad_{U(g)^*}.
\]
The Choi correspondence is injective, so invariance implies that this map equals \(\Phi\), which is covariance.
\phantomsection\label{strong:proof:035:end}\end{proof}

\begin{corollary}[Choi packet decomposition]\label{strong:corollary:009}\label{strong:cor:choi_packets}
The Choi matrix of a covariant channel is block diagonal with respect to the isotypic decomposition of \(U\otimes\overline U\).
\end{corollary}

\begin{proof}\phantomsection\label{strong:proof:036}
Apply \cref{strong:thm:isotypic_projectors} to the representation \(U\otimes\overline U\) and use \cref{strong:thm:choi_covariance}.
\phantomsection\label{strong:proof:036:end}\end{proof}

\begin{example}[Amplitude damping]\label{strong:example:002}\label{strong:ex:amplitude_damping}
For \(0\leq\gamma\leq1\), let
\[
K_0=
\begin{pmatrix}
1&0\\
0&\sqrt{1-\gamma}
\end{pmatrix},
\qquad
K_1=
\begin{pmatrix}
0&\sqrt\gamma\\
0&0
\end{pmatrix}.
\]
The channel
\[
\Phi_\gamma(\rho)=K_0\rho K_0^*+K_1\rho K_1^*
\]
is trace preserving. It sends the off diagonal packet \(\ket0\bra1\) to \(\sqrt{1-\gamma}\ket0\bra1\), while the excited population transfers partly into the ground population.
\end{example}

\begin{figure}[t]
\centering
\includegraphics[width=.82\textwidth]{08_channel_coherence.pdf}
\caption{Amplitude damping transfers population and contracts the coherence packet.}
\label{strong:fig:channel_coherence}
\end{figure}


\subsection{Quantum dynamical semigroups}\label{strong:subsec:lindblad}

\begin{standardfact}[Finite Lindblad generator]\label{strong:standardfact:002}\label{strong:fact:lindblad}
A norm continuous semigroup of completely positive trace preserving maps on a finite matrix algebra, in the Schr\"odinger picture, has a generator of the form
\[
\Gcal(X)
=
-\ii[H,X]
+
\sum_r
\left(
L_rXL_r^*
-
\frac12L_r^*L_rX
-
\frac12XL_r^*L_r
\right),
\]
where \(H=H^*\). Conversely, every operator of this form generates a completely positive trace preserving semigroup \cite{Lindblad1976}.
\end{standardfact}

\begin{theorem}[Charge preservation by a Lindblad generator]\label{strong:theorem:025}\label{strong:thm:lindblad_charge}
Let \(S\) be unitary and \(N\ge1\) an integer with \(S^N=I\). Assume
\[
SHS^*=H
\]
and, for every jump operator, there is a charge \(q_r\) such that
\[
SL_rS^*=\omega^{q_r}L_r.
\]
Then the Lindblad generator \(\Gcal\) preserves every coherence order:
\[
\Gcal\Lcal_p=\Lcal_p\Gcal.
\]
The quantum dynamical semigroup \(\e^{t\Gcal}\) has the same property.
\end{theorem}

\begin{proof}\phantomsection\label{strong:proof:037}
Let \(X\) have coherence order \(p\). The Hamiltonian commutator has order \(p\) because \(H\) has charge zero. The jump term \(L_rXL_r^*\) has charge
\[
q_r+p-q_r=p.
\]
Also \(L_r^*L_r\) has charge zero, so both anticommutator terms have charge \(p\). Thus \(\Gcal\) preserves \(\Ran\Lcal_p\). The exponential series then shows that \(\e^{t\Gcal}\) preserves the same packets.
\phantomsection\label{strong:proof:037:end}\end{proof}

\begin{corollary}[Independent coherence evolution]\label{strong:corollary:010}\label{strong:cor:no_coherence_mixing}
Under the hypotheses of \cref{strong:thm:lindblad_charge}, each coherence component of an initial operator evolves within its own order. In particular, a density matrix in the zero order sector remains in that sector.
\end{corollary}

\begin{proof}\phantomsection\label{strong:proof:038}
Decompose the initial operator by \cref{strong:thm:coherence_projectors} and apply the invariant packet statement to each component.
\phantomsection\label{strong:proof:038:end}\end{proof}

The theorem does not say that every diagonal quantum master equation closes to a classical Markov chain on packet populations. Scalar population closure needs additional isotropy inside each packet.

\section{Quantum error correction as a compression obstruction}\label{strong:section:008}\label{strong:sec:qec}

Let \(P\) be the orthogonal projector onto a code space
\[
\Ccal=P\Hcal,
\qquad
k=\dim\Ccal\ge1.
\]
The compression of \(A\in\End(\Hcal)\) to the code is \(PAP\in\End(\Ccal)\). The zero code has only the zero supported operator and no normalized code state; that branch is treated in the mathematical completion in Part III.

\begin{definition}[Code compression obstruction]\label{strong:definition:005}\label{strong:def:code_obstruction}
Let
\[
q_P:\End(\Ccal)\longrightarrow\End(\Ccal)/(\CC I_{\Ccal})
\]
be the quotient map. Define
\[
\Ob_P(A)=q_P(PAP).
\]
Equivalently, one may use the traceless representative
\[
\ob_P(A)
=
PAP-
\frac{\Tr(PAP)}{k}P.
\]
\end{definition}

The class vanishes exactly when \(A\) acts as a scalar after compression to the code. Such an operator cannot distinguish normalized code vectors by expectation value.

Exact correction of an error span means that one completely positive trace preserving recovery $\Rcal$ satisfies $\Rcal(F\rho F^*)=q_F\rho$ for every $F$ in that span and every code-supported matrix $\rho$, with a nonnegative scalar $q_F$ depending on $F$ but not on the logical input. For every normalized channel with Kraus operators in the same span, the same recovery restores the code input with total scalar one. Individual unnormalized errors are not required to have unit recovery weight.

\begin{theorem}[Knill Laflamme as obstruction vanishing]\label{strong:theorem:026}\label{strong:thm:kl_obstruction}
Let \(E_1,\ldots,E_r\) be an error family. The code \(\Ccal\) corrects their linear span exactly if and only if
\[
\Ob_P(E_a^*E_b)=0
\]
for all \(a,b\). Equivalently, there is a matrix \(c=(c_{ab})\) such that
\[
PE_a^*E_bP=c_{ab}P
\]
for all \(a,b\).
\end{theorem}

\begin{proof}\phantomsection\label{strong:proof:039}
The equivalence between the quotient statement and the scalar compression identity is immediate from the definition. The latter identity is the finite Knill Laflamme criterion \cite{KnillLaflamme1997,SchumacherNielsen1996}.
\phantomsection\label{strong:proof:039:end}\end{proof}

The next result gives the packet construction behind the criterion.

\begin{theorem}[Orthogonal syndrome construction]\label{strong:theorem:027}\label{strong:thm:syndrome_construction}
Assume the Knill Laflamme identities. Then the coefficient matrix \(c\) is positive semidefinite. Choose a unitary matrix \(W\) such that
\[
WcW^*=\diag(\lambda_1,\ldots,\lambda_r),
\qquad
\lambda_j\geq0,
\]
and define
\[
F_j=\sum_a\overline{W_{ja}}E_a.
\]
For every \(j\) with \(\lambda_j>0\), the operator
\[
V_j=\lambda_j^{-1/2}F_jP
\]
is an isometry from the code into \(\Hcal\). Its range projector
\[
Q_j=V_jV_j^*
\]
is a syndrome packet, and the nonzero \(Q_j\) are pairwise orthogonal.
\end{theorem}

\begin{proof}\phantomsection\label{strong:proof:040}
For \(z\in\CC^r\), let \(F_z=\sum_az_aE_a\). Then
\[
PF_z^*F_zP=(z^*cz)P\geq0,
\]
so \(c\geq0\). The unitary change of error coordinates gives
\[
PF_j^*F_\ell P
=
\lambda_j\delta_{j\ell}P.
\]
Thus
\[
V_j^*V_\ell=\delta_{j\ell}P.
\]
Each \(V_j\) is an isometry on \(P\Hcal\), and distinct ranges are orthogonal. Their range projectors are therefore pairwise orthogonal.
\phantomsection\label{strong:proof:040:end}\end{proof}

\begin{theorem}[Exact recovery on the error span]\label{strong:theorem:028}\label{strong:thm:exact_recovery}
Under the hypotheses of \cref{strong:thm:syndrome_construction}, define on the syndrome support
\[
Q=\sum_{\lambda_j>0}Q_j
\]
the completely positive map
\[
\Rcal_0(X)
=
\sum_{\lambda_j>0}
V_j^*Q_jXQ_jV_j.
\]
It extends to a quantum channel \(\Rcal\) on \(\End(\Hcal)\). If
\[
\Phi(X)=\sum_aE_aXE_a^*
\]
is trace preserving on code states, then
\[
\Rcal\Phi(\rho)=\rho
\]
for every density matrix \(\rho=P\rho P\).
\end{theorem}

\begin{proof}\phantomsection\label{strong:proof:041}
The map \(\Rcal_0\) is completely positive. It is trace preserving on operators supported in \(Q\Hcal\) because
\[
\Tr\bigl(V_j^*Q_jXQ_jV_j\bigr)
=
\Tr(Q_jXQ_j).
\]
Fix a code state \(\rho_0\) and write \(Q^\perp=I-Q\). Then
\[
\Rcal(X)
=
\Rcal_0(QXQ)
+
\Tr(Q^\perp X)\rho_0
\]
is completely positive and trace preserving. It extends \(\Rcal_0\) on operators supported in \(Q\Hcal\).

Unitary mixing of a Kraus family does not change the channel, so
\[
\Phi(\rho)
=
\sum_jF_j\rho F_j^*
=
\sum_{\lambda_j>0}
\lambda_jV_j\rho V_j^*.
\]
Recovery gives
\[
\Rcal\Phi(\rho)
=
\left(\sum_j\lambda_j\right)\rho.
\]
Trace preservation on a code state implies \(\sum_j\lambda_j=1\).
\phantomsection\label{strong:proof:041:end}\end{proof}

\begin{corollary}[Syndrome measurement preserves logical information]\label{strong:corollary:011}\label{strong:cor:syndrome_preserves_logic}
For a channel trace preserving on code states as in the preceding theorem, the completed projective measurement consisting of \((Q_j)_{\lambda_j>0}\) and \(I-\sum_{\lambda_j>0}Q_j\) identifies the orthogonal error syndrome without resolving the encoded state inside the code. The complement outcome has probability zero on the corrupted code state.
\end{corollary}

\begin{proof}\phantomsection\label{strong:proof:042}
On syndrome \(j\), the encoded state is carried by the isometry \(V_j\). The outcome probability is \(\lambda_j\), independent of the normalized code state, while the postmeasurement state is \(V_j\rho V_j^*\).
\phantomsection\label{strong:proof:042:end}\end{proof}

\begin{definition}[Approximate Knill Laflamme defect]\label{strong:definition:006}\label{strong:def:approximate_kl}
For a finite error family define, with the maximum of the empty family set to zero,
\[
\varepsilon_{\rm KL}
=
\max_{a,b}
\inf_{c\in\CC}
\norm{PE_a^*E_bP-cP}.
\]
\end{definition}

\begin{proposition}[Meaning of the approximate defect]\label{strong:proposition:003}\label{strong:prop:approximate_kl_meaning}
For each pair \((a,b)\), choose a minimizing scalar \(c_{ab}\). For unit code vectors \(\psi,\phi\),
\[
\abs{
\ip{E_a\psi}{E_b\phi}
-c_{ab}\ip{\psi}{\phi}
}
\leq
\varepsilon_{\rm KL}.
\]
Thus \(\varepsilon_{\rm KL}\) controls both state dependence of error overlaps and nonorthogonality of nominal syndrome images.
\end{proposition}

\begin{proof}\phantomsection\label{strong:proof:043}
The left side is
\[
\abs{
\ip{\psi}{(PE_a^*E_bP-c_{ab}P)\phi}
},
\]
which is bounded by the operator norm.
\phantomsection\label{strong:proof:043:end}\end{proof}

No recovery fidelity estimate is inferred from this norm alone. Approximate error correction requires a separate continuity theorem and a declared channel metric.

\begin{example}[Three qubit repetition code]\label{strong:example:003}\label{strong:ex:bit_flip_code}
Let
\[
\ket{0_L}=\ket{000},
\qquad
\ket{1_L}=\ket{111},
\]
and take the errors
\[
I,
\quad
X\otimes I\otimes I,
\quad
I\otimes X\otimes I,
\quad
I\otimes I\otimes X.
\]
Their images of the code are four pairwise orthogonal two dimensional syndrome packets whose sum is the full eight dimensional Hilbert space.
\end{example}

\begin{figure}[t]
\centering
\includegraphics[width=.68\textwidth]{09_qec_syndromes.pdf}
\caption{The syndrome packets of the three qubit repetition code. Each packet has dimension two and carries the full logical state.}
\label{strong:fig:qec_syndromes}
\end{figure}

\begin{standardfact}[Correctable observable algebras]\label{strong:standardfact:003}\label{strong:fact:operator_qec}
Let \(\Phi\) be a finite completely positive trace preserving channel with Kraus operators \((E_a)\). A finite dimensional unital star subalgebra \(\Mcal\subseteq P\End(\Hcal)P\) is correctable when a completely positive trace preserving recovery \(\Rcal\) satisfies \(P(\Phi^*\circ\Rcal^*)(X)P=X\) for every \(X\in\Mcal\), where stars on maps denote Heisenberg adjoints. This holds exactly when every compressed product \(PE_a^*E_bP\) belongs to the commutant \(\Mcal'\) on the code. The scalar code criterion is the special case \(\Mcal=P\End(\Hcal)P\) \cite{BenyKempfKribs2007}.
\end{standardfact}

\section{Feshbach effective dynamics and hidden packet memory}\label{strong:section:009}\label{strong:sec:feshbach}

Let \(P\) be an orthogonal projector, put \(Q=I-P\), and write a self adjoint Hamiltonian as
\[
H=
\begin{pmatrix}
A&B\\
B^*&D
\end{pmatrix}
\]
relative to \(P\Hcal\oplus Q\Hcal\).

\begin{definition}[Feshbach map and self energy]\label{strong:definition:007}\label{strong:def:feshbach_map}
For \(z\notin\Spec D\), define
\[
\Sigma(z)=B(z-D)^{-1}B^*
\]
and
\[
F_P(z)=z-A-\Sigma(z).
\]
\end{definition}

\begin{theorem}[Exact projected resolvent]\label{strong:theorem:029}\label{strong:thm:feshbach_resolvent}
If \(z\notin\Spec D\) and \(F_P(z)\) is invertible, then
\[
P(z-H)^{-1}P=F_P(z)^{-1}
\]
as an operator on \(P\Hcal\).
\end{theorem}

\begin{proof}\phantomsection\label{strong:proof:044}
The lower right block of \(z-H\) is \(z-D\). Block Gaussian elimination gives
\[
\begin{pmatrix}
I&B(z-D)^{-1}\\
0&I
\end{pmatrix}
(z-H)
\begin{pmatrix}
I&0\\
(z-D)^{-1}B^*&I
\end{pmatrix}
=
\begin{pmatrix}
F_P(z)&0\\
0&z-D
\end{pmatrix}.
\]
Invert the factors and read the upper left block.
\phantomsection\label{strong:proof:044:end}\end{proof}

\begin{theorem}[Spectral equivalence]\label{strong:theorem:030}\label{strong:thm:feshbach_spectrum}
Let \(z\notin\Spec D\). Then
\[
z\in\Spec H
\quad\Longleftrightarrow\quad
0\in\Spec F_P(z).
\]
The map
\[
x\longmapsto
\begin{pmatrix}
x\\
(z-D)^{-1}B^*x
\end{pmatrix}
\]
is an isomorphism from \(\Ker F_P(z)\) onto \(\Ker(z-H)\).
\end{theorem}

\begin{proof}\phantomsection\label{strong:proof:045}
The equation \((z-H)(x,y)^{\mathsf T}=0\) is
\[
(z-A)x-By=0,
\qquad
-B^*x+(z-D)y=0.
\]
The second equation gives \(y=(z-D)^{-1}B^*x\). Substitution into the first gives \(F_P(z)x=0\). The construction is reversible.
\phantomsection\label{strong:proof:045:end}\end{proof}

\begin{figure}[t]
\centering
\includegraphics[width=.78\textwidth]{10_feshbach_spectrum.pdf}
\caption{The zeros of the Feshbach determinant recover the eigenvalues of the full finite Hamiltonian away from the hidden spectrum.}
\label{strong:fig:feshbach_spectrum}
\end{figure}

\begin{theorem}[Upper half plane sign]\label{strong:theorem:031}\label{strong:thm:self_energy_sign}
If \(\operatorname{Im}z>0\), then
\[
\operatorname{Im}\Sigma(z)
=
-(\operatorname{Im}z)
B(z-D)^{-1}(\overline z-D)^{-1}B^*
\leq0.
\]
Thus \(-\Sigma\) is an operator valued Herglotz function.
\end{theorem}

\begin{proof}\phantomsection\label{strong:proof:046}
For \(R(z)=(z-D)^{-1}\), the resolvent identity gives
\[
R(z)-R(z)^*
=(\overline z-z)R(z)R(z)^*.
\]
Divide by \(2\ii\) and sandwich the result between \(B\) and \(B^*\).
\phantomsection\label{strong:proof:046:end}\end{proof}

\begin{proposition}[Uniform self energy bounds]\label{strong:proposition:004}\label{strong:prop:self_energy_bounds}
Let
\[
\delta=\dist(z,\Spec D)>0,
\qquad
g=\norm B.
\]
Then
\[
\norm{\Sigma(z)}\leq\frac{g^2}{\delta}
\]
and
\[
\norm{\Sigma'(z)}\leq\frac{g^2}{\delta^2}.
\]
For real \(E,E'\) in an interval at distance at least \(\delta\) from \(\Spec D\),
\[
\norm{\Sigma(E)-\Sigma(E')}
\leq
\frac{g^2}{\delta^2}\abs{E-E'}.
\]
\end{proposition}

\begin{proof}\phantomsection\label{strong:proof:047}
The self adjoint resolvent satisfies
\[
\norm{(z-D)^{-1}}=\frac1{\dist(z,\Spec D)}.
\]
Differentiate the resolvent to obtain
\[
\Sigma'(z)=-B(z-D)^{-2}B^*.
\]
The first two estimates follow. The last is the mean value estimate along the real interval.
\phantomsection\label{strong:proof:047:end}\end{proof}

\begin{proposition}[Low energy expansion]\label{strong:proposition:005}\label{strong:prop:low_energy_expansion}
Assume \(D\) is invertible and \(\abs z\norm{D^{-1}}<1\). Then
\[
\Sigma(z)
=-
\sum_{n=0}^{\infty}
z^nBD^{-n-1}B^*
\]
with norm convergence. In particular,
\[
F_P(z)
=
 z-A+BD^{-1}B^*+O\bigl(\abs z\norm B^2\norm{D^{-1}}^2\bigr).
\]
\end{proposition}

\begin{proof}\phantomsection\label{strong:proof:048}
Write
\[
(z-D)^{-1}
=-D^{-1}(I-zD^{-1})^{-1}
=-\sum_{n=0}^{\infty}z^nD^{-n-1}
\]
and sandwich by \(B\) and \(B^*\).
\phantomsection\label{strong:proof:048:end}\end{proof}

\begin{proposition}[Bare truncation error identity]\label{strong:proposition:006}\label{strong:prop:bare_truncation}
If both \(F_P(z)\) and \(z-A\) are invertible, then
\[
P(z-H)^{-1}P-(z-A)^{-1}
=
F_P(z)^{-1}\Sigma(z)(z-A)^{-1}.
\]
\end{proposition}

\begin{proof}\phantomsection\label{strong:proof:049}
Use the inverse difference identity
\[
X^{-1}-Y^{-1}=X^{-1}(Y-X)Y^{-1}
\]
with \(X=F_P(z)\) and \(Y=z-A\).
\phantomsection\label{strong:proof:049:end}\end{proof}

\subsection{Time memory}\label{strong:subsec:time_memory}

Write a Schrödinger state as \(\psi=x+y\), where \(x=P\psi\) and \(y=Q\psi\).

\begin{theorem}[Exact retained memory equation]\label{strong:theorem:032}\label{strong:thm:memory_equation}
The retained coordinate satisfies
\[
\begin{aligned}
\ii\dot x(t)
={}&Ax(t)+B\e^{-\ii tD}y(0)\\
&-\ii\int_0^tK(t-s)x(s)\dd s,
\end{aligned}
\]
where
\[
K(t)=B\e^{-\ii tD}B^*.
\]
For \(\operatorname{Re}s>0\),
\[
\int_0^\infty\e^{-st}K(t)\dd t
=
B(s+\ii D)^{-1}B^*.
\]
\end{theorem}

\begin{proof}\phantomsection\label{strong:proof:050}
The hidden equation is
\[
\ii\dot y=B^*x+Dy.
\]
Variation of constants gives
\[
y(t)=\e^{-\ii tD}y(0)
-\ii\int_0^t\e^{-\ii(t-s)D}B^*x(s)\dd s.
\]
Substitute into \(\ii\dot x=Ax+By\). The Laplace formula follows by diagonalizing \(D\) or integrating the norm convergent finite dimensional exponential.
\phantomsection\label{strong:proof:050:end}\end{proof}

\begin{theorem}[Positive definite quantum memory]\label{strong:theorem:033}\label{strong:thm:positive_memory}
The kernel \(K(t-s)\) is positive definite. For times \(t_1,\ldots,t_m\) and retained vectors \(v_1,\ldots,v_m\),
\[
\sum_{j,k=1}^{m}
\ip{v_j}{K(t_j-t_k)v_k}
\geq0.
\]
\end{theorem}

\begin{proof}\phantomsection\label{strong:proof:051}
Set
\[
w=\sum_{k=1}^{m}\e^{\ii t_kD}B^*v_k.
\]
Then
\[
\norm w^2
=
\sum_{j,k}
\ip{B^*v_j}{\e^{-\ii(t_j-t_k)D}B^*v_k}
=
\sum_{j,k}
\ip{v_j}{K(t_j-t_k)v_k}.
\]
\phantomsection\label{strong:proof:051:end}\end{proof}

\begin{corollary}[Finite oscillatory memory expansion]\label{strong:corollary:012}\label{strong:cor:oscillatory_memory}
If
\[
D=\sum_{r}E_rQ_r
\]
is the spectral decomposition of the hidden Hamiltonian, then
\[
K(t)=\sum_r\e^{-\ii tE_r}BQ_rB^*.
\]
Every coefficient \(BQ_rB^*\) is positive semidefinite.
\end{corollary}

\begin{proof}\phantomsection\label{strong:proof:052}
Insert the spectral decomposition into the exponential.
\phantomsection\label{strong:proof:052:end}\end{proof}

\begin{figure}[t]
\centering
\includegraphics[width=.76\textwidth]{11_memory_kernel.pdf}
\caption{A finite hidden spectrum produces an oscillatory memory kernel rather than an instantaneous retained Hamiltonian.}
\label{strong:fig:memory_kernel}
\end{figure}

\begin{theorem}[Exact local closure criterion]\label{strong:theorem:034}\label{strong:thm:local_closure}
Assume \(y(0)=0\). The retained equation is local and equal to
\[
\ii\dot x=Ax
\]
for every retained initial state exactly when \(B=0\). Equivalently, exact local closure holds exactly when the retained and hidden packets reduce \(H\).
\end{theorem}

\begin{proof}\phantomsection\label{strong:proof:053}
If \(B=0\), the retained and hidden equations decouple. Conversely, suppose that every solution with \(y(0)=0\) satisfies \(x(t)=\e^{-\ii tA}x(0)\). The full block equation gives
\[
\ddot x(0)
=
-\bigl(A^2+BB^*\bigr)x(0),
\]
while the declared local equation gives
\[
\ddot x(0)=-A^2x(0).
\]
Since this holds for every retained initial vector, \(BB^*=0\). Hence \(B=0\).
\phantomsection\label{strong:proof:053:end}\end{proof}

\subsection{Stability across models}\label{strong:subsec:feshbach_stability}

\begin{theorem}[Stability of the self energy]\label{strong:theorem:035}\label{strong:thm:self_energy_stability}
Let \(D=D^*\) and \(\widetilde D=\widetilde D^*\) act on one fixed hidden Hilbert space, and let \(B,\widetilde B\) map that hidden space to one fixed retained Hilbert space. Define
\[
\Sigma(z)=B(z-D)^{-1}B^*,
\qquad
\widetilde\Sigma(z)=\widetilde B(z-\widetilde D)^{-1}\widetilde B^*.
\]
Assume
\[
\norm B,\norm{\widetilde B}\leq g
\]
and
\[
\dist(z,\Spec D),
\dist(z,\Spec\widetilde D)
\geq\delta>0.
\]
Then
\[
\norm{\Sigma(z)-\widetilde\Sigma(z)}
\leq
\frac{2g}{\delta}\norm{B-\widetilde B}
+
\frac{g^2}{\delta^2}\norm{D-\widetilde D}.
\]
\end{theorem}

\begin{proof}\phantomsection\label{strong:proof:054}
Add and subtract \(\widetilde B(z-D)^{-1}B^*\) and \(\widetilde B(z-D)^{-1}\widetilde B^*\). The first two differences give the coupling term. For the last, use
\[
(z-D)^{-1}-(z-\widetilde D)^{-1}
=(z-D)^{-1}(D-\widetilde D)(z-\widetilde D)^{-1}.
\]
Apply the resolvent bounds.
\phantomsection\label{strong:proof:054:end}\end{proof}

\begin{corollary}[Uniform Feshbach control]\label{strong:corollary:013}\label{strong:cor:uniform_feshbach}
For a family represented on fixed retained and hidden Hilbert spaces, with self adjoint hidden Hamiltonians, a common hidden spectral separation \(\delta\), and a common coupling bound \(g\), the self energies are uniformly bounded and uniformly Lipschitz on the declared spectral region. If the couplings and hidden Hamiltonians converge in norm, then the self energies converge uniformly there.
\end{corollary}

\begin{proof}\phantomsection\label{strong:proof:055}
Apply \cref{strong:prop:self_energy_bounds,strong:thm:self_energy_stability} uniformly.
\phantomsection\label{strong:proof:055:end}\end{proof}

This result identifies the finite hypotheses needed for cutoff independent Feshbach control. It does not prove that a chosen geometric refinement tower satisfies them.

\section{Euclidean transfer and finite reflection positivity}\label{strong:section:010}\label{strong:sec:euclidean}

Let \(H=H^*\) on a nonzero finite dimensional Hilbert space and let
\[
E_0=\min\Spec H.
\]
For a lattice time step \(a>0\), define
\[
T=\e^{-a(H-E_0I)}.
\]

\begin{theorem}[Finite Euclidean transfer]\label{strong:theorem:036}\label{strong:thm:euclidean_transfer}
The operator \(T\) is positive definite, self adjoint, and contractive. Its spectrum lies in \((0,1]\), and
\[
T^n=\e^{-na(H-E_0I)}
\]
for every \(n\in\NN\).
\end{theorem}

\begin{proof}\phantomsection\label{strong:proof:056}
Diagonalize \(H\). Each eigenvalue of \(T\) is
\[
\e^{-a(E_j-E_0)}
\]
and therefore lies in \((0,1]\). The power law follows from the functional calculus.
\phantomsection\label{strong:proof:056:end}\end{proof}

\begin{theorem}[Hamiltonian reconstruction from transfer]\label{strong:theorem:037}\label{strong:thm:transfer_logarithm}
Let \(T\) be positive definite, self adjoint, and contractive. Then
\[
H_T=-a^{-1}\log T
\]
is a nonnegative self adjoint operator and
\[
T^n=\e^{-naH_T}.
\]
It is the unique nonnegative self adjoint operator with \(T=\e^{-aH_T}\).
\end{theorem}

\begin{proof}\phantomsection\label{strong:proof:057}
Apply the real logarithm to the positive spectrum of \(T\). The spectral calculus gives existence and the power identity. Injectivity of the scalar exponential on \([0,\infty)\) gives uniqueness.
\phantomsection\label{strong:proof:057:end}\end{proof}

\begin{theorem}[Finite reflection positivity]\label{strong:theorem:038}\label{strong:thm:finite_reflection_positivity}
For vectors \(v_1,\ldots,v_m\in\Hcal\) and nonnegative integers \(n_1,\ldots,n_m\), the matrix
\[
G_{jk}=\ip{v_j}{T^{n_j+n_k}v_k}
\]
is positive semidefinite. Equivalently,
\[
\sum_{j,k}\overline{c_j}c_k
\ip{v_j}{T^{n_j+n_k}v_k}
\geq0
\]
for every \(c\in\CC^m\).
\end{theorem}

\begin{proof}\phantomsection\label{strong:proof:058}
Since \(T\) is positive and self adjoint,
\[
\ip{v_j}{T^{n_j+n_k}v_k}
=
\ip{T^{n_j}v_j}{T^{n_k}v_k}.
\]
The quadratic form is the squared norm of
\[
\sum_kc_kT^{n_k}v_k.
\]
\phantomsection\label{strong:proof:058:end}\end{proof}

\begin{corollary}[Reflection positive correlators]\label{strong:corollary:014}\label{strong:cor:reflection_correlators}
Fix a vector \(\Omega\) and operators \(A_1,\ldots,A_m\). The matrix
\[
C_{jk}
=
\ip{A_j\Omega}{T^{n_j+n_k}A_k\Omega}
\]
is positive semidefinite.
\end{corollary}

\begin{proof}\phantomsection\label{strong:proof:059}
Apply \cref{strong:thm:finite_reflection_positivity} with \(v_j=A_j\Omega\).
\phantomsection\label{strong:proof:059:end}\end{proof}

\begin{figure}[t]
\centering
\includegraphics[width=.68\textwidth]{12_reflection_positivity.pdf}
\caption{Eigenvalues of a finite reflection positive Gram matrix. Small numerical negative values would signal either roundoff or failure of the transfer assumptions.}
\label{strong:fig:reflection_positivity}
\end{figure}

\begin{proposition}[Ground packet convergence]\label{strong:proposition:007}\label{strong:prop:ground_packet}
Let \(P_0\) be the spectral projector of \(H\) at \(E_0\), and let
\[
\Delta=\min\{E_j-E_0:E_j>E_0\}.
\]
If an excited eigenvalue exists, then
\[
\norm{T^n-P_0}=\e^{-na\Delta}.
\]
\end{proposition}

\begin{proof}\phantomsection\label{strong:proof:060}
In the spectral decomposition, \(T^n-P_0\) has eigenvalues zero on the ground space and \(\e^{-na(E_j-E_0)}\) on the excited eigenspaces. The largest is \(\e^{-na\Delta}\).
\phantomsection\label{strong:proof:060:end}\end{proof}

\begin{proposition}[Symmetry resolved transfer]\label{strong:proposition:008}\label{strong:prop:symmetry_transfer}
If \([H,P_a]=0\) for a complete packet family, then
\[
T=\sum_aP_aTP_a
\]
and each reflection positive Gram matrix decomposes as the sum of its packet Gram matrices.
\end{proposition}

\begin{proof}\phantomsection\label{strong:proof:061}
The functional calculus preserves commutation. Decompose every vector into packet components and use orthogonality.
\phantomsection\label{strong:proof:061:end}\end{proof}

Finite reflection positivity is a necessary structural property of transfer data. It is not by itself the continuum Osterwalder Schrader theorem, which also requires a suitable algebra of Euclidean fields, covariance, regularity, clustering, and a controlled continuum limit \cite{OsterwalderSchrader1973,OsterwalderSchrader1975,GlimmJaffe1987}.

\section{Finite spectral geometry and diffusion dimension}\label{strong:section:011}\label{strong:sec:spectral_geometry}

A finite spectral datum consists of a finite dimensional Hilbert space, a nonnegative self adjoint operator \(L\), and any declared packet family. One may take \(L=D^2\) for a finite Dirac type operator or a graph Laplacian. Related geometric viewpoints appear in \cite{Connes1994,Vassilevich2003}.

\begin{definition}[Heat trace and spectral dimension]\label{strong:definition:008}\label{strong:def:spectral_dimension}
Let the eigenvalues of \(L\) be \(\lambda_1,\ldots,\lambda_N\), with at least one zero eigenvalue. Define
\[
Z_L(t)=\Tr(\e^{-tL})
\]
and
\[
d_L(t)
=-2\frac{\dd\log Z_L(t)}{\dd\log t}
=
2t
\frac{\Tr(L\e^{-tL})}{\Tr(\e^{-tL})}.
\]
\end{definition}

\begin{proposition}[Finite scale limits]\label{strong:proposition:009}\label{strong:prop:finite_spectral_limits}
For every finite nonnegative \(L\) with nontrivial kernel,
\[
\lim_{t\downarrow0}d_L(t)=0
\]
and
\[
\lim_{t\to\infty}d_L(t)=0.
\]
For every \(t>0\),
\[
0\leq d_L(t)\leq2t\norm L.
\]
\end{proposition}

\begin{proof}\phantomsection\label{strong:proof:062}
The formula in \cref{strong:def:spectral_dimension} is twice \(t\) times a weighted average of the eigenvalues. This gives the bound and the limit at zero. At large time, the denominator converges to \(\dim\Ker L\), while every positive eigenvalue contribution in the numerator decays exponentially.
\phantomsection\label{strong:proof:062:end}\end{proof}

Thus a finite system cannot have a nonzero spectral dimension at the extreme ultraviolet or infrared ends. A geometric dimension can appear only in a refinement window.

\begin{proposition}[Heat packet decomposition]\label{strong:proposition:010}\label{strong:prop:heat_packets}
If \([L,P_a]=0\) for every packet, then
\[
Z_L(t)=\sum_aZ_{L,a}(t),
\qquad
Z_{L,a}(t)=\Tr(P_a\e^{-tL}).
\]
The global spectral dimension is the heat weight average
\[
d_L(t)
=
\sum_{a:Z_{L,a}(t)>0}
\frac{Z_{L,a}(t)}{Z_L(t)}
 d_{L,a}(t),
\]
where
\[
d_{L,a}(t)
=2t
\frac{\Tr(P_aL\e^{-tL})}{Z_{L,a}(t)}
\]
for nonempty packet heat traces.
\end{proposition}

\begin{proof}\phantomsection\label{strong:proof:063}
Commutation gives a block diagonal heat operator. Add the traces and the corresponding first spectral moments, then divide by the total heat trace.
\phantomsection\label{strong:proof:063:end}\end{proof}

\subsection{Periodic lattices}\label{strong:subsec:periodic_lattices}

Let \(N\geq2\) and let \(\Lambda_N=(\ZZ/N\ZZ)^d\) be the unit flat torus grid with spacing \(1/N\). The standard periodic difference Laplacian has eigenvalues
\[
\lambda_N(k)
=
4N^2
\sum_{j=1}^{d}
\sin^2\left(\frac{\pi k_j}{N}\right),
\]
where \(k\) runs through centered residue representatives. This is the standard periodic finite difference Laplacian \cite{Strikwerda2004}.

\begin{theorem}[Low mode dispersion in dimension \(d\)]\label{strong:theorem:039}\label{strong:thm:dispersion_bound}
For every centered mode \(k\),
\[
0
\leq
4\pi^2\abs{k}^2-\lambda_N(k)
\leq
\frac{4\pi^4}{3N^2}
\sum_{j=1}^{d}k_j^4.
\]
\end{theorem}

\begin{proof}\phantomsection\label{strong:proof:064}
For \(x\in[0,\pi/2]\),
\[
\sin x\leq x
\]
and
\[
\sin x\geq x-\frac{x^3}{6}.
\]
Thus
\[
0\leq x^2-\sin^2x\leq\frac{x^4}{3}.
\]
Apply this with \(x=\pi\abs{k_j}/N\), multiply by \(4N^2\), and sum over \(j\).
\phantomsection\label{strong:proof:064:end}\end{proof}

\begin{theorem}[Fixed time heat trace convergence]\label{strong:theorem:040}\label{strong:thm:heat_trace_convergence}
For every \(t>0\),
\[
\sum_{k\in\Lambda_N^*}
\e^{-t\lambda_N(k)}
\longrightarrow
\sum_{k\in\ZZ^d}
\e^{-4\pi^2t\abs{k}^2}
\]
as \(N\to\infty\). The first spectral moments also converge:
\[
\sum_{k\in\Lambda_N^*}
\lambda_N(k)\e^{-t\lambda_N(k)}
\longrightarrow
\sum_{k\in\ZZ^d}
4\pi^2\abs{k}^2\e^{-4\pi^2t\abs{k}^2}.
\]
Consequently, the discrete spectral dimensions converge to the continuum torus spectral dimension at every fixed positive time.
\end{theorem}

\begin{proof}\phantomsection\label{strong:proof:065}
For a fixed integer mode, \cref{strong:thm:dispersion_bound} gives
\[
\lambda_N(k)\longrightarrow4\pi^2\abs{k}^2.
\]
For centered modes, the inequality
\[
\sin x\geq\frac{2x}{\pi},
\qquad
0\leq x\leq\frac\pi2,
\]
gives
\[
\lambda_N(k)\geq16\abs{k}^2.
\]
Also \(\lambda_N(k)\leq4\pi^2\abs{k}^2\). Extend the discrete summands by zero outside the centered residue box. The heat trace terms are dominated by
\[
\e^{-16t\abs{k}^2},
\]
and the moment terms by
\[
4\pi^2\abs{k}^2\e^{-16t\abs{k}^2}.
\]
Both dominators are summable on \(\ZZ^d\). Dominated convergence proves the two limits. Their quotient gives spectral dimension convergence.
\phantomsection\label{strong:proof:065:end}\end{proof}

\begin{standardfact}[Theta transformation]\label{strong:standardfact:004}\label{strong:fact:theta_transform}
The continuum flat torus heat trace satisfies
\[
Z_d(t)
=
\sum_{k\in\ZZ^d}\e^{-4\pi^2t\abs{k}^2}
=
(4\pi t)^{-d/2}
\sum_{n\in\ZZ^d}\e^{-\abs{n}^2/(4t)}.
\]
This is the Poisson summation formula for a Gaussian.
\end{standardfact}

\begin{theorem}[Continuum torus small time dimension]\label{strong:theorem:041}\label{strong:thm:continuum_dimension}
Let
\[
d_{\rm torus}(t)
=-2\frac{\dd\log Z_d(t)}{\dd\log t}.
\]
Then
\[
d_{\rm torus}(t)
=
 d+O\left(t^{-1}\e^{-1/(4t)}\right)
\]
as \(t\downarrow0\).
\end{theorem}

\begin{proof}\phantomsection\label{strong:proof:066}
By \cref{strong:fact:theta_transform},
\[
Z_d(t)
=(4\pi t)^{-d/2}(1+R(t)),
\]
where
\[
R(t)=\sum_{n\neq0}\e^{-\abs{n}^2/(4t)}
=O\bigl(\e^{-1/(4t)}\bigr).
\]
Differentiation gives
\[
tR'(t)=O\left(t^{-1}\e^{-1/(4t)}\right).
\]
Insert this in the logarithmic derivative.
\phantomsection\label{strong:proof:066:end}\end{proof}

\begin{corollary}[Rigorous intermediate dimension window]\label{strong:corollary:015}\label{strong:cor:dimension_window}
For every \(\varepsilon>0\), there is a diffusion time \(t_\varepsilon>0\) and an integer \(N_\varepsilon\) such that
\[
\abs{d_{L_N}(t_\varepsilon)-d}<\varepsilon
\]
for every \(N\geq N_\varepsilon\).
\end{corollary}

\begin{proof}\phantomsection\label{strong:proof:067}
Choose \(t_\varepsilon\) so that \cref{strong:thm:continuum_dimension} is within \(\varepsilon/2\) of \(d\). Then apply \cref{strong:thm:heat_trace_convergence} to choose \(N_\varepsilon\) so that the discrete spectral dimension is within \(\varepsilon/2\) of the continuum value.
\phantomsection\label{strong:proof:067:end}\end{proof}

\begin{figure}[t]
\centering
\includegraphics[width=.78\textwidth]{13_spectral_dimension.pdf}
\caption{Finite periodic lattices show an intermediate dimension window. The ultraviolet and infrared limits return to zero because the state space is finite.}
\label{strong:fig:spectral_dimension}
\end{figure}

The order of limits matters. Taking \(t\downarrow0\) at fixed \(N\) gives zero. Taking \(N\to\infty\) at fixed positive \(t\), then taking \(t\downarrow0\), gives the continuum dimension.

\section{Constraint closure and anomaly packets}\label{strong:section:012}\label{strong:sec:constraints}

Let \(C_1,\ldots,C_m\) be self adjoint constraint operators on a finite dimensional Hilbert space. The distinction between constraint action and constraint closure is standard in constrained Hamiltonian theory \cite{Dirac1964,HenneauxTeitelboim1992}. Write
\[
\mathfrak c=\operatorname{span}_{\CC}\{C_1,\ldots,C_m\}
\subseteq\End(\Hcal).
\]

\begin{definition}[Universal closure obstruction]\label{strong:definition:009}\label{strong:def:closure_obstruction}
Let
\[
q_{\mathfrak c}:\End(\Hcal)\longrightarrow\End(\Hcal)/\mathfrak c
\]
be the quotient map. Define
\[
\Ob_{ab}^{\rm cl}
=
q_{\mathfrak c}([C_a,C_b]).
\]
\end{definition}

\begin{theorem}[Finite constraint closure criterion]\label{strong:theorem:042}\label{strong:thm:constraint_closure}
The constraint span \(\mathfrak c\) is closed under commutators exactly when
\[
\Ob_{ab}^{\rm cl}=0
\]
for all \(a,b\).
\end{theorem}

\begin{proof}\phantomsection\label{strong:proof:068}
The quotient class vanishes exactly when \([C_a,C_b]\in\mathfrak c\). This holds for every pair exactly when \(\mathfrak c\) is a Lie algebra under the commutator.
\phantomsection\label{strong:proof:068:end}\end{proof}

If structure coefficients \(f_{ab}^{\phantom{ab}c}\) are prescribed, define the fixed structure anomaly
\[
\mathfrak A_{ab}
=
[C_a,C_b]
-\ii\sum_cf_{ab}^{\phantom{ab}c}C_c.
\]
Exact realization of the prescribed algebra is the stronger condition \(\mathfrak A_{ab}=0\).

\begin{proposition}[Physical compression is not a closure test]\label{strong:proposition:011}\label{strong:prop:physical_compression_trivial}
Let
\[
\Hcal_{\phys}=\bigcap_a\Ker C_a
\]
and let \(P_{\phys}\) be its orthogonal projector. Then
\[
P_{\phys}[C_a,C_b]P_{\phys}=0
\]
for every pair \(a,b\), whether or not the constraint span closes.
\end{proposition}

\begin{proof}\phantomsection\label{strong:proof:069}
Self adjointness and \(C_aP_{\phys}=0\) give \(P_{\phys}C_a=0\). Both terms in the compressed commutator therefore vanish.
\phantomsection\label{strong:proof:069:end}\end{proof}

This is why the quotient in \cref{strong:def:closure_obstruction} is needed. Compression to physical states tests action on the common kernel, not closure of the operator algebra.

\begin{figure}[t]
\centering
\includegraphics[width=.88\textwidth]{15_constraint_obstruction.pdf}
\caption{The closure obstruction lives in a quotient by the constraint span. Physical compression is automatic for self adjoint constraints and cannot replace this quotient.}
\label{strong:fig:constraint_obstruction}
\end{figure}

\begin{theorem}[Invariance of the physical space]\label{strong:theorem:043}\label{strong:thm:physical_invariance}
Let \(H=H^*\). Suppose
\[
[H,C_a]
=
\sum_cb_a^{\phantom a c}C_c
\]
for every \(a\). Then \(H\) preserves \(\Hcal_{\phys}\), and so does \(\e^{-\ii tH}\) for every real \(t\).
\end{theorem}

\begin{proof}\phantomsection\label{strong:proof:070}
For \(\psi\in\Hcal_{\phys}\),
\[
C_aH\psi
=HC_a\psi-[H,C_a]\psi
=0.
\]
Thus \(H\psi\in\Hcal_{\phys}\). The exponential preserves every invariant finite dimensional subspace.
\phantomsection\label{strong:proof:070:end}\end{proof}

\subsection{Constraint refinement}\label{strong:subsec:constraint_refinement}

Let \((\Hcal_\ell,C_{\ell,a})\) and \((\Hcal_{\ell+1},C_{\ell+1,a})\) be two levels, and let
\[
R_\ell:\Hcal_\ell\longrightarrow\Hcal_{\ell+1}
\]
be an isometry.

\begin{theorem}[Exact constraint transport]\label{strong:theorem:044}\label{strong:thm:constraint_transport}
Assume
\[
C_{\ell+1,a}R_\ell=R_\ell C_{\ell,a}
\]
for every \(a\). Then
\[
[C_{\ell+1,a},C_{\ell+1,b}]R_\ell
=
R_\ell[C_{\ell,a},C_{\ell,b}].
\]
If the same structure coefficients are used at both levels, then
\[
\mathfrak A_{\ell+1,ab}R_\ell
=
R_\ell\mathfrak A_{\ell,ab}.
\]
The isometry maps coarse physical states into fine physical states.
\end{theorem}

\begin{proof}\phantomsection\label{strong:proof:071}
Apply the intertwining identity twice to each product in the commutator. The anomaly identity follows by subtracting the same linear combination of constraints. If \(C_{\ell,a}\psi=0\), then
\[
C_{\ell+1,a}R_\ell\psi=R_\ell C_{\ell,a}\psi=0.
\]
\phantomsection\label{strong:proof:071:end}\end{proof}

\begin{proposition}[Approximate constraint transport]\label{strong:proposition:012}\label{strong:prop:approx_constraint_transport}
Assume
\[
\norm{C_{\ell+1,a}R_\ell-R_\ell C_{\ell,a}}
\leq\varepsilon_a
\]
and
\[
\norm{C_{\ell,a}},\norm{C_{\ell+1,a}}\leq M.
\]
Then a coarse physical unit vector \(\psi\) satisfies
\[
\norm{C_{\ell+1,a}R_\ell\psi}\leq\varepsilon_a.
\]
Also,
\[
\begin{aligned}
&\norm{
[C_{\ell+1,a},C_{\ell+1,b}]R_\ell
-R_\ell[C_{\ell,a},C_{\ell,b}]
}\\
&\hspace{4em}\leq
2M(\varepsilon_a+\varepsilon_b).
\end{aligned}
\]
\end{proposition}

\begin{proof}\phantomsection\label{strong:proof:072}
The first claim follows by applying the intertwining defect to \(\psi\). Write
\[
E_a=C_{\ell+1,a}R_\ell-R_\ell C_{\ell,a}.
\]
Expanding the two commutators gives four error terms:
\[
E_aC_{\ell,b},
\quad
C_{\ell+1,a}E_b,
\quad
-E_bC_{\ell,a},
\quad
-C_{\ell+1,b}E_a.
\]
Their norms give the bound.
\phantomsection\label{strong:proof:072:end}\end{proof}

Approximate transport controls finite violations. It does not turn a nonzero anomaly into an exact continuum algebra without a limit theorem and a common operator domain.

\section{Quantum refinement towers}\label{strong:section:013}\label{strong:sec:towers}

\begin{definition}[Finite quantum refinement tower]\label{strong:definition:010}\label{strong:def:quantum_tower}
A finite quantum refinement tower consists of Hilbert spaces \(\Hcal_\ell\), isometries
\[
R_\ell:\Hcal_\ell\longrightarrow\Hcal_{\ell+1},
\]
self adjoint Hamiltonians \(H_\ell\), packet projectors \(P_{\ell,a}\), constraints \(C_{\ell,r}\), and positive transfer operators \(T_\ell\). Exact compatibility means that every declared object intertwines with \(R_\ell\). Approximate compatibility is recorded by the corresponding operator norm defects.
\end{definition}

\begin{figure}[t]
\centering
\includegraphics[width=.82\textwidth]{14_continuum_tower.pdf}
\caption{A refinement tower carries Hamiltonians, packets, constraints, and transfer data. Each family has its own compatibility defect.}
\label{strong:fig:continuum_tower}
\end{figure}

\subsection{Exact bounded direct limits}\label{strong:subsec:exact_direct_limit}

Let \(\Hcal_\infty\) be the Hilbert direct limit of the isometric system, and let
\[
\iota_\ell:\Hcal_\ell\longrightarrow\Hcal_\infty
\]
be the canonical isometries.

\begin{theorem}[Bounded Hamiltonian direct limit]\label{strong:theorem:045}\label{strong:thm:bounded_direct_limit}
Assume
\[
H_{\ell+1}R_\ell=R_\ell H_\ell
\]
and
\[
\sup_\ell\norm{H_\ell}<\infty.
\]
Then there is a unique bounded self adjoint operator \(H_\infty\) on \(\Hcal_\infty\) such that
\[
H_\infty\iota_\ell=\iota_\ell H_\ell
\]
for every \(\ell\). It satisfies
\[
\e^{-\ii tH_\infty}\iota_\ell
=
\iota_\ell\e^{-\ii tH_\ell}
\]
and
\[
\e^{-tH_\infty}\iota_\ell
=
\iota_\ell\e^{-tH_\ell}
\]
when the Hamiltonians are nonnegative.
\end{theorem}

\begin{proof}\phantomsection\label{strong:proof:073}
Define \(H_\infty\) on the algebraic union of the embedded levels by
\[
H_\infty\iota_\ell x=\iota_\ell H_\ell x.
\]
Compatibility makes this well defined. The uniform norm bound makes it bounded on the dense union, so it extends uniquely. It is symmetric because every \(H_\ell\) is self adjoint. A bounded symmetric operator defined on the whole Hilbert space is self adjoint. The functional calculus intertwining follows first for polynomials and then for exponentials by uniform convergence.
\phantomsection\label{strong:proof:073:end}\end{proof}

\begin{theorem}[Exact packet direct limit]\label{strong:theorem:046}\label{strong:thm:packet_direct_limit}
Assume the packet families have one fixed finite address set at every level and each packet address \(a\) is compatible in the sense
\[
P_{\ell+1,a}R_\ell=R_\ell P_{\ell,a}.
\]
Then there is an orthogonal projector \(P_{\infty,a}\) on \(\Hcal_\infty\) satisfying
\[
P_{\infty,a}\iota_\ell=\iota_\ell P_{\ell,a}.
\]
Finite completeness and orthogonality pass to the direct limit.
\end{theorem}

\begin{proof}\phantomsection\label{strong:proof:074}
Use the same construction as in \cref{strong:thm:bounded_direct_limit}. Every projector has norm at most one, so no additional boundedness assumption is needed. The projector identities hold on the dense algebraic union and hence on the completion.
\phantomsection\label{strong:proof:074:end}\end{proof}

\subsection{Norm summable approximation}\label{strong:subsec:norm_summable}

The next theorem uses a common Hilbert space. It is appropriate for bounded effective Hamiltonians, not for an unbounded field Hamiltonian.

\begin{theorem}[Norm summable Hamiltonian tower]\label{strong:theorem:047}\label{strong:thm:norm_summable_tower}
Let \((H_n)\) be bounded self adjoint operators on one Hilbert space and assume
\[
\sum_{n=0}^{\infty}\norm{H_{n+1}-H_n}<\infty.
\]
Then \(H_n\) converges in norm to a bounded self adjoint operator \(H\). For every \(T>0\),
\[
\sup_{\abs t\leq T}
\norm{\e^{-\ii tH_n}-\e^{-\ii tH}}
\leq
T\norm{H_n-H}.
\]
If every \(H_n\geq0\), then
\[
\sup_{0\leq t\leq T}
\norm{\e^{-tH_n}-\e^{-tH}}
\leq
T\norm{H_n-H}.
\]
\end{theorem}

\begin{proof}\phantomsection\label{strong:proof:075}
The summability condition makes \((H_n)\) Cauchy in operator norm. The limit is self adjoint. Duhamel's identity gives
\[
\e^{-\ii tH_n}-\e^{-\ii tH}
=-\ii\int_0^t
\e^{-\ii(t-s)H_n}(H_n-H)\e^{-\ii sH}\dd s.
\]
Both unitary factors have norm one. The Euclidean estimate is identical because the nonnegative semigroups are contractions.
\phantomsection\label{strong:proof:075:end}\end{proof}

\begin{corollary}[Limit of packet leakage]\label{strong:corollary:016}\label{strong:cor:limit_leakage}
Let \(P_n\) be orthogonal projectors and \(H_n\) bounded self adjoint operators on one Hilbert space, with \(P_n\to P\) and \(H_n\to H\) in norm. Then for every bounded time interval,
\[
(I-P_n)\e^{-\ii tH_n}P_n
\longrightarrow
(I-P)\e^{-\ii tH}P
\]
uniformly in time.
\end{corollary}

\begin{proof}\phantomsection\label{strong:proof:076}
Add and subtract terms so that one factor changes at a time. The projector terms converge uniformly because \(P_n\to P\) in norm. Duhamel's identity gives
\[
\sup_{\abs t\leq T}
\norm{\e^{-\ii tH_n}-\e^{-\ii tH}}
\leq
T\norm{H_n-H},
\]
which proves uniform convergence of the middle factor.
\phantomsection\label{strong:proof:076:end}\end{proof}

\begin{theorem}[Reflection positivity is closed under limits]\label{strong:theorem:048}\label{strong:thm:reflection_limit}
Suppose \(C_n(j,k)\) are reflection positive kernels on \(\NN\), meaning every finite matrix
\[
\bigl(C_n(j_r,j_s)\bigr)_{r,s}
\]
is positive semidefinite. If
\[
C_n(j,k)\longrightarrow C(j,k)
\]
for every pair, then \(C\) is reflection positive.
\end{theorem}

\begin{proof}\phantomsection\label{strong:proof:077}
For a fixed finite set of times and coefficients, the corresponding quadratic forms are nonnegative and converge to the quadratic form of \(C\). The limit is nonnegative.
\phantomsection\label{strong:proof:077:end}\end{proof}

\begin{theorem}[Constraint closure under norm convergence]\label{strong:theorem:049}\label{strong:thm:constraint_limit}
Let \(C_{n,a}\to C_a\) in norm for every \(a\), and fix structure coefficients \(f_{ab}^{\phantom{ab}c}\). If
\[
\norm{
[C_{n,a},C_{n,b}]
-\ii\sum_cf_{ab}^{\phantom{ab}c}C_{n,c}
}
\longrightarrow0,
\]
then
\[
[C_a,C_b]
=\ii\sum_cf_{ab}^{\phantom{ab}c}C_c.
\]
\end{theorem}

\begin{proof}\phantomsection\label{strong:proof:078}
Operator multiplication and the commutator are continuous in operator norm. Pass to the limit in the anomaly identity.
\phantomsection\label{strong:proof:078:end}\end{proof}

\begin{theorem}[Heat trace convergence from trace norm]\label{strong:theorem:050}\label{strong:thm:trace_norm_heat}
Let \(L_n,L\) be nonnegative self adjoint operators on one Hilbert space. For a fixed \(t>0\), assume their heat operators are trace class and that
\[
\e^{-tL_n}\longrightarrow\e^{-tL}
\]
in trace norm. Then
\[
\Tr(\e^{-tL_n})\longrightarrow\Tr(\e^{-tL}).
\]
If also the heat moment operators are trace class and
\[
L_n\e^{-tL_n}\longrightarrow L\e^{-tL}
\]
in trace norm and the limiting heat trace is nonzero, then the spectral dimensions converge at that time.
\end{theorem}

\begin{proof}\phantomsection\label{strong:proof:079}
The trace is continuous in trace norm. Apply this to the heat operators and to the first heat moments, then take the quotient in \cref{strong:def:spectral_dimension}.
\phantomsection\label{strong:proof:079:end}\end{proof}

\begin{remark}\label{strong:remark:003}\label{strong:rem:unbounded_limit}
The theorems above produce bounded continuum Hamiltonians. A physically useful continuum Hamiltonian is often unbounded. Its construction generally requires strong resolvent convergence, convergence of closed quadratic forms, or another domain sensitive method \cite{Kato1995,ReedSimon1972}. None of those conclusions follows from finite matrix convergence without additional hypotheses.
\end{remark}

\section{Continuum milestones and exact boundary}\label{strong:section:014}\label{strong:sec:boundary}

The finite theorems establish exact identities and bounded limit criteria. They do not collapse the steps between a finite transfer model and a continuum theory. This section records the separate mathematical tasks.

\begin{definition}[Anomaly free continuum constraint system]\label{strong:definition:011}\label{strong:def:anomaly_free_continuum}
A continuum constraint system consists of a Hilbert space, a common invariant dense domain \(\Dcal\), and operators \(C_a\) on \(\Dcal\). It is anomaly free for prescribed structure coefficients when
\[
[C_a,C_b]\psi
=
\ii\sum_cf_{ab}^{\phantom{ab}c}C_c\psi
\]
for every \(\psi\in\Dcal\), and when the associated gauge evolution is defined on the physical quotient or physical subspace.
\end{definition}

Norm closure of bounded finite constraints, as in \cref{strong:thm:constraint_limit}, is one special case. It does not treat unbounded constraints, changing domains, structure functions, or renormalized operator products.

\begin{definition}[Reflection positive Euclidean limit]\label{strong:definition:012}\label{strong:def:reflection_positive_limit}
A reflection positive Euclidean limit requires a limiting algebra of observables, Euclidean time translations, a reflection operation, and correlation functions whose reflected quadratic forms are nonnegative. A field theoretic reconstruction also requires the regularity and covariance assumptions needed to construct a Hilbert space, a vacuum, and a self adjoint time translation generator.
\end{definition}

\Cref{strong:thm:reflection_limit} shows that positivity of each fixed finite Gram matrix survives pointwise convergence. It does not prove existence, locality, clustering, or the remaining Osterwalder Schrader axioms.

\begin{definition}[Einstein infrared criterion]\label{strong:definition:013}\label{strong:def:einstein_criterion}
A quantum geometric tower meets an Einstein infrared criterion only if a controlled long wavelength effective action exists and has an expansion whose leading geometric terms take the form
\[
\Gamma_{\rm eff}[g]
=
\int
\sqrt{\abs g}
\left(
\Lambda
+
\frac{1}{16\pi G}R
+
\text{controlled higher order terms}
\right)
\dd^dx.
\]
The expansion must respect the limiting gauge identities and must include an estimate for the omitted terms on the declared infrared regime.
\end{definition}

A spectral dimension close to four does not imply this criterion. A heat trace records eigenvalue growth. It does not determine the tensor structure, gauge identities, sign, or interaction coefficients of an effective action.

\begin{definition}[Massless spin two criterion]\label{strong:definition:014}\label{strong:def:spin_two_criterion}
A positive massless spin two sector requires a physical two point function or resolvent with a massless pole in the transverse traceless sector, positive residue on physical states, absence of negative norm modes, and compatibility with the limiting gauge constraints. Interactions must preserve the constraint structure.
\end{definition}

A low energy eigenvalue, a spectral dimension plateau, or a finite rank two tensor packet does not establish a massless spin two particle.

\begin{figure}[t]
\centering
\includegraphics[width=.96\textwidth]{16_continuum_boundary.pdf}
\caption{The finite identities occupy the first step. Every later step needs a separate convergence or reconstruction theorem.}
\label{strong:fig:continuum_boundary}
\end{figure}

\begin{table}[t]
\centering
\small
\begin{tabular}{>{\raggedright\arraybackslash}p{.25\textwidth}>{\raggedright\arraybackslash}p{.31\textwidth}>{\raggedright\arraybackslash}p{.34\textwidth}}
\toprule
Milestone & Finite result available here & Missing continuum input \\
\midrule
Symmetry and constraints & exact packets and quotient obstruction & common unbounded domain and renormalized closure \\
Effective dynamics & exact Feshbach map and uniform gap bounds & cutoff independent spectral separation and locality \\
Euclidean transfer & positive finite transfer and finite reflection positivity & convergent Euclidean fields and reconstruction axioms \\
Spectral dimension & lattice heat trace convergence & relation to a physical metric and interacting observables \\
Lorentzian theory & bounded unitary direct limits & unbounded generator, causality, locality, and field algebra \\
Einstein regime & a definition of the required expansion & derivation and error control for the effective action \\
Spin two sector & a precise spectral criterion & positive physical pole and interaction consistency \\
\bottomrule
\end{tabular}
\caption{What the finite results establish and what remains outside the article.}
\label{strong:tab:continuum_milestones}
\end{table}

\section{Conclusion}\label{strong:section:015}\label{strong:sec:conclusion}

A complete packet family gives one coordinate language for finite quantum structure. Group characters produce symmetry packets. Charged operators obey exact selection rules. Circuits become coherent sums over packet paths. The canonical moving connection separates geometric transport from physical mixing. Density matrices and channels use two sided packets, while channel covariance becomes block preservation on operator space and symmetry of the Choi matrix.

Compression to a code gives a natural error correction obstruction. Its vanishing is the Knill Laflamme condition. Diagonalizing the resulting coefficient matrix produces orthogonal syndrome packets and an exact recovery on the error span. Schur Feshbach elimination gives an equally exact description of hidden Hamiltonian packets. Their frequency response is the self energy, and their time response is a positive definite memory kernel.

Positive Euclidean transfer gives finite reflection positivity. Periodic lattice heat traces converge to the continuum torus at every fixed positive diffusion time, and the continuum small time spectral dimension is recovered after the refinement limit. Constraint closure is measured by a quotient of operator space, not by compression to the physical kernel.

The bounded tower theorems show how exact compatibility and norm summability produce limiting quantum evolutions. They also show the boundary of the argument. Unbounded field Hamiltonians, continuum gauge constraints, reflection positive field measures, the Einstein infrared action, and a massless spin two sector require additional theorems. The finite packet calculus organizes those questions but does not settle them.



\startreadingpart{Mathematical proof completions and retained derivations}
\section{Exact packets and arbitrary circuit length}\label{supp1:section:001}
Let $R$ be a unital ring, not necessarily commutative, and let $I$ be a finite address set. An exact packet family is a family $P_a\in R$ such that
\[
P_aP_b=\begin{cases}P_a,&a=b,\\0,&a\ne b,\end{cases}
\qquad \sum_{a\in I}P_a=1.
\]
Adjoint conditions are added when needed. For an operator $U$, define the transfer corner
\[
U_{ba}=P_bUP_a.
\]
The target index comes first. Completeness and distributivity give
\[
\sum_{b,a}U_{ba}=U,\qquad
(UV)_{ba}=\sum_c U_{bc}V_{ca}.
\]
The second identity uses $P_c^2=P_c$ when adjacent transfer corners are multiplied.

\subsection{Circuit paths: source Theorem 4.1}
For $U=U_m\cdots U_1$, the checked result is
\[
P_{a_m}UP_{a_0}
=\sum_{a_1,\ldots,a_{m-1}}
P_{a_m}U_mP_{a_{m-1}}\cdots P_{a_1}U_1P_{a_0}.
\]
The Lean implementation defines a recursive path sum and proves the identity by induction over an arbitrary finite gate list. It is not restricted to a fixed circuit depth. The empty list has circuit operator $1$ and path value $P_bP_a$. Summing the source and target indices reconstructs the full circuit.



\subsection{Adjoint conservation: source Theorem 4.2}
Assume a star operation with $P_a^*=P_a$. Then
\[
(U_{ba})^*=P_aU^*P_b.
\]
The two checked conservation laws are
\begin{align*}
U^*U=1&\quad\Longrightarrow\quad
\sum_b(U_{ba})^*U_{bc}=\delta_{ac}P_a,\\
UU^*=1&\quad\Longrightarrow\quad
\sum_aU_{ba}(U_{da})^*=\delta_{bd}P_b.
\end{align*}
For the first law, expand each corner, collapse $P_b^2$ and sum:
\[
\sum_bP_aU^*P_bUP_c
=P_aU^*\left(\sum_bP_b\right)UP_c=P_aP_c.
\]
The second law is analogous and is separately checked. Concrete complex-matrix versions use conjugate transpose as the adjoint. No commutative multiplication rule is used for operator products.





\section{Symmetry and measurement}\label{supp1:section:002}
\subsection{Exact symmetry: source Theorem 5.1}
If $HP_a=P_aH$ for every packet, then
\[
P_bHP_a=0\quad(b\ne a),\qquad H=\sum_aP_aHP_a.
\]
In a complex Banach algebra, commutation also propagates to the actual exponential:
\[
e^{zH}P_a=P_ae^{zH}\qquad(z\in\C).
\]
Consequently,
\[
e^{zH}=\sum_aP_ae^{zH}P_a,
\qquad P_be^{zH}P_a=0\quad(b\ne a).
\]
This covers real-time evolution by taking $z=-it$. The Lean proof invokes the exponential commutation theorem, not an abstract flow supplied with the desired invariant-packet property as a hypothesis. Self-adjointness of $H$ is not needed for the algebraic reducing property, so this result applies more broadly than the Hermitian specialization in the paper.



This is an exact symmetry result. It does not supply the quantitative broken-symmetry leakage estimate of Theorem 5.2, whose Duhamel and norm arguments remain outside this block.

\subsection{Measurements: source Theorem 7.1}
For finite complex matrices, define packet dephasing by
\[
\deph(X)=\sum_aP_aXP_a.
\]
The code proves idempotence and trace duality:
\[
\deph(\deph(X))=\deph(X),\qquad
\tr\bigl(\deph(E)\rho\bigr)=\tr\bigl(E\deph(\rho)\bigr).
\]
Trace duality follows term by term from cyclicity:
\[
\tr(P_aEP_a\rho)=\tr(EP_a\rho P_a).
\]
Therefore a packet-diagonal effect, $\deph(E)=E$, satisfies
\[
\boxed{\tr(E\rho)=\tr\bigl(E\deph(\rho)\bigr).}
\]
The proof holds for arbitrary finite complex matrices $E$ and $\rho$. Density matrices and positive measurement effects are valid specializations rather than assumptions needed for this trace identity. Taking $E=1$ also gives trace preservation of packet dephasing.



This does not establish the full Kraus-channel theorem or the Knill--Laflamme recovery criterion. Those require additional positivity, auxiliary-system or recovery-map arguments.



\section{Corrected Feshbach resolvent}\label{supp1:section:003}
Write a finite complex matrix in retained and eliminated coordinates:
\[
H=\begin{pmatrix}A&B\\C&D\end{pmatrix},\qquad
zI-H=\begin{pmatrix}zI_P-A&-B\\-C&zI_Q-D\end{pmatrix}.
\]
Set
\[
D_z=zI_Q-D,\qquad S_z=zI_P-A-BD_z^{-1}C.
\]
The effective Hamiltonian is $H_{\mathrm{eff}}(z)=A+BD_z^{-1}C$, so
$S_z=zI_P-H_{\mathrm{eff}}(z)$.

\subsection{Correct statement: source Theorem 9.1}
Assume both $D_z$ and $zI-H$ are invertible. Then $S_z$ is invertible and
\[
\boxed{\bigl((zI-H)^{-1}\bigr)_{PP}=S_z^{-1}.}
\]
The left-hand side is the literal retained-coordinate submatrix of the full inverse. This matches the paper's compressed operator when written in an orthonormal basis of the retained subspace. The proof is valid for arbitrary complex blocks of compatible sizes and does not require Hermiticity.

The checked invertibility criterion is
\[
D_z\text{ invertible}\quad\Longrightarrow\quad
\bigl(zI-H\text{ invertible}\ \Longleftrightarrow\ S_z\text{ invertible}\bigr).
\]
Thus one may equivalently assume invertibility of $D_z$ and $S_z$.

\textbf{Why the original assumption is insufficient.}
Take $z=0$ and $H=\diag(0,-1)$, retaining the first coordinate. Then
\[
D_z=1,\qquad zI-H=\begin{pmatrix}0&0\\0&1\end{pmatrix}.
\]
The eliminated block is invertible, but the full matrix is singular. Hence the original hypothesis does not ensure the existence of its stated resolvent. The Lean counterexample proves invertibility of the scalar $1$ and noninvertibility of this full $2\times2$ matrix.



\subsection{Provenance and limits}
The general block-inverse proof was reused from the checked fifth block of Paper VI and freshly compiled in the Paper VII namespace. The extra counterexample was added in this release. This is reuse of a proof, not a claim of a new Schur identity.

The unnumbered time-domain memory equation in Section 9.2, cutoff-uniform resolvent control and continuum limits are not proved by this finite inverse identity. The source amendment makes the two invertibility assumptions explicit without changing the intended effective Hamiltonian.



\section{Dispersion with an explicit spacing condition}\label{supp1:section:004}
The grid eigenvalue in source Theorem 11.1 is
\[
\lambda_{k,h}=\frac{4}{h^2}\sin^2\left(\frac{kh}{2}\right).
\]
The original low-mode condition is $0\le kh/2\le1$. A nonzero spacing is also required for the quotient to have its intended meaning. For a physical grid, one normally takes $h>0$.

\subsection{Checked stronger statement}
For every $k,h\in\R$ with $h\ne0$,
\[
\boxed{0\le k^2-\lambda_{k,h}\le\frac{k^4h^2}{12}.}
\]
The code proves this without the low-mode restriction. The bound may be loose away from the low-mode regime, but it remains valid.

The proof begins with the all-real estimate
\[
0\le x^2-\sin^2x\le\frac{x^4}{3}.
\]
The lower bound follows from $|\sin x|\le|x|$. For the upper bound, use
\[
|x-\sin x|\le\frac{|x|^3}{6},\qquad
|x+\sin x|\le2|x|.
\]
Factoring the difference of squares gives
\[
0\le x^2-\sin^2x
=|x-\sin x|\,|x+\sin x|
\le\frac{x^4}{3}.
\]
With $x=kh/2$ and $h\ne0$,
\begin{align*}
k^2-\lambda_{k,h}
&=\frac4{h^2}\left(x^2-\sin^2x\right),\\
\frac4{h^2}\frac{x^4}{3}
&=\frac{k^4h^2}{12}.
\end{align*}
Every algebraic transformation and inequality in the active source was accepted by Lean.



\subsection{What this does not establish}
This is an exact deterministic error bound for the displayed dispersion relation. It is not a verification of the original paper's floating-point convergence-rate experiment, spectral-dimension plateau measurements or selection of a physical spacetime geometry. Those computations were not rerun in this release.



\section{A constraint diagnostic that vanishes automatically}\label{supp1:section:005}
Section 12 defines the physical projector using the exact common kernel of the constraints. In finite dimensions this gives $C_aP=0$ for every $a$. Absorb the factor $i$ and any structure coefficient into $f_c$, and write
\[
\mathcal A_{ab}=C_aC_b-C_bC_a-\sum_cf_cC_c.
\]
\begin{proposition}[Exact common-kernel compression]\label{supp1:proposition:001}
If $C_cP=0$ for every $c$, then
\[
\mathcal A_{ab}P=0,\qquad P\mathcal A_{ab}P=0.
\]
No idempotence or adjoint property of $P$ is needed for this implication.
\end{proposition}
\begin{proof}\phantomsection\label{supp1:proof:001}
Associativity gives
\[
\mathcal A_{ab}P=C_a(C_bP)-C_b(C_aP)-\sum_cf_c(C_cP)=0.
\]
Left multiplication by $P$ gives the second identity.
\phantomsection\label{supp1:proof:001:end}\end{proof}

\subsection{Compressed vanishing without operator closure}
Over the rationals, take
\[
P=\begin{pmatrix}1&0&0\\0&0&0\\0&0&0\end{pmatrix},\quad
X=\begin{pmatrix}0&0&0\\0&0&1\\0&1&0\end{pmatrix},\quad
Z=\begin{pmatrix}0&0&0\\0&1&0\\0&0&-1\end{pmatrix}.
\]
These displayed matrices are symmetric. The code checks that $P^2=P$, $P\ne0$ and $XP=ZP=0$, but
\[
[X,Z]=\begin{pmatrix}0&0&0\\0&0&-2\\0&2&0\end{pmatrix}\ne0,
\qquad P[X,Z]P=0.
\]
The commutator therefore violates the zero-structure-coefficient closure relation outside the physical range, despite exact vanishing of its physical compression.



\subsection{Interpretation for the paper}
Compressed anomaly vanishing is not sufficient to certify full operator closure. In the exact finite common-kernel setup it is automatic. A nonzero computed compression would instead contradict those exact assumptions or expose an approximation or implementation defect.

This does not establish that a nontrivial compatible refinement tower exists. It does not prove or disprove the full continuum program. A stronger closure claim must test more than this compression, and unbounded constraints need their operator domains specified. This audit is additional to the six checked source items and is not counted as a seventh original theorem.




\section{One seed for the computational basis}\label{supp2:section:001}
\textbf{Manuscript item 2.4.}\quad \lean{ComputationalSeed.lean}\hfill 14 declarations

For $n\geq 0$, let
\[
 A_n=(\mathbb Z/2\mathbb Z)^{\{0,\ldots,n-1\}},\qquad
 V_n=\{v:A_n\to\C\}.
\]
For an address $a\in A_n$, define the standard coordinate vector
\[
 \ket a(x)=\begin{cases}1&x=a,\\0&x\ne a.\end{cases}
\]
The seed is $\psi_0=\ket0$. Define a translation operator by
\[
 (T_a v)(x)=v(x-a).
\]
At position $j$, the local Pauli $X$ is $X_j=T_{e_j}$, where $e_j$ has a single nonzero binary coordinate. The formal circuit uses the identity when $a_j=0$ and $X_j$ when $a_j=1$.

\begin{statement}\label{supp2:statement:001}
For every binary address $a$,
\[
 W_a=\prod_{j=0}^{n-1}X_j^{a_j}=T_a,
 \qquad W_a\psi_0=\ket a.
\]
The family $(W_a\psi_0)_{a\in A_n}$ is linearly independent, spans $V_n$, and has exactly $2^n$ elements. It is the standard computational basis.
\end{statement}

\begin{proof}\phantomsection\label{supp2:proof:001}
Composition gives $T_aT_b=T_{a+b}$, and direct evaluation gives $T_a\ket b=\ket{a+b}$. Every binary address satisfies $a+a=0$, so each translation is an involution. In particular, the local operator changes precisely the addressed bit:
\[
 X_j\ket a=\ket{a\text{ with coordinate }j\text{ replaced by }a_j+1}.
\]
The selected factor at $j$ equals translation by the address supported at $j$ with value $a_j$. Induction on a list proves that the product of translations is translation by their sum. The sum of these single-coordinate addresses is $a$. Thus $W_a=T_a$ and $W_a\psi_0=\ket a$.

Lean identifies the resulting family with \lean{Pi.basisFun} over $\C$. Its standard basis theorems supply linear independence and full span. The finite function-space cardinality gives $|A_n|=2^n$.
\phantomsection\label{supp2:proof:001:end}\end{proof}

\textbf{Formal scope.} This is an actual gate-word construction, not an assumption that the addressed states already form a basis. The empty circuit for $n=0$ is included. Orthogonality belongs to the standard computational-basis interpretation; the formal endpoint is identification with that standard coordinate basis, together with independent explicit span and linear-independence theorems.





\section{Tensor and polyphase refinement}\label{supp2:section:002}
\textbf{Manuscript item 3.2.}\quad \lean{TensorRefinement.lean}\hfill 11 declarations

Let $F$ be an arbitrary finite set of factors and let $D_i$ be a finite coordinate index for factor $i$. For matrix families $A_i$, define the coordinate tensor product on the product index space by
\[
 \mathsf T(A)(x,y)=\prod_{i\in F} A_i(x_i,y_i).
\]
The proof first establishes its multiplication, identity, adjoint and finite-sum laws. Thus the construction is not restricted to two factors or to qubits.

Suppose $N>0$ and each factor has a complete orthogonal self-adjoint projector family $(P^{(i)}_j)_{j\in\mathbb Z/N\mathbb Z}$. For a fine address $a:F\to\mathbb Z/N\mathbb Z$, set
\[
 Q_a=\mathsf T\bigl(i\mapsto P^{(i)}_{a_i}\bigr),
 \qquad
 C_r=\sum_{\sum_i a_i=r}Q_a.
\]

\begin{statement}\label{supp2:statement:002}
The fine and coarse families are complete orthogonal self-adjoint projector systems. In particular,
\[
 C_rC_s=\delta_{rs}C_r,\qquad
 \sum_r C_r=I,\qquad C_r^*=C_r.
\]
Each coarse projector is exactly the displayed sum of its fine descendants.
\end{statement}

\begin{proof}\phantomsection\label{supp2:proof:002}
The finite product-of-sums identity gives
\[
 \mathsf T(A)\mathsf T(B)=\mathsf T(i\mapsto A_iB_i).
\]
Consequently, if two fine addresses differ in any coordinate, the corresponding product contains a zero factor. If the addresses agree, the product is the same fine projector. The adjoint passes through the finite coordinate product, so $Q_a^*=Q_a$. Summing over all address functions and using the same product-of-sums identity yields
\[
 \sum_a Q_a=\mathsf T\left(i\mapsto\sum_jP^{(i)}_j\right)=I.
\]
A separate algebraic regrouping theorem applies to any finite orthogonal family and any address map into a finite set. In the double sum for $C_rC_s$, orthogonality removes all terms with different fine addresses. A surviving address can lie in both classes only when $r=s$. Completeness follows by summing each fine projector exactly once. Adjoint preservation follows term by term. Specializing the address map to total residue proves the claim.
\phantomsection\label{supp2:proof:002:end}\end{proof}

\textbf{Formal scope.} Residue grouping is proved for every positive modulus and arbitrary finite factor dimensions. It assumes the supplied local projector laws. The separate cyclic-character construction of those local projectors, item 3.1, is not thereby checked.





\section{Differentiated moving coordinates}\label{supp2:section:003}
\textbf{Manuscript item 6.1.}\quad \lean{MovingFrames.lean}\hfill 9 declarations

Fix a real time $t$ and a finite complex coordinate space. Let $U(s)$ be a unitary matrix curve, differentiable entrywise at $t$, and let $\psi(s)$ be differentiable componentwise at $t$. Write $\dot U$ and $\dot\psi$ for those derivatives. Suppose
\[
 i\dot\psi=H\psi(t).
\]
Only the value of the Hamiltonian at the chosen time is needed; no derivative of $H$ is assumed.

\begin{statement}\label{supp2:statement:003}
The explicitly defined coordinate curve
\[
 y(s)=U(s)^*\psi(s)
\]
is differentiable at $t$, with
\[
 \dot y=\dot U^*\psi(t)+U(t)^*\dot\psi.
\]
It satisfies
\[
 \boxed{\quad i\dot y=H_{\mathrm{eff}}y(t),\qquad
 H_{\mathrm{eff}}=U(t)^*HU(t)-iU(t)^*\dot U.\quad}
\]
The connection term $-iU(t)^*\dot U$ is Hermitian. If $H^*=H$, then $H_{\mathrm{eff}}^*=H_{\mathrm{eff}}$.
\end{statement}

\begin{proof}\phantomsection\label{supp2:proof:003}
Finite entrywise product rules prove the matrix and matrix-vector differentiation formulas. Since the parameter is real, conjugation commutes with differentiation, giving the stated derivative of $U^*$ and hence of $y$.

Differentiate $U^*U=I$ using those rules and uniqueness of derivatives:
\[
 \dot U^*U+U^*\dot U=0.
\]
Thus $(U^*\dot U)^*=-U^*\dot U$, so multiplication by $-i$ produces a Hermitian matrix.

The identity $UU^*=I$ gives $\psi=Uy$ for the constructed coordinate curve. Differentiating and inserting the original equation yields
\[
 i\dot U y+iU\dot y=HUy.
\]
Multiplying by $U^*$ and using $U^*U=I$ gives the transformed equation, with the negative sign on the connection term. Finally, conjugation by $U$ preserves Hermiticity of $H$, and the sum of the two Hermitian terms is Hermitian.
\phantomsection\label{supp2:proof:003:end}\end{proof}

\textbf{Formal scope.} The main certificate constructs $y$ and its derivative rather than assuming a separate derivative witness for the desired answer. Both unitary identities are explicit. The proofs are pointwise and apply at any time where the stated hypotheses hold. They do not assert existence of a global solution to an evolution equation.





\section{The canonical packet connection}\label{supp2:section:004}
\textbf{Manuscript item 6.3.}\quad \lean{CanonicalConnection.lean}\hfill 9 declarations

Let $P_a(s)$ be a finite family in a normed real algebra with a continuous compatible adjoint. Assume actual derivatives $D_a=\dot P_a(t)$ exist and, for all parameter values under consideration,
\[
 P_aP_b=\delta_{ab}P_a,\qquad \sum_aP_a=I,\qquad P_a^*=P_a.
\]
In the supplied Lean theorem, the curves are defined on $\R$ and the projector laws hold at every real parameter. Matrix and bounded-operator settings satisfying the typeclass hypotheses fit this algebraic formulation.

\begin{statement}\label{supp2:statement:004}
Define $C=\sum_aD_aP_a(t)$. Then
\[
 \boxed{\quad C=\frac12\sum_a\comm{D_a}{P_a(t)},\qquad C^*=-C,\quad}
\]
and, for every $b$,
\[
 \boxed{\comm{C}{P_b(t)}=D_b.}
\]
\end{statement}

\begin{proof}\phantomsection\label{supp2:proof:004}
The derivative product rule and uniqueness of derivatives give
\[
 D_aP_b+P_aD_b=\delta_{ab}D_a.
\]
Differentiating completeness and adjoint invariance gives
\[
 \sum_aD_a=0,\qquad D_a^*=D_a.
\]
These are proved derivative consequences, not additional axioms supplied in place of differentiation.

On the diagonal, $D_aP_a+P_aD_a=D_a$. Summing shows
\[
 \sum_aP_aD_a=-\sum_aD_aP_a=-C.
\]
This proves both the half-commutator formula and $C^*=-C$.

Orthogonality directly yields $CP_b=D_bP_b$. For the left product, multiply the differentiated identity for $(b,a)$ on the right by $P_a$ and sum over $a$. Completeness reduces the result to
\[
 D_b+P_bC=D_bP_b.
\]
Subtracting from the right-product identity gives $CP_b-P_bC=D_b$.
\phantomsection\label{supp2:proof:004:end}\end{proof}

\textbf{Formal scope.} The derivative layer is separate from the purely algebraic consequences, making the dependency on genuine derivative hypotheses visible. This proves the one-parameter canonical connection identity. It does not supply the parameter-manifold range bundle, exterior differential calculus, or curvature theorem of item 6.4.





\section{Finite Kraus channels and full amplification}\label{supp2:section:005}
\textbf{Manuscript item 7.2.}\quad \lean{KrausChannels.lean}\hfill 7 declarations

For a finite family of complex matrices $K_a$, define
\[
 \Phi(\rho)=\sum_a K_a\rho K_a^*.
\]
The formalization constructs this as a complex-linear map and proves the following properties directly in the finite matrix representation.

\begin{statement}\label{supp2:statement:005}
The map $\Phi$ is completely positive. If
\[
 \sum_aK_a^*K_a=I,
\]
then $\tr\Phi(\rho)=\tr\rho$ for every matrix $\rho$. In particular, positive semidefinite matrices of trace one are mapped to positive semidefinite matrices of trace one.
\end{statement}

\begin{proof}\phantomsection\label{supp2:proof:005}
Matrix multiplication distributes across addition and scalar multiplication, which gives complex linearity. If $\rho\succeq0$, each term $K_a\rho K_a^*$ is positive semidefinite, hence so is their finite sum.

To prove complete positivity rather than positivity alone, let $J$ be an arbitrary finite auxiliary index set. For a matrix $X$ indexed by $J\times n$, define its blockwise amplification by
\[
 [\Phi^{(J)}(X)]_{(i,u),(j,v)}
 =\Phi\bigl(X_{(i,\cdot),(j,\cdot)}\bigr)_{u,v}.
\]
The explicit index calculation proves the identity
\[
 \boxed{\quad
 \Phi^{(J)}(X)=\sum_a (I_J\otimes K_a)X(I_J\otimes K_a)^*.
 \quad}
\]
Thus this is the actual identity-tensor amplification, not an unrelated map assumed to be positive. Every summand on the right is positive when $X$ is positive, and this holds for every finite $J$.

Finally, cyclicity and linearity of the trace give
\begin{align*}
 \tr\Phi(\rho)
 &=\sum_a\tr(K_a\rho K_a^*)\\
 &=\tr\left(\left(\sum_aK_a^*K_a\right)\rho\right)
 =\tr\rho.
\end{align*}
Positivity together with trace preservation yields the density-matrix conclusion.
\phantomsection\label{supp2:proof:005:end}\end{proof}

\textbf{Formal scope.} Complete positivity is quantified over all finite auxiliary dimensions, not sampled dimensions. Normalization is required only for trace preservation and the density-matrix conclusion. The converse representation theorem for arbitrary completely positive maps is not asserted here.






\section{Actual orbits and sharp seed counts}\label{supp3:section:001}
Let $C\in M_{d\times m}(\C)$ be the coefficient matrix of one seed. Define its orbit by actual matrix multiplication:
\[
\mathcal O(C)=\{AC:A\in M_d(\C)\}.
\]
This set is already a linear subspace. The source proves its equality with the linear span of the same orbit, rather than assuming an abstract orbit-dimension formula.

A matrix $X$ belongs to $\mathcal O(C)$ exactly when each of its rows belongs to $\operatorname{Row}(C)$. One direction follows by expanding $AC$. For the converse, choose the coefficients producing each row separately and place them in the corresponding row of $A$. An explicit linear equivalence therefore gives
\[
\mathcal O(C)\simeq \operatorname{Row}(C)^d,
\qquad
\boxed{\dim_\C\mathcal O(C)=d\,\rank C.}
\]
If $d>0$, equality with the full $d\times m$ coefficient space is equivalent to $\rank C=m$. Consequently a cyclic seed requires $m\le d$. The nonempty-row assumption is essential for canceling the factor $d$.

\subsection{Several seeds: both necessity and construction}
For seeds $C_1,\ldots,C_k$, generation means the concrete statement
\[
\forall X\in M_{d\times m}(\C),\quad
\exists A_1,\ldots,A_k\in M_d(\C):\quad
X=\sum_{a=1}^k A_aC_a.
\]
Stack the $C_a$ vertically into a $kd\times m$ matrix $D$. The source proves the exact stacked-action identity and the equivalence
\[
\text{the seeds generate}\quad\Longleftrightarrow\quad\rank D=m.
\]
Thus $m\le kd$ is necessary. It is sufficient as well: embed the $m$ coordinate columns into $kd$ coordinate positions to obtain a full-column-rank matrix $D$, then split its rows back into $k$ seeds. This is an existence proof of actual generating matrices, not only an inequality on possible dimensions.

\subsection{Simultaneous direct-sum generation}
For the stipulated block representation
\[
\mathcal A=\bigoplus_i\bigl(M_{d_i}(\C)\otimes I_{m_i}\bigr),\qquad d_i>0,
\]
the proof establishes that $k$ seeds generate all blocks simultaneously exactly when $m_i\le kd_i$ for every $i$. It follows that
\[
\boxed{\srank_{\mathcal A}(\mathcal H)=\max_i\left\lceil\frac{m_i}{d_i}\right\rceil.}
\]
The final Lean proposition is an \code{IsLeast} statement over the set of feasible seed counts. It proves both attainment and minimality. For an empty block family the finite supremum is defined as $0$, matching the zero module.





\section{Schmidt rank and the reduced state}\label{supp3:section:002}
Use standard product-basis coordinates $\psi_{ij}$ and let $C_{ij}=\psi_{ij}$. An explicit linear equivalence identifies the product-index state vector with its coefficient matrix. The implementation checks the actual Kronecker action:
\[
\operatorname{coeff}\bigl((A\otimes I_B)\psi\bigr)=AC.
\]
It also proves that the transported matrix orbit equals the span of these Kronecker-transformed state vectors. The state-space dimension formula is therefore connected to the physical local action, not only to a separate matrix model.

\subsection{An independent definition of Schmidt rank}
The source defines Schmidt rank as the least number of product-vector summands. A decomposition of length $s$ is the concrete equality
\[
C_{ij}=\sum_{a=1}^{s}u_a(i)v_a(j).
\]
It factors $C=UV$ through an $s$-dimensional intermediate space, so $\rank C\le s$. Conversely, choose a basis of the row space of $C$. Expanding each row in this basis constructs a decomposition with exactly $\rank C$ product terms. Thus
\[
\boxed{r(\psi)=\rank C.}
\]
The definition uses a least natural number, and the equality with matrix rank is proved. The proof does not define Schmidt rank to be matrix rank. It does not construct an orthonormal Schmidt decomposition or prove a new singular-value decomposition theorem.

The matrix-orbit dimension and the coefficient equivalence now yield
\[
\boxed{\dim_\C\spanof\{(A\otimes I_B)\psi:A\in M_{d_A}(\C)\}
=d_A r(\psi).}
\]

\subsection{The actual partial trace and its support}
The reduced matrix is defined by tracing the addressed index of the pure-state density matrix:
\[
(\rho_B)_{j\ell}=\sum_i\psi_{ij}\overline{\psi_{i\ell}}.
\]
For the coefficient convention used here, the exact Gram identity is
\[
\rho_B=C^{\mathsf T}(C^{\mathsf T})^*,
\qquad \rank\rho_B=\rank C=r(\psi).
\]
The transpose is important. With this convention, writing $C^*C$ instead would conjugate the reduced matrix. The Lean source checks the entrywise identity and positivity of the resulting matrix.

Full support is expressed as surjectivity of the reduced matrix's linear map. Assuming $d_A>0$, the checked equivalence is
\[
\boxed{\mathcal O_A(\psi)=\mathcal H_A\otimes\mathcal H_B
\quad\Longleftrightarrow\quad \Ran\rho_B=\mathcal H_B.}
\]
Normalization of $\psi$ is unnecessary for these dimension and support statements. The zero state has rank $0$ under the minimum-product-length definition.





\section{Constructed cyclic character projectors}\label{supp3:section:003}
Let $N>0$ and let $S$ be a unitary complex matrix with $S^N=I$. The implementation uses the genuine standard additive character of $\mathbb Z/N\mathbb Z$ and proves its exponential-root form:
\[
\omega=\exp(2\pi i/N),\qquad
\chi(r)=\omega^{\operatorname{val}(r)}.
\]
Here $\operatorname{val}(r)$ is the canonical representative in $\{0,\ldots,N-1\}$. The representation is constructed as
\[
\rho(r)=S^{\operatorname{val}(r)}.
\]
Reduction of exponents modulo $N$ is proved from $S^N=I$, giving $\rho(r+s)=\rho(r)\rho(s)$ and $\rho(0)=I$. No representation or character-orthogonality assumption is inserted in the final matrix theorem.

\subsection{Exact coefficient formula}
The packet definition is a finite sum over residue classes:
\[
P_j=\frac1N\sum_{r\in\mathbb Z/N\mathbb Z}\chi(-rj)\rho(r).
\]
The source proves
\[
\chi(-rj)=\bigl(\omega^{\operatorname{val}(r)\operatorname{val}(j)}\bigr)^{-1},
\]
so this is precisely the manuscript's $N$-term Fourier projector formula, with canonical representatives.

\subsection{Algebraic identities}
Primitivity of the standard character yields the checked finite sum
\[
\sum_r\chi(tr)=\begin{cases}N,&t=0,\\0,&t\ne0.\end{cases}
\]
Reindexing finite sums then gives
\[
\boxed{P_jP_k=\delta_{jk}P_j,\qquad \sum_jP_j=I,\qquad SP_j=\omega^jP_j.}
\]
Multiplicativity and completeness are first proved for a general representation in a complex algebra, then instantiated with the actual powers of $S$. This organization isolates reusable finite-sum algebra without weakening the final hypotheses.

\subsection{Orthogonal, not merely idempotent}
Unitarity supplies $\rho(r)^*=\rho(-r)$, while the character obeys $\overline{\chi(r)}=\chi(-r)$. Reindexing $r\mapsto-r$ proves $P_j^*=P_j$. The final concrete matrix theorem includes Hermiticity together with multiplication, completeness, and eigenaction. Thus the result establishes orthogonal spectral packets, not only algebraic idempotents.

The proof includes the degenerate modulus $N=1$: finite order then forces $S=I$ and the sole projector is $I$. It does not need an artificial assumption $N>1$.





\section{The actual Euclidean transfer exponential}\label{supp3:section:004}
Let $H$ be a Hermitian matrix on a nonempty finite complex state space. The source defines
\[
T_e(t)=\exp\bigl(-t(H-eI)\bigr)
\]
using Lean's power-series exponential. It does not introduce an unrelated transfer family or assume that a spectral expansion equals this exponential.

\subsection{Connecting power series and spectral calculus}
The shifted generator is proved Hermitian. The implementation then identifies the actual matrix exponential with continuous functional calculus:
\[
\boxed{T_e(t)=f_t(H),\qquad f_t(x)=e^{-t(x-e)}.}
\]
The formal proof uses composition, scalar multiplication, subtraction, and the established agreement between the functional-calculus exponential and the power-series exponential. Spectral mapping is consequently available for the operator that was actually defined:
\[
\Spec_\R T_e(t)=\{e^{-t(x-e)}:x\in\Spec_\R H\}.
\]
Hermiticity also identifies the full complex spectrum with the image of these real eigenvalues. A separate checked theorem states the complex-spectrum interval containment explicitly.

\subsection{A genuine minimum eigenvalue}
The ground energy is the minimum of the finite, nonempty eigenvalue set. The proof establishes both
\[
E_0\in\Spec_\R H,\qquad E_0\le x\quad(x\in\Spec_\R H).
\]
The nonempty-state-space assumption is recorded in the theorem because the minimum of an empty spectrum is not defined by the manuscript's formula.

\subsection{Positivity, semigroup composition, and the interval}
Every eigenvalue of $T_e(t)$ is strictly positive, for arbitrary real $e,t$, since the real exponential is positive. For $e=E_0$ and $t\ge0$ it is also at most one. Therefore
\[
\boxed{T_{E_0}(t)>0,\qquad \Spec T_{E_0}(t)\subseteq(0,1].}
\]
The two scalar multiples of $H-eI$ commute, so the actual power-series exponential obeys
\[
\boxed{T_e(t+s)=T_e(t)T_e(s),\qquad T_e(0)=I.}
\]
The semigroup identity is proved for all real $t,s$. Nonnegative time is needed for the upper spectral bound, not for the composition identity or strict positivity.

These are finite-dimensional operator results. No continuum transfer limit, gravitational path integral, or physical spacetime model is established by this theorem. Numerical figures in the original paper were not recomputed in this release.






\section{Quantitative leakage from the actual exponential}\label{supp4:section:001}
Let $A$ be a unital complex C-star algebra, let $H=H^*$, and let $P=P^*=P^2$. Put $Q=1-P$. The checked statement is
\[
\boxed{\|Q\exp(-itH)P\|\le |t|\,\|HP-PH\|\qquad(t\in\R).}
\]
The matrix specialization uses the Hilbert-space operator norm, implemented by \path{Matrix.Norms.L2Operator}. It is not an entrywise norm or the row-sum operator norm. The theorem includes the zero projector and empty matrix index type. There is no nonzero-time hypothesis.

\subsection{Derivative proof rather than an assumed Duhamel estimate}
For a skew-adjoint generator $G$, define the actual conjugation path
\[
T_P(t)=\exp(tG)P\exp(-tG).
\]
The formalization derives its derivative from the power-series exponential and the product rule:
\[
T_P'(t)=\exp(tG)(GP-PG)\exp(-tG).
\]
Both exponential factors are proved unitary. Left and right multiplication by a unitary preserve the C-star norm. Therefore
\[
\|T_P'(t)\|=\|GP-PG\|,
\qquad
\|T_P(t)-T_P(s)\|\le |t-s|\,\|GP-PG\|.
\]
The second inequality follows from the normed-space mean value inequality and covers either ordering of $s$ and $t$. It is not a finite-difference estimate or a numerical time discretization.

Multiplication by $\exp(tG)$ gives
\[
(T_P(t)-P)\exp(tG)=\exp(tG)P-P\exp(tG).
\]
For every orthogonal projector, $\|P\|\le1$ and $\|1-P\|\le1$. The exact corner identity
\[
(1-P)UP=(1-P)(UP-PU)P
\]
then bounds leakage by the exponential commutator. Taking $G=-iH$ gives the stated theorem because $\|[G,P]\|=\|[H,P]\|$.

\subsection{Checked consequences}
Unitarity also gives the capped estimate
\[
\|Q\exp(-itH)P\|\le\min\{1,\ |t|\,\|[H,P]\|\}.
\]
If $H=H_0+V$, both summands are self-adjoint, and $[H_0,P]=0$, then
\[
\|Q\exp(-it(H_0+V))P\|\le2|t|\,\|V\|.
\]
The latter uses $\|[V,P]\|\le2\|V\|$; its constant is not claimed to be optimal. Exact commutation gives zero leakage. The work does not include unbounded generators or time-dependent Hamiltonians.



\section{Actual projected differentiation and curvature}\label{supp4:section:002}
Let $E,F$ be real normed spaces. Let $P:E\to\mathcal L(F,F)$ be differentiable with $P(x)^2=P(x)$ for every $x$. Let $s:E\to F$ be $C^2$ with $P(x)s(x)=s(x)$, and let $X,Y:E\to E$ be differentiable vector fields. The checked connection is the actual expression
\[
\nabla_Xs(x)=P(x)\,Ds(x)[X(x)].
\]
It is proved to preserve the range condition. It is linear in the vector-field direction and obeys the scalar Leibniz rule
\[
\nabla_X(cs)=Dc[X]\,s+c\nabla_Xs
\]
for differentiable real scalar functions $c$. No differentiated projector identities are inserted as hypotheses in the main curvature theorem.

\subsection{Differentiated projector and section identities}
Differentiating $P^2=P$ and $Ps=s$ produces
\begin{align*}
DP[v]P+PDP[v]&=DP[v],\\
PDP[v]P&=0,\\
Ds[v]&=DP[v]s+P Ds[v].
\end{align*}
The zero sandwich follows algebraically from the first identity. All products preserve their noncommutative order.

\subsection{The vector-field bracket is retained}
Define the actual curvature action by
\[
R_P(X,Y)s=\nabla_X\nabla_Ys-\nabla_Y\nabla_Xs-\nabla_{[X,Y]}s.
\]
The source uses the actual Fr\'echet-derivative Lie bracket. Symmetry of second derivatives is supplied by the $C^2$ hypothesis, not postulated separately. Expanding the connection gives
\[
R_P(X,Y)s=P\bigl(DP[X]Ds[Y]-DP[Y]Ds[X]\bigr).
\]
Substituting the section identity and removing the zero sandwich terms proves
\[
\boxed{R_P(X,Y)s=P\bigl(DP[X]DP[Y]-DP[Y]DP[X]\bigr)P s.}
\]
Thus the coefficient is precisely $P(\dd P\wedge\dd P)P$ with the convention
$(\dd P\wedge\dd P)(v,w)=DP[v]DP[w]-DP[w]DP[v]$.

\subsection{Tensor properties and the remaining global boundary}
The coefficient is proved bilinear, alternating, and supported on the range on both sides. The chain rule proves the pullback identity
\[
F_{P\circ f}(x)(v,w)=F_P(f(x))(Df(x)v,Df(x)w).
\]
The main theorem has global regularity hypotheses on normed parameter spaces. The pullback identity alone does not supply localization, smooth local trivializations, or gluing. The global range-bundle construction and connection gluing are supplied in the final completion below.



\section{A constructive and total syndrome recovery}\label{supp4:section:003}
Let $P$ be an orthogonal code projector. The input certificate consists of a finite family $V_a$ satisfying
\[
V_aP=V_a,
\qquad
V_a^*V_b=\delta_{ab}P.
\]
These are explicit partial-isometry and orthogonal-syndrome hypotheses. The block does not claim to have obtained them from arbitrary Knill--Laflamme Gram data.

\subsection{Constructing the recovery on the whole ambient space}
Define the syndrome projectors and their complement by
\[
S_a=V_aV_a^*,\qquad S=\sum_aS_a,\qquad Q=1-S.
\]
The source proves self-adjointness, $S_aS_b=\delta_{ab}S_a$, $S^2=S$, and $Q^2=Q=Q^*$. It also proves $QV_a=0$ and $V_a^*Q=0$.

The recovery is not left undefined outside the syndrome subspaces. Its explicit Kraus operators are $R_a=V_a^*$ together with $R_{\mathrm{extra}}=Q$, giving
\[
\mathcal R(\rho)=\sum_aV_a^*\rho V_a+Q\rho Q.
\]
They satisfy exact completeness:
\[
\sum_aR_a^*R_a+R_{\mathrm{extra}}^*R_{\mathrm{extra}}
=\sum_aV_aV_a^*+Q=S+Q=I.
\]
For finite complex matrices, trace preservation and complete positivity at every finite auxiliary dimension are proved for this actual map. The amplified map is identified with the Kraus family $I_J\otimes R_a$, including the complement term. Positivity is not merely checked on a single matrix size.

\subsection{Exact correction, including coherent errors}
For every code-supported operator $P\rho P=\rho$, the central identity is
\[
\boxed{\mathcal R(V_a\rho V_b^*)=\delta_{ab}\rho.}
\]
Consequently normalized mixtures satisfy
\[
\mathcal R\!\left(\sum_ap_aV_a\rho V_a^*\right)=\rho
\quad\text{when}\quad\sum_ap_a=1.
\]
The identity is algebraic and is proved for complex weights. Physical mixtures use the nonnegative real case. For a coherent error $E_z=\sum_a z_aV_a$, the stronger cross-term calculation yields
\[
\mathcal R(E_z\rho E_z^*)=\left(\sum_a z_a\overline{z_a}\right)\rho.
\]
Thus normalized coherent superpositions are corrected as well. This is an explicit recovery stage, not a replacement definition that makes the full Knill--Laflamme criterion true by construction.




\section{The precise finite channel statement}\label{supp5:section:001}
Let $\mathcal H=\C^n$ be a finite coordinate space and let $P=P^*=P^2$. A matrix is code supported when $P\rho P=\rho$. For a finite family of square matrices $E_a$, put
\[
\mathcal N_E(\rho)=\sum_a E_a\rho E_a^*.
\]
The input channel normalization is the actual matrix identity
\[
\sum_a E_a^*E_a=I.
\]
A finite Kraus recovery has the form $\mathcal R(\sigma)=\sum_r R_r\sigma R_r^*$, with $\sum_rR_r^*R_r=I$. This representation, together with the checked amplification results, provides complete positivity and trace preservation on the entire ambient matrix space.

\subsection{Necessary and sufficient conditions}
Under the stated normalization, the following equivalence is proved:
\[
\boxed{\begin{gathered}
\exists\,\mathcal R\text{ of the above finite Kraus form,}\\[-2pt]
\mathcal R\mathcal N_E(\rho)=\rho
\quad\text{for every }\rho\succeq0,\ \Tr\rho=1,\ P\rho P=\rho
\end{gathered}}
\quad\Longleftrightarrow\quad
\boxed{\begin{gathered}
\exists C\succeq0,\\
PE_a^*E_bP=C_{ab}P\quad\text{for all }a,b.
\end{gathered}}
\]
The code-space statement is not limited to a selected basis of states. The formalization separately proves equivalence with correction of \emph{every} complex matrix supported on the code. It also proves preservation with arbitrary finite reference systems, as detailed in the physical-state bridge below.



\subsection{Degenerate cases are not silently excluded}
There is no positive-definiteness assumption on $C$. Zero eigenvalues, redundant error operators, empty positive-eigenvalue support and the zero projector are included. If $P=0$, the only supported matrix is zero and no trace-one code density matrix exists. The proofs handle that case separately rather than dividing by the code rank or by $\Tr P$.

For an unnormalized error family, the correct statement is generally
\[
\mathcal R\mathcal N_E(\rho)=\Tr(C)\rho,
\]
not unconditional equality to $\rho$. Normalization forces $\Tr C=1$ on a nonzero code. The unrestricted-family formulation and the same-decoder span result appear in the coherent-error discussion below.



\section{From Gram data to an actual total recovery}\label{supp5:section:002}
Assume $C\succeq0$ and $PE_a^*E_bP=C_{ab}P$. First restrict the error operators to the code, $F_a=E_aP$. The finite Hermitian spectral theorem supplies a unitary $U$ such that
\[
U^*CU=\operatorname{diag}(\lambda_r),\qquad \lambda_r\ge0.
\]


Define the transformed errors by an actual finite sum,
\[
B_r=\sum_a U_{ar}F_a.
\]
The new module proves, rather than assumes,
\[
B_rP=B_r,\qquad B_r^*B_s=\delta_{rs}\lambda_rP,
\qquad \sum_rB_r\rho B_r^*=\sum_aF_a\rho F_a^*.
\]
For code-supported $\rho$, the restricted and unrestricted error sums are equal. Thus the unitary change of error coordinates preserves the actual channel on the code.

\subsection{Null directions are removed by proof}
If $\lambda_r=0$, then $B_r^*B_r=0$. Positivity of the ordinary complex matrix inner product gives $B_r=0$. This is the checked declaration \path{DiagonalErrors.zero_weight_error}. The normalized syndrome operators are defined only on
\[
A_+=\{r:\lambda_r>0\},\qquad V_r=\lambda_r^{-1/2}B_r.
\]
Consequently no inverse square root at a null eigenvalue is needed. Lean proves
\[
V_rP=V_r,\qquad V_r^*V_s=\delta_{rs}P.
\]
The previous block's recovery is now instantiated from these derived syndromes, rather than from an assumed orthogonal error family.

\subsection{A channel on the whole ambient space}
Put $Q=I-\sum_{r\in A_+}V_rV_r^*$. The recovery is
\[
\boxed{\mathcal R(\sigma)=\sum_{r\in A_+}V_r^*\sigma V_r+Q\sigma Q.}
\]
The complement is an orthogonal projector and
\[
\sum_{r\in A_+}V_rV_r^*+Q^*Q=I.
\]
Hence the recovery is trace preserving. Its finite Kraus form also proves positivity after amplification by the identity on every finite auxiliary space. The complement term matters: the map is defined as a channel even outside the syndrome ranges.

Finally the checked calculation gives $\mathcal R\mathcal N_E(\rho)=\sum_r\lambda_r\rho=\Tr(C)\rho$. Under input-channel normalization and $P\ne0$, summing the diagonal Gram identities gives $\Tr(C)P=P$, hence $\Tr C=1$. The zero-code branch does not use this cancellation.



\section{Deriving the Gram matrix from recovery}\label{supp5:section:003}
Suppose an actual finite trace-preserving recovery corrects every code-supported operator. The proof does not assume the Knill--Laflamme inner-product condition or a pre-existing syndrome decomposition. Instead, it expands the composite channel into the finite Kraus family
\[
A_{ra}=R_rE_aP.
\]
For every ambient matrix $\rho$, correction on the code implies
\[
\sum_{r,a} A_{ra}\rho A_{ra}^*=P\rho P.
\]
This equality on all matrices is obtained from the code-supported sandwich $P\rho P$, using the projector identities.

\subsection{Rank-one rigidity of the corrected channel}
The module tests the map equality on actual matrix units and flattens each $A_{ra}$ into a row. If $X$ is the matrix of those rows and $p$ is the vector of entries of $P$, the resulting Gram identity is
\[
X^*X=\overline p\,p^{\mathsf T}.
\]
The checked rank-one Gram lemma proves that every row of $X$ is proportional to $p$. For nonzero $p$, a nonzero coordinate is selected and a two-coordinate test vector is placed in the null space of the Gram matrix. The implication $X^*Xv=0\Rightarrow Xv=0$ then gives the required proportionality. When $p=0$, the Gram identity forces $X=0$ directly.

Applied to the corrected Kraus family, this yields scalars $z_{ra}$ with
\[
R_rE_aP=z_{ra}P.
\]
The essential declarations are \path{rank_one_gram_rigidity}, \path{kraus_rank_one_rigidity} and \path{corrected_kraus_family}. No rank-one Choi characterization is inserted as an axiom or as the desired conclusion in hypothesis form.

\subsection{Trace-preserving recovery supplies the original Gram data}
Let $Z_{ra}=z_{ra}$. The recovery normalization gives
\begin{align*}
PE_a^*E_bP
&=PE_a^*\Bigl(\sum_rR_r^*R_r\Bigr)E_bP\\
&=\sum_r(R_rE_aP)^*(R_rE_bP)\\
&=\Bigl(\sum_r\overline{z_{ra}}z_{rb}\Bigr)P
=(Z^*Z)_{ab}P.
\end{align*}
Thus $C=Z^*Z$ is positive semidefinite and satisfies the required compressed identities. The theorem \path{gram_of_recovery} implements this calculation with the original error and recovery matrices.

A second rigidity argument handles a nonnegative scalar multiple of the identity code map. It uses $\sqrt t\,P$ as the rank-one map generator and includes $t=0$. That argument gives necessity for unnormalized error families without pretending that they are channels.



\section{Density matrices, all operators and finite references}\label{supp5:section:004}
The operator-space formulation used in necessity must agree with the physical density-state formulation. This bridge is proved explicitly for complex-linear maps on finite matrices.

\subsection{Normalization is not an extra assumption}
A nonzero positive semidefinite matrix has strictly positive real trace. If a complex-linear map fixes every normalized code density matrix, apply that hypothesis to
\[
\frac{\rho}{\Tr\rho}
\]
for each nonzero positive code-supported $\rho$. Linearity then gives equality on $\rho$ itself. The zero matrix is handled separately. The source proves the real-trace, nonzero-trace and scalar-cancellation facts rather than assuming that normalization is valid.

\subsection{Complex polarization reaches the whole operator space}
Write $q(x)=xx^*$. The module proves the exact identity
\[
4xy^*=q(x+y)-q(x-y)
+i\bigl(q(x+iy)-q(x-iy)\bigr).
\]
Agreement on all positive pure outer products therefore implies agreement on every outer product. Actual matrix units then span the full matrix space. Applying this argument to the sandwich by $P$ proves equality on every code-supported operator, including non-Hermitian ones.

The result \path{density_corrects_iff_operators} establishes the equivalence for the actual composite Kraus maps. It does not replace the set of density matrices by an assumed basis of logical states.

\subsection{An arbitrary finite reference system}
For any finite reference index set $A$, the amplified channel is defined using the actual Kronecker Kraus operators,
\[
\widetilde{\mathcal N}_E(\rho)
=\sum_a(I_A\otimes E_a)\rho(I_A\otimes E_a)^*.
\]
Its reference blocks are proved equal to the original channel applied blockwise. If
\[
(I_A\otimes P)\rho(I_A\otimes P)=\rho,
\]
then every reference block is code supported. Applying the checked operator-space correction to those blocks gives
\[
\boxed{\widetilde{\mathcal R}\,\widetilde{\mathcal N}_E(\rho)=\rho.}
\]
This includes entangled density matrices supported on the reference--code space, and in fact holds for every jointly supported matrix. There is no fixed auxiliary dimension or test-state sampling.

The declaration \path{recovery_with_reference} proves the equality. The recovery's complete positivity is separately established by the actual finite Kraus amplification, so preservation of states is not inferred from a merely positive, non-completely-positive map.



\section{Coherent combinations and normalization}\label{supp5:section:005}
A correction proof for only one selected Kraus presentation is weaker than a proof that one decoder works throughout the error span. The new module connects the constructed spectral recovery back to the original error operators.

\subsection{Cross terms for the original errors}
For arbitrary positive Gram data, the same recovery satisfies
\[
\boxed{\mathcal R(E_a\rho E_b^*)=C_{ba}\rho\qquad(P\rho P=\rho).}
\]
The index order $ba$ is intentional: the defining Gram relation is $PE_a^*E_bP=C_{ab}P$. The proof first derives scalar corrected Kraus composites for the constructed recovery and then uses its trace-preserving completeness to identify the original Gram entries.

For $E(c)=\sum_a c_aE_a$, finite-sum expansion gives
\[
\mathcal R\bigl(E(c)\rho E(c)^*\bigr)
=\left(\sum_{a,b}c_a\overline{c_b}C_{ba}\right)\rho
=(c^*Cc)\rho.
\]
The double-sum formula is the literal checked conclusion; the final quadratic-form notation follows by renaming finite indices. It is not restricted to orthogonal errors or real coefficients.

\subsection{Every normalized channel in the same span}
Suppose $F_j=\sum_a c_{ja}E_a$ and $\sum_jF_j^*F_j=I$. The decoder is \emph{not} reconstructed for the $F_j$. Summing the coherent identities gives a scalar multiple $q\rho$. Trace preservation of both the $F$ channel and the recovery forces $q=1$ on a nonzero code. The proof uses positivity of $P$ to justify cancellation of its trace and handles $P=0$ separately. Therefore
\[
\boxed{\mathcal R\Bigl(\sum_jF_j\rho F_j^*\Bigr)=\rho.}
\]
This is \path{Data.recovery_span_channel}, a quantified statement over arbitrary finite $F$ families and their complex coefficients.

\subsection{The unnormalized error-family equivalence}
For arbitrary finite $E$, define exact family recovery to mean that a finite trace-preserving Kraus recovery and a single $t\ge0$ exist such that
\[
\mathcal R\mathcal N_E(\rho)=t\rho
\quad\text{for every code-supported }\rho.
\]
Then \path{error_family_correctable_iff_gram} proves that this condition is equivalent to the existence of $C\succeq0$ satisfying the Gram identities. Sufficiency chooses $t=\Tr C$; necessity constructs $C$ from the scaled rank-one rigidity argument.

The scale is essential. Zero errors can give $t=0$ and cannot be inverted to a nonzero normalized state. An input with a larger total trace is not silently declared a channel. For normalized input channels, the separate channel theorem removes the scale exactly as described in the sufficiency argument above.



\section{Exact examples, trust audit and remaining work}\label{supp5:section:006}
\subsection{Singular, nonorthogonal errors are accepted}
On a qubit with full code $P=I$, take
\[
E_0=\frac35I,\qquad E_1=\frac45I,\qquad
C=\begin{pmatrix}9/25&12/25\\12/25&16/25\end{pmatrix}.
\]
Lean checks the channel normalization, the Gram identities, positivity of $C$, $\det C=0$, and the nonzero off-diagonal entry. It then applies the completed density-state criterion to prove correctability. This is an exact rational-complex example, not a floating-point tolerance check.

\subsection{A genuine channel can still be uncorrectable}
For full qubit dephasing, $E_0=\operatorname{diag}(1,0)$ and $E_1=\operatorname{diag}(0,1)$, the Kraus normalization is checked. But $E_0^*E_0$ cannot be a scalar multiple of $I$, as its two diagonal entries differ. Necessity therefore proves that no finite recovery channel corrects the full code space. The zero-code case is also checked independently.


\section{The global projected connection}\label{supp6:section:001}
Let $M$ be a smooth real manifold without boundary, $F$ a complete normed real vector space and $P:M\to\End(F)$ a smooth family of bounded idempotents. For smooth range sections $s$, define
\[
\nabla_Xs=P\,\dd s[X].
\]
The range bundle is constructed from $P$ rather than assumed. Its local connection formulas glue, and the actual covariant commutator satisfies
\[
\boxed{(\nabla_X\nabla_Y-\nabla_Y\nabla_X-\nabla_{[X,Y]})s
=P\bigl(\dd P[X]\dd P[Y]-\dd P[Y]\dd P[X]\bigr)P\,s.}
\]
The derivative and Lie bracket are Mathlib's actual manifold operators. No global coordinate chart is assumed. The result requires idempotence, not an additional orthogonality hypothesis.


\section{Local products constructed from the projector}\label{supp6:section:002}
The fiber over $x$ is the actual fixed subspace
\[
\mathcal E_x=\{v\in F:P(x)v=v\}=\Ran P(x).
\]
The total space is the subset $\mathcal E=\{(x,v):P(x)v=v\}\subseteq M\times F$, with the induced topology. The construction includes actual local homeomorphisms of this space, not merely abstract transition functions.

\subsection{An invertible intertwiner near each anchor}
For idempotents $p,q$, set
\[
T(p,q)=pq+(I-p)(I-q).
\]
Lean proves $T(p,p)=I$ and $pT(p,q)=T(p,q)q$. For an anchor $a\in M$, define
\[
U_a=\{x:T(P(x),P(a))\text{ is invertible}\}.
\]
The invertible elements form an open subset of the bounded-operator algebra. Consequently $U_a$ is open and contains $a$. The family $(U_a)$ covers the whole base. Invertibility is \emph{derived locally}, not assumed at every pair of points.

\subsection{The mutually inverse local maps}
Write $T_a(x)=T(P(x),P(a))$. For $x\in U_a$,
\[
B_a(x):\mathcal E_a\longrightarrow\mathcal E_x,
\quad v\longmapsto T_a(x)v,
\]
\[
A_a(x):\mathcal E_x\longrightarrow\mathcal E_a,
\quad w\longmapsto P(a)T_a(x)^{-1}w.
\]
The intertwining identity and its inverse counterpart show that the displayed targets are correct. Lean proves both inverse laws, continuity, linearity in the fiber and smooth dependence on the base. Thus
\[
\mathcal E|_{U_a}\ \cong\ U_a\times\mathcal E_a
\]
is an actual base-preserving local product homeomorphism.

\subsection{Smooth overlap maps and rank}
The coordinate transition is $g_{ab}(x)=A_b(x)B_a(x)$. It is smooth on $U_a\cap U_b$, agrees with the actual change between local homeomorphisms, and satisfies
\[
g_{aa}=I,\qquad g_{bc}g_{ab}=g_{ac}.
\]
A continuous linear equivalence between $\mathcal E_a$ and $\mathcal E_x$ also gives local constancy of fiber dimension. No single constant rank across disconnected components is assumed.





\section{A connection on the constructed bundle}\label{supp6:section:003}
The global formula is defined using the actual manifold derivative:
\[
\nabla_Xs(x)=P(x)\,D_Ms(x)[X(x)].
\]
Here $D_M$ is represented by \path{mvfderiv}. A second definition, \path{GlobalConnection.covariant}, has values in the actual fiber:
\[
\nabla s(x):T_xM\longrightarrow\mathcal E_x.
\]
It is not only an ambient expression asserted to land in the range.

\subsection{Connection axioms and smoothness}
The new proofs establish range preservation, additivity in the section, linearity over functions in the vector-field argument, and the scalar Leibniz identity
\[
\nabla_X(fs)=\dd f[X]\,s+f\nabla_Xs.
\]
Smooth $P$, smooth sections and smooth vector fields give smooth covariant derivatives. Equality of section germs gives equality of covariant derivatives at the point, so the operation is local.

\subsection{Smooth sections are genuinely intrinsic}
Let $x\in U_a$. If $P(y)s(y)=s(y)$ in a neighborhood of $x$, then
\[
s\text{ is smooth at }x
\quad\Longleftrightarrow\quad
\bigl(y\mapsto A_a(y)s(y)\bigr)\text{ is smooth at }x.
\]
The range condition is only local. The proof uses the actual inverse local maps and equality of germs. It follows that smoothness in bundle coordinates does not depend on the chosen chart.

\subsection{The local connection and its change of coordinates}
For a coordinate function $u$ valued in $\mathcal E_a$, define
\[
\nabla^{(a)}_Xu=A_a\bigl(P\,D_M(B_a u)[X]\bigr).
\]
For every overlap point, Lean proves the operator identity
\[
\boxed{\nabla^{(b)}_X(g_{ab}u)=g_{ab}\nabla^{(a)}_Xu.}
\]
This is proved for arbitrary local coordinate functions. It does not assume that $u$ was obtained from a globally smooth section or that a smooth extension has been supplied. Together with the local smoothness criterion, this closes the gluing step left open in the fifth block.





\section{From differentiated projectors to the two-form}\label{supp6:section:004}
The curvature is defined by the covariant commutator, including the \emph{actual manifold Lie bracket}. It is not defined to equal the desired projector expression.

\subsection{The flat ambient identity is proved in charts}
For a smooth vector-valued function $s$, the proof identifies directional derivatives with the pullback into each extended chart. It then proves
\[
D_Ms([X,Y])=D_M(D_Ms[Y])[X]-D_M(D_Ms[X])[Y].
\]
This uses the chart pullback of the actual vector fields and the ordinary differential identity, then transfers the result back to the manifold. The point-local auxiliary module permits use within arbitrary open chart domains.

\subsection{Projector and section identities are derived}
Differentiating $P^2=P$ in a tangent direction $v$ yields
\[
(\dd P[v])P+P\dd P[v]=\dd P[v],
\qquad P\dd P[v]P=0.
\]
Differentiating $Ps=s$ gives
\[
\dd s[v]=\dd P[v]s+P\dd s[v].
\]
Both statements are consequences of actual derivative rules. They are not additional assumptions.

\subsection{Cancellation of the second derivatives}
The product rule and the ambient Lie-bracket identity reduce curvature to
\[
R^\nabla(X,Y)s=P\dd P[X]\dd s[Y]-P\dd P[Y]\dd s[X].
\]
Substitute the differentiated range-section identity. The terms containing $P\dd P[v]P$ vanish, leaving
\[
R^\nabla(X,Y)s=
P\bigl(\dd P[X]\dd P[Y]-\dd P[Y]\dd P[X]\bigr)P s.
\]
Lean also proves that this depends only on $X(x)$, $Y(x)$ and $s(x)$, not on their extensions. The operator-valued form is alternating, additive and homogeneous, and lands in the actual fiber. Constant projector bundles have zero curvature form.





\subsection{Geometry and scalar conventions}
The main global curvature theorem uses a smooth real manifold without boundary, modeled on a complete normed real space, and bounded real-linear endomorphisms of a complete normed real space. Finite-dimensional parameter and quantum spaces fall within the corresponding real differential calculus. Complex quantum spaces can be viewed through their underlying real normed spaces; this block does not introduce a separate complex-bundle typeclass wrapper.

The atlas and global connection constructions also support the more general model-with-corners infrastructure where their hypotheses apply. The \emph{curvature completion claim} is specifically for the boundaryless manifold theorem. A boundary/corner curvature generalization is not claimed.

\subsection{Explicit smooth atlas, not an assumed one}
The implementation exposes the actual total space, local product homeomorphisms, linear transitions, smoothness proofs and cocycle laws. It does not register these as a \path{ContMDiffVectorBundle} instance. This is an interface choice: no range-bundle atlas, local trivialization, overlap identity or curvature conclusion is inserted as a premise.

The smooth-section equivalence and gauge law make the atlas usable intrinsically. The connection is also defined with values in \path{fiber (P x)}. Six proof-bearing constructed objects, including both local equivalences and the fiber-valued derivative, have separate axiom audits.



\section{Earlier compact derivations retained as alternatives}\label{legacy:section:001}\label{legacy:alternatives}
The following derivations retain valid mathematical arguments from the original manuscript. Their hypotheses and notation agree with the corrected finite core. The earlier abbreviated error-correction argument is replaced by the complete normalized and scaled recovery proofs above, and its substantive conclusions are included there.
\subsection{Duhamel proof of fixed-packet leakage}\label{legacy:duhamel}
Let $H=H^*$ on a finite-dimensional complex Hilbert space, $P=P^*=P^2$, $Q=I-P$, and $t\in\mathbb R$.
\begin{proof}[Alternative proof of the commutator leakage bound]\phantomsection\label{legacy:proof:001}
Because $QP=0$,
\[
Qe^{-itH}P=Q[e^{-itH},P]P.
\]
The Duhamel identity gives
\[
[e^{-itH},P]=-i\int_0^t e^{-i(t-s)H}[H,P]e^{-isH}\,ds.
\]
Both unitary factors have norm one, so taking norms and using the integration interval's length $|t|$ yields
$\|Qe^{-itH}P\|\le |t|\,\|[H,P]\|$.
\phantomsection\label{legacy:proof:001:end}\end{proof}
\subsection{Exterior-form proof of the projected curvature}\label{legacy:exterior}
Let $P$ be a smooth idempotent field and extend the projected connection to range-bundle-valued differential forms by $\nabla\omega=P\,d\omega$. Thus $\nabla^2$ denotes the exterior covariant square, including the exterior derivative's Lie-bracket term when evaluated on vector fields.
\begin{proof}[Alternative proof of the curvature formula]\phantomsection\label{legacy:proof:002}
For a section $s$ with $Ps=s$,
\[
\nabla^2s=P\,d(P\,ds)=P(dP)\wedge ds.
\]
Since $ds=d(Ps)=(dP)s+Pds$ and $P(dP)P=0$, the surviving term is
\[
P(dP\wedge dP)Ps.
\]
\phantomsection\label{legacy:proof:002:end}\end{proof}
\subsection{Restricted low-mode proof}\label{legacy:lowmode}
Let $h\ne0$, $x=kh/2$, and $0\le x\le1$.
\begin{proof}[Alternative proof in the original low-mode regime]\phantomsection\label{legacy:proof:003}
Since $0\le\sin x\le x$, the error is nonnegative. On $[0,1]$,
\[
\sin x\ge x-\frac{x^3}{6}\ge0.
\]
Therefore
\[
\sin^2x\ge\left(x-\frac{x^3}{6}\right)^2\ge x^2-\frac{x^4}{3}.
\]
Hence $x^2-\sin^2x\le x^4/3$. Multiplication by $4/h^2$ gives
\[
0\le k^2-\frac4{h^2}\sin^2(kh/2)\le\frac{k^4h^2}{12}.
\]
\phantomsection\label{legacy:proof:003:end}\end{proof}

\clearpage\FloatBarrier
\section*{References}
\addcontentsline{toc}{section}{References}
\renewcommand{\refname}{}
\bibliographystyle{amsplain}
\bibliography{references}
\end{document}
